\documentclass[12pt,a4paper,notitlepage]{article}
\usepackage{amsmath}
\usepackage{amssymb}
\usepackage{amsbsy}
\usepackage{mathtools}
\usepackage{bbm}
\usepackage{bm}
\usepackage{float}
\usepackage[french]{babel}
\usepackage{tikz}

\usepackage{palatino}

 \usepackage[active]{srcltx}
\usepackage{scrtime}

% Notations communes avec les transparents et avec la fiche d'exercices.
% -*- mode: latex -*-
%
% Notations communes à la fiche d'exercices et à ses éléments de correction.
%
% Les définitions reprennent à l'identique celles de cours/preamble.tex
% (l. 110-123), de façon que les documents de travaux dirigés emploient
% exactement les mêmes symboles que les transparents, sans avoir à charger tout
% le préambule beamer (qui suppose \coursfile défini et appelle git via
% \write18).
%
% Convention structurante du cours, à respecter dans tous les énoncés : le
% « modèle de la nature » (le DGP) s'écrit avec des lettres grecques (\beta,
% \varepsilon), le « modèle de l'économètre » avec des latines grasses
% (\mathbf b) et des erreurs \epsilon. L'estimateur des MCO est \hat{\mathbf b}.
%
% Il n'existe pas de macro \E ni \Var dans ce dépôt : l'espérance et la variance
% s'écrivent \mathbb E[...] et \mathbb V[...]. Ne pas en introduire.

\newcommand{\trace}{\mathrm{tr}}
\newcommand{\tracarg}[1]{\mathrm{tr}\left\{#1\right\}}
\newcommand{\iid}[2]{\mathrm{iid}\left(#1,#2\right)}
\newcommand{\normal}[2]{\mathcal N\left(#1,#2\right)}
\newcommand{\epsvar}{\sigma_{\varepsilon}^2}
\newcommand{\sample}{\mathcal Y_T}
\newcommand{\samplet}[1]{\mathcal Y_{#1}}

% Convergences. Le cours indexe \plim et \dlim sur T ; les exercices écrits en
% coupe individuelle ont besoin de la variante en N, d'où les deux jeux.
\newcommand{\plimT}{\xrightarrow[T\rightarrow\infty]{\text{proba}}}
\newcommand{\dlimT}{\xrightarrow[T\rightarrow\infty]{\text{loi}}}
\newcommand{\plimN}{\xrightarrow[N\rightarrow\infty]{\text{proba}}}
\newcommand{\dlimN}{\xrightarrow[N\rightarrow\infty]{\text{loi}}}


\newcounter{qnumber}
\setcounter{qnumber}{0}

\newcounter{enumber}
\setcounter{enumber}{0}

\newcommand{\question}{\textbf{(\addtocounter{qnumber}{1}\theqnumber)}\,}
\newcommand{\exercice}{\refstepcounter{enumber}\textbf{\textsc{Exercice} \theenumber}.\,}

% Ouvre un exercice en remettant le compteur de questions à zéro. Les deux
% compteurs sont indépendants de ceux de fiche-exercices.tex : toute insertion
% d'un exercice dans la fiche doit être répercutée ici, faute de quoi les deux
% documents se désynchronisent silencieusement.
\newcommand{\nouvelexercice}{\setcounter{qnumber}{0}\bigskip\exercice}

\setlength{\parindent}{0cm}

% Absorbe les petits dépassements de justification (lignes trop longues) en
% répartissant l'élasticité sur les espaces plutôt que dans la marge.
\setlength{\emergencystretch}{3em}

\begin{document}

\title{\textsc{Économétrie approfondie}\\ \small{(Éléments de correction des exercices)}}
\date{Le \today\ à \thistime}

\maketitle

\thispagestyle{empty}

Ce document donne des éléments de correction pour les 35 exercices de la fiche.
La numérotation des exercices et des questions suit celle de la fiche. On
respecte la convention du cours~: le modèle de la nature est écrit avec des
lettres grecques, celui de l'économètre avec des lettres latines, et
l'estimateur des MCO est noté $\hat{\mathbf b}$.

\bigskip

\section*{Chapitre I~-- Les MCO quand tout va bien}

\nouvelexercice\label{ex:rupture-fixe} Le modèle empirique ne contient qu'une constante~: la matrice
des explicatives se réduit au vecteur $\iota=(1,\dots,1)'$.

\question L'estimateur des MCO minimise
$\mathcal S(c) = \sum_{t=1}^T (y_t-c)^2$. La condition du premier ordre s'écrit
$-2\sum_{t=1}^T (y_t-\hat c_T) = 0$, d'où~:
\[
  \hat c_T = \frac{1}{T}\sum_{t=1}^T y_t = \bar y
\]
On retrouve bien $\hat c_T = (\iota'\iota)^{-1}\iota'\mathbf y$. La condition du
second ordre est satisfaite ($\mathcal S''=2T>0$)~: il s'agit d'un minimum.

\question En substituant le DGP et en remarquant que
$\sum_{t=1}^T\mathbbm{1}_{\{t>\kappa\}}(t) = T-\kappa$ (il y a $T-\kappa$ dates
postérieures à $\kappa$)~:
\[
  \begin{split}
    \mathbb E\left[\hat c_T\right]
    &= \frac{1}{T}\sum_{t=1}^T \mathbb E\left[\mu_1 + \mu_2\mathbbm{1}_{\{t>\kappa\}}(t) + \varepsilon_t\right]\\
    &= \mu_1 + \frac{\mu_2}{T}\left(T-\kappa\right)
  \end{split}
\]
L'estimateur n'estime donc \emph{ni} $\mu_1$, \emph{ni} $\mu_1+\mu_2$~: il
estime la moyenne de la série sur la période, c'est-à-dire une moyenne pondérée
des deux régimes, les poids étant les fréquences $\kappa/T$ et $(T-\kappa)/T$.
Le biais par rapport à $\mu_1$ vaut $\mu_2(T-\kappa)/T$.

\question La partie déterministe du DGP ne contribue pas à la variance~:
\[
  \mathbb V\left[\hat c_T\right]
  = \mathbb V\left[\frac{1}{T}\sum_{t=1}^T \varepsilon_t\right]
  = \frac{1}{T^2}\sum_{t=1}^T \mathbb V\left[\varepsilon_t\right]
  = \frac{\epsvar}{T}
\]
en utilisant l'indépendance des $\varepsilon_t$. Cette variance tend vers zéro
lorsque $T\rightarrow\infty$.

\question Comme $\kappa$ est \emph{fixe}, $(T-\kappa)/T\rightarrow 1$. Par
ailleurs $\bar\varepsilon\plimT 0$ (loi des grands nombres). Donc~:
\[
  \hat c_T \plimT \mu_1+\mu_2
\]
\textbf{Interprétation.} Asymptotiquement, la sous-période antérieure à la
rupture devient négligeable~: elle ne compte que $\kappa$ observations sur $T$.
Tout se passe comme si l'échantillon était entièrement engendré par le second
régime. L'estimateur converge, mais vers le mauvais paramètre si l'on cherchait
à estimer $\mu_1$.

\nouvelexercice\label{ex:rupture-proportionnelle} Seule la date de rupture change~: elle est désormais
proportionnelle à la taille de l'échantillon.

\question Inchangé~: $\hat c_T = \bar y$.

\question Il y a maintenant $T-[\kappa T]$ dates postérieures à la rupture,
donc~:
\[
  \mathbb E\left[\hat c_T\right] = \mu_1 + \mu_2\,\frac{T-[\kappa T]}{T}
\]

\question Inchangée également, $\mathbb V[\hat c_T] = \epsvar/T$~: la
modification ne porte que sur la partie déterministe du DGP, qui est de variance
nulle.

\question Comme $[\kappa T]/T \rightarrow \kappa$~:
\[
  \hat c_T \plimT \mu_1+(1-\kappa)\mu_2
\]

\question \textbf{Comparaison.} Dans l'exercice~\ref{ex:rupture-fixe}, la rupture intervient à une
date \emph{fixe}~: la proportion d'observations du premier régime, $\kappa/T$,
tend vers zéro et le premier régime disparaît asymptotiquement. Dans l'exercice
2, la rupture intervient à une date \emph{proportionnelle}~: chaque régime
conserve une part fixe de l'échantillon ($\kappa$ et $1-\kappa$), et la limite
est la moyenne pondérée des deux régimes. La leçon est que le comportement
asymptotique d'un estimateur dépend de la façon dont on fait croître
l'échantillon, et pas seulement du DGP à $T$ donné.

\nouvelexercice\label{ex:mco-matriciel}

\question On empile les $N$ observations. En posant
$\mathbf y = (y_1,\dots,y_N)'$, $\bm\epsilon = (\epsilon_1,\dots,\epsilon_N)'$,
$\mathbf b = (b_0,b_1,\dots,b_K)'$ et~:
\[
  X =
  \begin{pmatrix}
    1 & x_{1,1} & \ldots & x_{K,1}\\
    1 & x_{1,2} & \ldots & x_{K,2}\\
    \vdots & \vdots & & \vdots\\
    1 & x_{1,N} & \ldots & x_{K,N}
  \end{pmatrix}
\]
$X$ est de dimension $N\times(K+1)$ (chaque ligne rassemble les explicatives
d'une observation, chaque colonne les $N$ observations d'une variable). Le
modèle s'écrit $\mathbf y = X\mathbf b + \bm\epsilon$.

\question L'estimateur des MCO minimise la somme des carrés des résidus~:
\[
  \begin{split}
    \hat{\mathbf b} &= \arg\min_{\mathbf b}\ (\mathbf y - X\mathbf b)'(\mathbf y - X\mathbf b)\\
                    &= \arg\min_{\mathbf b}\ \mathbf y'\mathbf y - 2\mathbf b'X'\mathbf y + \mathbf b'X'X\mathbf b
  \end{split}
\]
(les scalaires $\mathbf b'X'\mathbf y$ et $\mathbf y'X\mathbf b$ sont égaux, et
$\mathbf y'\mathbf y$ ne dépend pas de $\mathbf b$). La condition du premier
ordre s'écrit~:
\[
  -2X'\mathbf y + 2X'X\hat{\mathbf b} = 0
  \quad\Longleftrightarrow\quad
  X'X\hat{\mathbf b} = X'\mathbf y
\]
Ce sont les \emph{équations normales}. Si $X$ est de plein rang colonne, $X'X$
est inversible et~:
\[
  \hat{\mathbf b} = \left(X'X\right)^{-1}X'\mathbf y
\]
La condition du second ordre est satisfaite car le hessien $2X'X$ est défini
positif sous l'hypothèse de plein rang.

\question En explicitant les produits~:
\[
  X'X =
  \begin{pmatrix}
    N & \sum_i x_{1,i} & \ldots & \sum_i x_{K,i}\\
    \sum_i x_{1,i} & \sum_i x_{1,i}^2 & \ldots & \sum_i x_{1,i}x_{K,i}\\
    \vdots & \vdots & & \vdots\\
    \sum_i x_{K,i} & \sum_i x_{K,i}x_{1,i} & \ldots & \sum_i x_{K,i}^2
  \end{pmatrix},
  \quad
  X'\mathbf y =
  \begin{pmatrix}
    \sum_i y_i\\
    \sum_i x_{1,i}y_i\\
    \vdots\\
    \sum_i x_{K,i}y_i
  \end{pmatrix}
\]
La matrice $X'X$ est symétrique~: c'est (à un facteur $N$ près) la matrice des
moments d'ordre deux non centrés des explicatives.

\question Pour $K=1$~:
\[
  X'X = \begin{pmatrix} N & \sum_i x_i\\ \sum_i x_i & \sum_i x_i^2\end{pmatrix},
  \qquad
  \det(X'X) = N\sum_i x_i^2 - \left(\sum_i x_i\right)^2
\]
\textit{Démonstration de la stricte positivité du déterminant.} On développe~:
\[
  \begin{split}
    \sum_{i=1}^N\sum_{j=1}^N (x_i-x_j)^2
    &= \sum_i\sum_j \left( x_i^2 - 2x_ix_j + x_j^2 \right)\\
    &= N\sum_i x_i^2 - 2\left(\sum_i x_i\right)\left(\sum_j x_j\right) + N\sum_j x_j^2\\
    &= 2N\sum_i x_i^2 - 2\left(\sum_i x_i\right)^2
  \end{split}
\]
d'où~:
\[
  \det(X'X) = \frac{1}{2}\sum_{i=1}^N\sum_{j=1}^N (x_i-x_j)^2
\]
Cette quantité est une somme de carrés, donc positive~; elle est
\emph{strictement} positive dès que $N>1$ et que les $x_i$ ne sont pas tous
égaux. C'est exactement la condition de plein rang~: si tous les $x_i$ étaient
égaux, la seconde colonne de $X$ serait proportionnelle à la première.

\smallskip

En inversant la matrice $2\times 2$~:
\[
  (X'X)^{-1} = \frac{1}{\det(X'X)}\begin{pmatrix}\sum_i x_i^2 & -\sum_i x_i\\ -\sum_i x_i & N\end{pmatrix}
\]
et donc~:
\[
  \hat b_1 = \frac{N\sum_i x_iy_i - \left(\sum_i x_i\right)\left(\sum_i y_i\right)}{N\sum_i x_i^2 - \left(\sum_i x_i\right)^2}
           = \frac{\sum_i (x_i-\bar x)(y_i-\bar y)}{\sum_i (x_i-\bar x)^2},
  \qquad
  \hat b_0 = \bar y - \hat b_1\bar x
\]
en utilisant $\sum_i (x_i-\bar x)^2 = \sum_i x_i^2 - N\bar x^2$ et
$\sum_i (x_i-\bar x)(y_i-\bar y) = \sum_i x_iy_i - N\bar x\bar y$. La pente est
le rapport de la covariance empirique à la variance empirique de $x$.

\question La première ligne des équations normales $X'\hat\epsilon = 0$
correspond à la colonne de uns~:
\[
  \iota'\hat\epsilon = \sum_{i=1}^N \hat\epsilon_i = 0
\]
La somme des résidus est donc nulle \emph{par construction}, dès lors que le
modèle contient une constante. C'est une propriété algébrique, vraie pour tout
échantillon, et non une hypothèse statistique.

\nouvelexercice\label{ex:tendance}

\question Comme à l'exercice~\ref{ex:rupture-fixe}, $\hat c_T = \bar y$.

\question En utilisant $\sum_{t=1}^T t = \frac{T(T+1)}{2}$~:
\[
  \mathbb E\left[\hat c_T\right]
  = \frac{1}{T}\sum_{t=1}^T \left(\mu_0 + \mu_1 t\right)
  = \mu_0 + \mu_1\,\frac{T+1}{2}
\]
L'estimateur est donc \textbf{biaisé}, le biais valant $\mu_1(T+1)/2$.

\question Le biais ne disparaît pas~: il \emph{diverge} avec $T$. C'est la
différence essentielle avec les exercices 1 et 2, où le biais restait borné par
$\mu_2$. \textbf{Interprétation.} La variable omise est ici une tendance
déterministe, non stationnaire~: sa moyenne sur l'échantillon croît sans borne.
Aucune quantité d'observations ne corrige l'erreur de spécification, au
contraire. Le cours signale d'ailleurs (chapitre~I, discussion de
$\mathcal H_2$) que le modèle à tendance déterministe est justement le
contre-exemple pour lequel $X'X/T$ ne converge pas vers une matrice finie~: les
énoncés asymptotiques usuels y demandent des normalisations différentes.

\nouvelexercice\label{ex:inversion-blocs}

\question On cherche $\mathcal A^{-1}$ sous forme de blocs
$\mathcal B = \begin{pmatrix} B_{11} & B_{12}\\ B_{21} & B_{22}\end{pmatrix}$
et on identifie les blocs de $\mathcal A\mathcal B = I$~:
\[
  \begin{cases}
    A_{11}B_{11} + A_{12}B_{21} = I & (a)\\
    A_{11}B_{12} + A_{12}B_{22} = 0 & (b)\\
    A_{21}B_{11} + A_{22}B_{21} = 0 & (c)\\
    A_{21}B_{12} + A_{22}B_{22} = I & (d)
  \end{cases}
\]
De $(b)$ on tire $B_{12} = -A_{11}^{-1}A_{12}B_{22}$. En reportant dans
$(d)$~:
\[
  \left(A_{22}-A_{21}A_{11}^{-1}A_{12}\right)B_{22} = I
  \quad\Longleftrightarrow\quad
  B_{22} = B \equiv \left(A_{22}-A_{21}A_{11}^{-1}A_{12}\right)^{-1}
\]
d'où $B_{12} = -A_{11}^{-1}A_{12}B$. De $(c)$ on tire
$B_{21} = -A_{22}^{-1}A_{21}B_{11}$, mais il est plus direct de partir de $(a)$~:
$B_{11} = A_{11}^{-1}\left(I - A_{12}B_{21}\right)$, puis de reporter dans
$(c)$~:
\[
  A_{21}A_{11}^{-1}\left(I-A_{12}B_{21}\right) + A_{22}B_{21} = 0
  \quad\Longleftrightarrow\quad
  \left(A_{22}-A_{21}A_{11}^{-1}A_{12}\right)B_{21} = -A_{21}A_{11}^{-1}
\]
soit $B_{21} = -BA_{21}A_{11}^{-1}$, et en reportant dans l'expression de
$B_{11}$~:
\[
  B_{11} = A_{11}^{-1}\left(I + A_{12}BA_{21}A_{11}^{-1}\right)
\]
ce qui est bien le résultat annoncé. La matrice
$A_{22}-A_{21}A_{11}^{-1}A_{12}$ est appelée \emph{complément de Schur} de
$A_{11}$~; son inversibilité est l'hypothèse faite dans l'énoncé.

\question Si $A_{12}=0$ et $A_{21}=0$, alors $B = A_{22}^{-1}$, $B_{12}=0$,
$B_{21}=0$ et $B_{11}=A_{11}^{-1}$~:
\[
  \mathcal A^{-1} = \begin{pmatrix} A_{11}^{-1} & 0\\ 0 & A_{22}^{-1}\end{pmatrix}
\]
L'inverse d'une matrice bloc-diagonale est la matrice bloc-diagonale des
inverses. C'est ce résultat qui, appliqué à $X'X$, donnera à l'exercice suivant
la séparabilité des estimateurs sous orthogonalité.

\nouvelexercice\label{ex:fwl} On note $\hat{\mathbf b} = (\hat{\mathbf b}_1', \hat{\mathbf b}_2')'$
la partition correspondante de l'estimateur, et $P_1 = X_1(X_1'X_1)^{-1}X_1'$.

\question Si $X_1'X_2 = 0$, la matrice $X'X$ est bloc-diagonale~:
\[
  X'X = \begin{pmatrix} X_1'X_1 & 0\\ 0 & X_2'X_2\end{pmatrix}
\]
D'après la question (2) de l'exercice précédent, son inverse est
bloc-diagonale, et comme $X'\mathbf y = (X_1'\mathbf y,\ X_2'\mathbf y)'$~:
\[
  \hat{\mathbf b}_1 = (X_1'X_1)^{-1}X_1'\mathbf y,
  \qquad
  \hat{\mathbf b}_2 = (X_2'X_2)^{-1}X_2'\mathbf y
\]
\textbf{Interprétation.} Sous orthogonalité, la régression multiple se scinde en
deux régressions séparées~: estimer $\beta_1$ ne demande pas de connaître $X_2$.
Autrement dit, \emph{omettre} $X_2$ ne biaise pas $\hat{\mathbf b}_1$ lorsque
$X_2$ est orthogonale à $X_1$ --- c'est le cas particulier sans biais du
problème de la variable omise (exercice~\ref{ex:variable-omise}).

\question Sans orthogonalité, la première ligne des équations normales
$X'X\hat{\mathbf b} = X'\mathbf y$ donne~:
\[
  X_1'X_1\hat{\mathbf b}_1 + X_1'X_2\hat{\mathbf b}_2 = X_1'\mathbf y
\]
soit~:
\[
  \hat{\mathbf b}_1 = (X_1'X_1)^{-1}X_1'\mathbf y - (X_1'X_1)^{-1}X_1'X_2\hat{\mathbf b}_2
                    = (X_1'X_1)^{-1}X_1'\left(\mathbf y - X_2\hat{\mathbf b}_2\right)
\]
On retrouve l'expression de la question précédente, \emph{corrigée} du terme
$(X_1'X_1)^{-1}X_1'X_2\hat{\mathbf b}_2$. Ce terme s'annule si et seulement si
$X_1'X_2 = 0$~: l'estimateur de $\beta_1$ dans la régression courte et dans la
régression longue coïncident exactement dans ce cas.

\question Par définition des résidus estimés~:
\[
  \hat\epsilon = \mathbf y - X\hat{\mathbf b}
               = \mathbf y - X(X'X)^{-1}X'\mathbf y
               = \left(I - X(X'X)^{-1}X'\right)\mathbf y
               = M\mathbf y
\]
avec $M = I - X(X'X)^{-1}X'$.

\question $M$ est symétrique (car $(X'X)^{-1}$ l'est), et~:
\[
  \begin{split}
    MM &= \left(I-X(X'X)^{-1}X'\right)\left(I-X(X'X)^{-1}X'\right)\\
       &= I - 2X(X'X)^{-1}X' + X\underbrace{(X'X)^{-1}X'X}_{=I}(X'X)^{-1}X'\\
       &= I - X(X'X)^{-1}X' = M
  \end{split}
\]
$M$ est donc idempotente~: c'est un projecteur. Elle projette sur l'orthogonal
de l'espace engendré par les colonnes de $X$.

\question
\[
  MX = \left(I-X(X'X)^{-1}X'\right)X = X - X\underbrace{(X'X)^{-1}X'X}_{=I} = X - X = 0
\]
Les résidus sont donc orthogonaux aux explicatives~:
$X'\hat\epsilon = X'M\mathbf y = (MX)'\mathbf y = 0$.

\question
\[
  \hat{\mathbf y} = X\hat{\mathbf b} = X(X'X)^{-1}X'\mathbf y = H\mathbf y
\]
avec $H = X(X'X)^{-1}X' = I-M$. Un calcul identique à celui de la question (4)
donne $HH = H$~: $H$ est également idempotente. Elle projette sur l'espace
engendré par les colonnes de $X$. On a de plus $HM = MH = 0$ et
$\mathbf y = \hat{\mathbf y} + \hat\epsilon$~: les MCO décomposent $\mathbf y$
en deux composantes orthogonales.

\question La seconde ligne des équations normales s'écrit~:
\[
  X_2'X_1\hat{\mathbf b}_1 + X_2'X_2\hat{\mathbf b}_2 = X_2'\mathbf y
\]
En y reportant l'expression de $\hat{\mathbf b}_1$ obtenue à la question (2)~:
\[
  X_2'X_1(X_1'X_1)^{-1}X_1'\left(\mathbf y - X_2\hat{\mathbf b}_2\right) + X_2'X_2\hat{\mathbf b}_2 = X_2'\mathbf y
\]
soit, en regroupant et en reconnaissant $P_1$~:
\[
  X_2'\left(I-P_1\right)X_2\,\hat{\mathbf b}_2 = X_2'\left(I-P_1\right)\mathbf y
\]
c'est-à-dire, avec $M_1 = I-P_1$~:
\[
  \hat{\mathbf b}_2 = \left(X_2'M_1X_2\right)^{-1}X_2'M_1\mathbf y
\]
Comme $M_1$ est symétrique et idempotente, on peut écrire
$X_2'M_1X_2 = (M_1X_2)'(M_1X_2)$ et $X_2'M_1\mathbf y = (M_1X_2)'(M_1\mathbf y)$,
d'où, en posant $\tilde X_2 = M_1X_2$ et $\tilde{\mathbf y} = M_1\mathbf y$~:
\[
  \hat{\mathbf b}_2 = \left(\tilde X_2'\tilde X_2\right)^{-1}\tilde X_2'\tilde{\mathbf y}
\]
\textbf{Interprétation (théorème de Frisch-Waugh-Lovell).} $\hat{\mathbf b}_2$
s'obtient en trois étapes~: (i) régresser $X_2$ sur $X_1$ et garder les résidus
$\tilde X_2$~; (ii) régresser $\mathbf y$ sur $X_1$ et garder les résidus
$\tilde{\mathbf y}$~; (iii) régresser $\tilde{\mathbf y}$ sur $\tilde X_2$. Le
coefficient $\hat{\mathbf b}_2$ mesure donc l'effet de $X_2$ \emph{une fois
retiré de $X_2$ et de $\mathbf y$ tout ce que $X_1$ explique}. C'est le sens
précis de l'expression «~à autres variables contrôlées~». L'intérêt pratique est
double~: on n'inverse que des matrices de taille $K_2$, et les propriétés
statistiques de $\hat{\mathbf b}_2$ sont inchangées.

\nouvelexercice\label{ex:mc-convergence} \textit{Exercice de simulation. On donne ici les résultats
attendus et leur lecture~; les scripts correspondants figurent dans
\texttt{exercices/E00.py} et \texttt{exercices/E01.py}.}

\question Pour $N=10$, la moyenne empirique des $50\,000$ valeurs de $\hat b$
est très proche de $\beta = \frac12$, et leur variance empirique est proche de
$\epsvar/\sum_i (x_i-\bar x)^2$. Ces deux quantités sont des estimations
Monte-Carlo de $\mathbb E[\hat b]$ et $\mathbb V[\hat b]$~: on simule la
distribution d'échantillonnage de l'estimateur en répétant l'expérience que
l'économètre, lui, ne réalise qu'une fois. La première illustre l'absence de
biais, la seconde la formule de variance du cours.

\question Lorsque $N$ passe de $10$ à $10\,000$, la moyenne reste centrée sur
$\frac12$ --- l'absence de biais ne dépend pas de la taille de l'échantillon ---
tandis que la variance est divisée par un facteur voisin de $10$ à chaque fois
que $N$ est multiplié par $10$. C'est la convergence~: $\mathbb V[\hat b]$ est
d'ordre $1/N$, donc $\hat b\plimN\beta$. La distribution se concentre autour de
la vraie valeur.

\question Avec une variable explicative stochastique, l'estimateur reste sans
biais~: puisque $\mathbb E[\varepsilon_i\mid x] = 0$, la loi des espérances
itérées donne
$\mathbb E[\hat b] = \mathbb E\left[\mathbb E[\hat b\mid x]\right] = \beta$
(c'est l'objet de l'exercice~\ref{ex:lie}). Les variances sont très légèrement
supérieures à celles du cas déterministe pour les petits $N$, car
$\mathbb V[\hat b] = \epsvar\,\mathbb E\left[\left(\sum_i(x_i-\bar x)^2\right)^{-1}\right]$
et la fonction $u\mapsto 1/u$ est convexe (inégalité de Jensen). L'écart
s'estompe quand $N$ croît.

\question La variance de $\hat b$ est approximativement
$\epsvar/(N\sigma_x^2)$. En passant de $\sigma_x^2 = 2$ à $0{,}2$, la variance
est multipliée par $10$~; en passant à $20$, elle est divisée par $10$.
\textbf{Interprétation.} Plus la variable explicative est dispersée, plus
l'estimation de la pente est précise~: on dispose d'un «~bras de levier~» plus
long pour identifier la droite. C'est le rapport signal sur bruit du chapitre~I.
Une explicative quasi constante ne permet pas d'identifier sa pente --- c'est
aussi, en germe, le problème de la multicolinéarité (exercice~\ref{ex:multicolinearite}).

\nouvelexercice\label{ex:mc-distribution} \textit{Exercice de simulation.}

\smallskip

La distribution de $\hat b - \frac12$ se concentre autour de zéro à mesure que
$N$ croît~: l'intervalle dans lequel tombent les valeurs simulées se resserre, et
à la limite toute la masse se retrouve concentrée au point $0$. Représentée telle
quelle, cette distribution ne renseigne donc que sur la convergence.

\smallskip

En revanche, la distribution de
$\sqrt N\left(\hat b-\frac12\right)$ est \emph{stable}~: les histogrammes
obtenus pour $N=10,\,100,\,1000,\,10\,000$ se superposent, et convergent vers
une loi normale centrée de variance $\epsvar/\sigma_x^2$. La normalisation par
$\sqrt N$ est exactement celle qui compense la décroissance en $1/N$ de la
variance~: elle rend visible la forme limite de la distribution.

\smallskip

\textbf{Deux remarques.} \emph{(i)} Lorsque la variable explicative est
déterministe et les erreurs gaussiennes, $\hat b$ est \emph{exactement} gaussien
pour tout $N$ (c'est le théorème du cours sur la loi de $\hat{\mathbf b}$)~:
l'histogramme est gaussien dès $N=10$. \emph{(ii)} Lorsque la variable
explicative est stochastique, la loi exacte à distance finie n'est plus
gaussienne, mais le théorème central limite rétablit la normalité
asymptotique~: c'est ce que montre la superposition des histogrammes
renormalisés. L'exercice~\ref{ex:tcl} reprend ce point avec des erreurs non gaussiennes.

\nouvelexercice\label{ex:normale-bivariee}
\setcounter{qnumber}{-1}
On pose $z_i = (x_i-\mu_i)/\sigma_i$ pour $i=1,2$.

\question Le paramètre $\rho$ est le coefficient de corrélation linéaire entre
$X_1$ et $X_2$~: $\rho = \mathrm{cov}(X_1,X_2)/(\sigma_1\sigma_2)$. Il mesure
l'intensité \emph{et} le sens de la dépendance linéaire. Dans le cas gaussien,
et seulement dans ce cas, $\rho=0$ équivaut à l'indépendance.

\question La matrice de variance-covariance et son déterminant valent~:
\[
  \Sigma = \begin{pmatrix}\sigma_1^2 & \rho\sigma_1\sigma_2\\ \rho\sigma_1\sigma_2 & \sigma_2^2\end{pmatrix},
  \qquad
  \det\Sigma = \sigma_1^2\sigma_2^2\left(1-\rho^2\right)
\]
\[
  \Sigma^{-1} = \frac{1}{\sigma_1^2\sigma_2^2(1-\rho^2)}
  \begin{pmatrix}\sigma_2^2 & -\rho\sigma_1\sigma_2\\ -\rho\sigma_1\sigma_2 & \sigma_1^2\end{pmatrix}
\]
La forme quadratique de l'exposant se développe en~:
\[
  Q = (x-\mu)'\Sigma^{-1}(x-\mu) = \frac{1}{1-\rho^2}\left( z_1^2 - 2\rho z_1z_2 + z_2^2 \right)
\]
d'où la densité jointe~:
\[
  f(x_1,x_2) = \frac{1}{2\pi\sigma_1\sigma_2\sqrt{1-\rho^2}}
  \exp\left\{-\frac{1}{2\left(1-\rho^2\right)}\left( z_1^2 - 2\rho z_1z_2 + z_2^2 \right)\right\}
\]

\question On complète le carré en $z_1$. En écrivant
$z_1^2-2\rho z_1z_2+z_2^2 = \left(z_1-\rho z_2\right)^2 + \left(1-\rho^2\right)z_2^2$,
l'exposant devient~:
\[
  -\frac{\left(z_1-\rho z_2\right)^2}{2\left(1-\rho^2\right)} - \frac{z_2^2}{2}
\]
et, comme $e^{A+B}=e^Ae^B$, la densité jointe se factorise~:
\[
  f(x_1,x_2) =
  \underbrace{\frac{1}{\sqrt{2\pi}\,\sigma_1\sqrt{1-\rho^2}}\exp\left\{-\frac{\left(z_1-\rho z_2\right)^2}{2\left(1-\rho^2\right)}\right\}}_{f_{X_1|X_2}(x_1|x_2)}
  \times
  \underbrace{\frac{1}{\sqrt{2\pi}\,\sigma_2}\exp\left\{-\frac{z_2^2}{2}\right\}}_{f_{X_2}(x_2)}
\]
Le second facteur est la densité d'une loi $\normal{\mu_2}{\sigma_2^2}$~: c'est
bien la marginale de $X_2$. Le premier, qui intègre à $1$ en $x_1$, est donc la
densité conditionnelle. En réécrivant
$z_1-\rho z_2 = \frac{1}{\sigma_1}\left(x_1 - \mu_1 - \rho\frac{\sigma_1}{\sigma_2}(x_2-\mu_2)\right)$,
on reconnaît une densité gaussienne. Les deux densités sont donc gaussiennes.

\question De la forme obtenue à la question précédente~:
\[
  \mathbb E\left[X_1|X_2=x_2\right] = \mu_1 + \rho\frac{\sigma_1}{\sigma_2}\left(x_2-\mu_2\right),
  \qquad
  \mathbb V\left[X_1|X_2=x_2\right] = \sigma_1^2\left(1-\rho^2\right)
\]
\textbf{Deux observations.} L'espérance conditionnelle est \emph{linéaire} en
$x_2$~: c'est la propriété qui fait de la régression linéaire le bon modèle dans
le monde gaussien. La variance conditionnelle, elle, ne dépend \emph{pas} de
$x_2$ (homoscédasticité) et est strictement inférieure à $\sigma_1^2$ dès que
$\rho\neq0$~: conditionner, c'est-à-dire disposer de l'information $X_2=x_2$,
réduit l'incertitude sur $X_1$, et d'autant plus que $|\rho|$ est grand.

\question $X_1$ est une combinaison linéaire de deux gaussiennes
indépendantes, donc gaussienne~:
\[
  X_1 = a + bX_2 + \varepsilon \sim \normal{a+b\mu_2}{b^2\sigma_2^2+\epsvar}
\]

\question Conditionnellement à $X_2=x_2$, le terme $a+bx_2$ est déterministe et
seul $\varepsilon$ reste aléatoire~:
\[
  X_1|X_2=x_2 \sim \normal{a+bx_2}{\epsvar}
  \qquad\text{soit}\qquad
  f_{X_1|X_2}(x_1|x_2) = \frac{1}{\sqrt{2\pi}\,\sigma_{\varepsilon}}\exp\left\{-\frac{\left(x_1-a-bx_2\right)^2}{2\epsvar}\right\}
\]
En identifiant avec la question (3), on retrouve
$b = \rho\sigma_1/\sigma_2$, $a = \mu_1-b\mu_2$ et
$\epsvar = \sigma_1^2(1-\rho^2)$~: le modèle de régression linéaire simple
\emph{est} la loi conditionnelle du couple gaussien. Les deux points de vue
coïncident.

\question \textbf{Galton et la régression vers la moyenne.} Notons $X_2$ la
taille des parents et $X_1$ celle des enfants. Si la distribution des tailles
est stationnaire d'une génération à l'autre, alors $\mu_1=\mu_2=\mu$ et
$\sigma_1=\sigma_2$. La question (3) donne alors~:
\[
  \mathbb E\left[X_1|X_2=x_2\right] = \mu + \rho\left(x_2-\mu\right) = (1-\rho)\mu + \rho\,x_2
\]
L'espérance de la taille de l'enfant est une moyenne pondérée de la taille des
parents et de la moyenne de la population, les poids étant $\rho$ et $1-\rho$.
Comme $0<\rho<1$, l'enfant de parents très grands ($x_2>\mu$) est \emph{en
moyenne} plus petit que ses parents, mais reste plus grand que la moyenne~; et
symétriquement pour des parents très petits.

\smallskip

\textbf{Interprétation.} Il n'y a là aucune loi biologique de «~retour à la
médiocrité~»~: le phénomène est un pur artefact statistique de la corrélation
imparfaite. Dès que $\rho<1$, l'espérance conditionnelle est nécessairement plus
proche de la moyenne que la valeur conditionnante. Le raisonnement est d'ailleurs
symétrique --- les parents d'enfants très grands sont eux aussi, en moyenne,
plus petits que leurs enfants --- ce qui interdit toute lecture causale. C'est de
cette étude que le mot «~régression~» est resté attaché à la méthode.

\nouvelexercice\label{ex:r2}

\question La première équation normale donne $\sum_t\hat\epsilon_t = 0$
(exercice~\ref{ex:mco-matriciel}), donc $\sum_t \hat y_t = \sum_t y_t$ et $\bar{\hat y} = \bar y$. En
écrivant $y_t-\bar y = (\hat y_t-\bar y) + \hat\epsilon_t$ et en élevant au
carré~:
\[
  SST = \sum_t (\hat y_t-\bar y)^2 + 2\sum_t(\hat y_t-\bar y)\hat\epsilon_t + \sum_t\hat\epsilon_t^{\,2}
\]
Le double produit est nul~:
$\sum_t \hat y_t\hat\epsilon_t = (X\hat{\mathbf b})'\hat\epsilon = \hat{\mathbf b}'\underbrace{X'\hat\epsilon}_{=0} = 0$
et $\bar y\sum_t\hat\epsilon_t = 0$. D'où $SST = SSR+SSE$. \textbf{La constante
intervient deux fois}~: pour garantir $\sum_t\hat\epsilon_t=0$ (donc
$\bar{\hat y}=\bar y$) et pour annuler le second terme du double produit.

\question $R^2 = 1-SSE/SST = (SST-SSE)/SST = SSR/SST$. Comme $SSR\geq0$ et
$SSE\geq0$ et que $SST = SSR+SSE$, on a $0\leq SSR\leq SST$, d'où
$0\leq R^2\leq 1$.

\question En utilisant $y_t-\bar y = (\hat y_t-\bar y)+\hat\epsilon_t$ et
l'orthogonalité~:
\[
  \sum_t (y_t-\bar y)(\hat y_t-\bar y) = \sum_t (\hat y_t-\bar y)^2 = SSR
\]
donc~:
\[
  \mathrm{corr}\left(\mathbf y,\hat{\mathbf y}\right)^2
  = \frac{\left(\sum_t (y_t-\bar y)(\hat y_t-\bar y)\right)^2}{\sum_t (y_t-\bar y)^2\sum_t(\hat y_t-\bar y)^2}
  = \frac{SSR^2}{SST\cdot SSR} = \frac{SSR}{SST} = R^2
\]
Le coefficient de détermination est donc le carré de la corrélation entre la
variable expliquée et sa prédiction.

\question Notons $SSE_K$ le minimum de la somme des carrés des résidus avec $K$
explicatives, et $SSE_{K+1}$ celui obtenu en ajoutant une variable. L'espace des
combinaisons linéaires engendré par $K+1$ colonnes \emph{contient} celui
engendré par les $K$ premières~: le minimum sur le plus grand ensemble ne peut
pas être plus élevé, donc $SSE_{K+1}\leq SSE_K$. Comme $SST$ ne dépend pas du
modèle, $R^2$ ne peut pas diminuer. \textbf{Conséquence pratique}~: le $R^2$ ne
peut pas servir à choisir entre deux modèles emboîtés, puisqu'il désigne
toujours le plus grand.

\question Comme $K\geq1$, on a $\frac{T-1}{T-K}\geq 1$, donc
$(1-R^2)\frac{T-1}{T-K}\geq 1-R^2$ et $\bar R^2\leq R^2$. Le $\bar R^2$ devient
négatif dès que $R^2 < \frac{K-1}{T-1}$, c'est-à-dire lorsque le modèle explique
moins qu'un modèle réduit à la constante, compte tenu du nombre de paramètres.
\textbf{Ce qu'il corrige}~: ajouter une variable fait mécaniquement croître
$R^2$ (question 4) mais aussi le facteur de pénalité $\frac{T-1}{T-K}$~; le
$\bar R^2$ n'augmente donc que si le gain d'ajustement compense le paramètre
consommé.

\question Sans constante, la première équation normale disparaît~:
$\sum_t\hat\epsilon_t\neq0$ en général, donc $\bar{\hat y}\neq\bar y$ et le
double produit de la question (1) ne s'annule plus. La décomposition
$SST = SSR+SSE$ est fausse. Il en résulte que $1-SSE/SST$ et $SSR/SST$ ne
coïncident plus, et que le premier peut devenir \emph{négatif} (rien ne borne
$SSE$ par $SST$). Le $R^2$ perd alors toute interprétation en «~part de variance
expliquée~», et les $R^2$ de deux modèles dont l'un a une constante et l'autre
non ne sont pas comparables.

\question Les quatre jeux d'Anscombe partagent moyennes, variances, corrélation,
droite de régression et $R^2$, alors que le premier est linéaire avec bruit, le
deuxième parfaitement quadratique, le troisième linéaire avec un point aberrant,
et le quatrième entièrement déterminé par un unique point à fort levier.
\textbf{Conclusion}~: le $R^2$ résume la qualité de l'ajustement \emph{linéaire}
en un seul nombre et ne dit rien de la validité de la spécification. Il ne
dispense ni de tracer les résidus, ni d'examiner les points influents.

\nouvelexercice\label{ex:s2}

\question En substituant le DGP dans $\hat\epsilon = M\mathbf y$ et en utilisant
$MX=0$ (exercice~\ref{ex:fwl})~:
\[
  \hat\epsilon = M\left(X\beta+\varepsilon\right) = \underbrace{MX}_{=0}\beta + M\varepsilon = M\varepsilon
\]
$M$ est symétrique et idempotente. Sa trace vaut~:
\[
  \trace(M) = \trace(I_T) - \tracarg{X(X'X)^{-1}X'} = T - \tracarg{(X'X)^{-1}X'X} = T-\trace(I_K) = T-K
\]
en utilisant $\trace(AB)=\trace(BA)$. Pour une matrice idempotente, trace et
rang coïncident~: $\mathrm{rang}(M)=T-K$.

\question Comme $M$ est symétrique idempotente,
$SSE = \hat\epsilon'\hat\epsilon = \varepsilon'M'M\varepsilon = \varepsilon'M\varepsilon$.
Donc~:
\[
  \mathbb E\left[SSE\right] = \tracarg{M\,\mathbb V[\varepsilon]} = \tracarg{M\,\epsvar I_T} = \epsvar\trace(M) = \epsvar(T-K)
\]
d'où $\mathbb E[s^2] = \mathbb E[SSE]/(T-K) = \epsvar$~: l'estimateur $s^2$ est
sans biais. \textbf{Remarque.} La division par $T-K$ et non par $T$ s'interprète
ainsi~: les résidus estimés vivent dans un sous-espace de dimension $T-K$, les
$K$ équations normales ayant consommé $K$ degrés de liberté.

\question $M$ étant symétrique idempotente de trace $T-K$~:
\[
  \frac{(T-K)s^2}{\epsvar} = \frac{\varepsilon'M\varepsilon}{\epsvar} \sim \chi^2(T-K)
\]

\question La variance d'une loi du khi-deux à $k$ degrés de liberté vaut $2k$,
donc $\mathbb V\left[(T-K)s^2/\epsvar\right] = 2(T-K)$, soit~:
\[
  \mathbb V\left[s^2\right] = \frac{2\sigma_{\varepsilon}^4}{T-K}
\]

\question On a $\hat{\mathbf b}-\beta = (X'X)^{-1}X'\varepsilon = B\varepsilon$
avec $B = (X'X)^{-1}X'$, et $SSE = \varepsilon'M\varepsilon$. Or~:
\[
  BM = (X'X)^{-1}X'M = (X'X)^{-1}\left(MX\right)' = 0
\]
puisque $M$ est symétrique et $MX=0$. Pour un vecteur gaussien, $B\varepsilon$
et $\varepsilon'M\varepsilon$ sont donc indépendants. Comme $s^2$ est fonction
de $SSE$ seul, $\hat{\mathbf b}$ et $s^2$ sont indépendants.

\question Une loi de Student à $k$ degrés de liberté est le rapport
$Z/\sqrt{U/k}$ où $Z\sim\normal{0}{1}$, $U\sim\chi^2(k)$ et $Z\perp U$. Or~:
\[
  \frac{\hat{\mathbf b}_i-\beta_i}{\sqrt{\epsvar\left[(X'X)^{-1}\right]_{ii}}}\sim\normal{0}{1},
  \qquad
  \frac{(T-K)s^2}{\epsvar}\sim\chi^2(T-K)
\]
et ces deux quantités sont indépendantes par la question (5). En formant le
rapport, $\epsvar$ disparaît~:
\[
  t = \frac{\hat{\mathbf b}_i-\beta_i}{\sqrt{s^2\left[(X'X)^{-1}\right]_{ii}}} \sim t(T-K)
\]
C'est précisément parce que $\epsvar$ est inconnu et remplacé par $s^2$ que la
loi n'est pas normale mais de Student.

\question $s^2$ est sans biais et
$\mathbb V[s^2] = 2\sigma_{\varepsilon}^4/(T-K)\rightarrow0$ lorsque
$T\rightarrow\infty$ ($K$ fixe). L'erreur quadratique moyenne tend donc vers
zéro. Or un estimateur sans biais dont la variance tend vers zéro se concentre
nécessairement autour de sa cible~: $s^2\plimT\epsvar$. (C'est l'inégalité de
Bienaymé-Tchebychev~: la probabilité de s'écarter de $\epsvar$ de plus de
$\delta$ est majorée par $\mathbb V[s^2]/\delta^2$, qui tend vers zéro.)

\nouvelexercice\label{ex:vraisemblance}

\question Sous $\varepsilon\sim\normal{0}{\epsvar I_T}$, le vecteur $\mathbf y$
est gaussien d'espérance $X\beta$ et de variance $\epsvar I_T$~:
\[
  L\left(\beta,\epsvar;\mathbf y\right)
  = \left(2\pi\epsvar\right)^{-T/2}\exp\left\{-\frac{\left(\mathbf y-X\beta\right)'\left(\mathbf y-X\beta\right)}{2\epsvar}\right\}
\]
\[
  \ell\left(\beta,\epsvar;\mathbf y\right)
  = -\frac{T}{2}\log(2\pi) - \frac{T}{2}\log\epsvar - \frac{1}{2\epsvar}\left(\mathbf y-X\beta\right)'\left(\mathbf y-X\beta\right)
\]

\question Pour $\epsvar$ fixé, maximiser $\ell$ en $\beta$ revient à
\emph{minimiser} $(\mathbf y-X\beta)'(\mathbf y-X\beta)$~: c'est exactement le
programme des MCO. D'où $\hat\beta_{\textsc{mv}} = (X'X)^{-1}X'\mathbf y = \hat{\mathbf b}$.
Puis~:
\[
  \frac{\partial\ell}{\partial\epsvar} = -\frac{T}{2\epsvar} + \frac{SSE}{2\sigma_{\varepsilon}^4} = 0
  \quad\Longleftrightarrow\quad
  \hat\sigma^2_{\varepsilon} = \frac{SSE}{T}
\]
\textbf{Remarque.} L'équivalence MCO/MV n'a rien d'automatique~: elle tient à la
normalité. Les MCO ne supposent pas de loi, le MV en suppose une.

\question D'après l'exercice~\ref{ex:s2}, $\mathbb E[SSE]=\epsvar(T-K)$, donc~:
\[
  \mathbb E\left[\hat\sigma^2_{\varepsilon}\right] = \frac{T-K}{T}\epsvar = \epsvar - \frac{K}{T}\epsvar
\]
L'estimateur du maximum de vraisemblance est \textbf{biaisé}, et sous-estime
systématiquement $\epsvar$. Le biais, $-K\epsvar/T$, tend vers zéro quand
$T\rightarrow\infty$~: l'estimateur est asymptotiquement sans biais.

\question En dérivant une seconde fois et en prenant l'opposé de l'espérance~:
\[
  \mathcal I(\theta) =
  \begin{pmatrix}
    \dfrac{X'X}{\epsvar} & 0\\[2mm]
    0 & \dfrac{T}{2\sigma_{\varepsilon}^4}
  \end{pmatrix}
\]
Le bloc croisé est nul car
$\mathbb E\left[\partial^2\ell/\partial\beta\,\partial\epsvar\right]
= -\mathbb E\left[X'\varepsilon\right]/\sigma_{\varepsilon}^4 = 0$.
\textbf{Interprétation.} La bloc-diagonalité signifie que les deux paramètres
sont \emph{orthogonaux} au sens de l'information~: connaître ou ignorer
$\epsvar$ ne change pas la précision asymptotique avec laquelle on estime
$\beta$, et réciproquement. On retrouvera cette propriété au chapitre~IV.

\question Les bornes de Cramer-Darmois-Fréchet-Rao sont les blocs diagonaux de
$\mathcal I(\theta)^{-1}$~:
\[
  \epsvar (X'X)^{-1}
  \qquad\text{et}\qquad
  \frac{2\sigma_{\varepsilon}^4}{T}
\]
L'estimateur des MCO de $\beta$ a exactement pour variance $\epsvar(X'X)^{-1}$~:
il \textbf{atteint} la borne, il est donc efficace parmi \emph{tous} les
estimateurs sans biais (et pas seulement parmi les estimateurs linéaires, comme
l'assure Gauss-Markov). En revanche
$\mathbb V[s^2] = 2\sigma_{\varepsilon}^4/(T-K) > 2\sigma_{\varepsilon}^4/T$~:
$s^2$ n'atteint pas sa borne à distance finie, mais s'en approche
asymptotiquement.

\question Posons $m = T-K$. On a $\mathbb E[SSE] = m\epsvar$ et
$\mathbb V[SSE] = 2m\sigma_{\varepsilon}^4$ (loi du khi-deux). Pour
$\sigma^2_{(c)} = SSE/c$~:
\[
  \frac{\mathrm{EQM}\left(\sigma^2_{(c)}\right)}{\sigma_{\varepsilon}^4}
  = \underbrace{\frac{2m}{c^2}}_{\text{variance}} + \underbrace{\left(\frac{m}{c}-1\right)^2}_{\text{biais}^2}
  = \frac{2m+m^2}{c^2} - \frac{2m}{c} + 1
\]
En annulant la dérivée par rapport à $c$~:
\[
  -\frac{2\left(2m+m^2\right)}{c^3} + \frac{2m}{c^2} = 0
  \quad\Longleftrightarrow\quad
  c = \frac{2m+m^2}{m} = m+2 = T-K+2
\]
La dérivée seconde est positive en ce point~: il s'agit bien d'un minimum.

\question Ni $s^2$ ($c=T-K$) ni $\hat\sigma^2_{\varepsilon}$ ($c=T$) ne sont
optimaux au sens de l'erreur quadratique moyenne, puisque l'optimum est
$c = T-K+2$. En comparant les deux~:
\[
  \mathrm{EQM}\left(\hat\sigma^2_{\varepsilon}\right) < \mathrm{EQM}\left(s^2\right)
  \quad\Longleftrightarrow\quad
  (T-K)(K-4) < 2K
\]
condition satisfaite en particulier pour $K\leq4$, quelle que soit la taille de
l'échantillon. \textbf{Interprétation.} L'absence de biais n'est pas un
critère d'optimalité~: l'estimateur du maximum de vraisemblance, pourtant
biaisé, bat souvent l'estimateur sans biais en erreur quadratique moyenne, parce
qu'il échange un peu de biais contre une réduction de variance. C'est
l'arbitrage biais-variance, dont on retrouvera une autre illustration à
l'exercice~\ref{ex:contraintes}.

\nouvelexercice\label{ex:tests}

\question Sous les conditions idéales,
$\hat{\mathbf b}\sim\normal{\beta}{\epsvar(X'X)^{-1}}$. Comme $R$ est
déterministe, $R\hat{\mathbf b}$ est gaussien, et sous
$\mathcal H_0: R\beta = r$~:
\[
  R\hat{\mathbf b}-r \sim \normal{0}{\epsvar R(X'X)^{-1}R'}
\]

\question Pour $m=1$, $R(X'X)^{-1}R'$ est un scalaire strictement positif et~:
\[
  \frac{R\hat{\mathbf b}-r}{\sqrt{\epsvar R(X'X)^{-1}R'}}\sim\normal{0}{1}
\]
Cette statistique n'est pas utilisable~: elle dépend de $\epsvar$, inconnu. On le
remplace par $s^2$. Par l'exercice~\ref{ex:s2}, $(T-K)s^2/\epsvar\sim\chi^2(T-K)$
indépendamment de $\hat{\mathbf b}$, de sorte que~:
\[
  t = \frac{R\hat{\mathbf b}-r}{\sqrt{s^2 R(X'X)^{-1}R'}} \sim t(T-K)
\]

\question Il suffit de prendre $R = (0,\dots,0,1,0,\dots,0)$, le $1$ étant en
$i$-ème position, et $r=0$. Alors $R(X'X)^{-1}R' = \left[(X'X)^{-1}\right]_{ii}$
et~:
\[
  t = \frac{\hat{\mathbf b}_i}{\sqrt{s^2\left[(X'X)^{-1}\right]_{ii}}} = \frac{\hat{\mathbf b}_i}{s_{\hat{\mathbf b}_i}}
\]
C'est la statistique de Student imprimée par les logiciels.

\question Sous $\mathcal H_0$, le vecteur $R\hat{\mathbf b}-r$ est gaussien
centré de variance $\epsvar R(X'X)^{-1}R'$, donc la forme quadratique~:
\[
  \frac{\left(R\hat{\mathbf b}-r\right)'\left(R(X'X)^{-1}R'\right)^{-1}\left(R\hat{\mathbf b}-r\right)}{\epsvar}\sim\chi^2(m)
\]
Elle est indépendante de $(T-K)s^2/\epsvar\sim\chi^2(T-K)$. Le rapport de deux
khi-deux indépendants divisés par leurs degrés de liberté suit une loi de
Fisher~; $\epsvar$ se simplifie~:
\[
  F \sim \mathcal F(m,\,T-K)
\]

\question Pour $m=1$, la forme quadratique se réduit à
$\left(R\hat{\mathbf b}-r\right)^2/\left(R(X'X)^{-1}R'\right)$ et $m=1$ au
dénominateur, de sorte que $F = t^2$. C'est cohérent avec le fait qu'une loi
$\mathcal F(1,k)$ est la loi du carré d'une $t(k)$. Les deux tests sont donc
équivalents pour une restriction unique --- à ceci près que le test de Student
permet une alternative unilatérale, ce que le test de Fisher ne permet pas.

\question On prend $R = \left(0\ \vdots\ I_{K-1}\right)$ et $r=0$, soit $m=K-1$.
Le modèle contraint se réduit à la seule constante, dont l'estimateur est
$\bar y$~: sa somme des carrés des résidus vaut donc $SST$. En utilisant la
formulation du test par comparaison des sommes de carrés~:
\[
  F = \frac{\left(SST-SSE\right)/(K-1)}{SSE/(T-K)}
    = \frac{SSR/(K-1)}{SSE/(T-K)}
\]
En divisant numérateur et dénominateur par $SST$ et en utilisant
$R^2 = SSR/SST$ et $1-R^2 = SSE/SST$~:
\[
  F = \frac{R^2}{1-R^2}\,\frac{T-K}{K-1}
\]

\question \textbf{Oui}, et c'est une situation classique. Elle survient lorsque
les explicatives sont fortement corrélées entre elles~: chaque coefficient est
alors estimé très imprécisément (variances individuelles élevées), de sorte
qu'aucun $t$ n'est significatif~; mais les variables expliquent conjointement
une part importante de la variance de $\mathbf y$, et le test joint rejette
nettement. Techniquement, le test de Fisher tient compte des \emph{covariances}
entre estimateurs, que les tests de Student pris isolément ignorent. C'est le
symptôme le plus caractéristique de la multicolinéarité (exercice~\ref{ex:multicolinearite}).

\nouvelexercice\label{ex:contraintes} On note $P = (X'X)^{-1}$ pour alléger.

\question Le lagrangien s'écrit~:
\[
  \mathcal L(\mathbf b,\bm\lambda) = \left(\mathbf y-X\mathbf b\right)'\left(\mathbf y-X\mathbf b\right) + \bm\lambda'\left(R\mathbf b-r\right)
\]
et les conditions du premier ordre~:
\[
  \begin{cases}
    \dfrac{\partial\mathcal L}{\partial\mathbf b} = -2X'\left(\mathbf y-X\tilde{\mathbf b}\right) + R'\bm\lambda = 0\\[3mm]
    \dfrac{\partial\mathcal L}{\partial\bm\lambda} = R\tilde{\mathbf b}-r = 0
  \end{cases}
\]

\question La première condition donne
$2X'X\tilde{\mathbf b} = 2X'\mathbf y - R'\bm\lambda$, soit~:
\[
  \tilde{\mathbf b} = \hat{\mathbf b} - \tfrac{1}{2}PR'\bm\lambda
\]
En prémultipliant par $R$ et en imposant la contrainte $R\tilde{\mathbf b}=r$~:
\[
  r = R\hat{\mathbf b} - \tfrac{1}{2}RPR'\bm\lambda
  \quad\Longleftrightarrow\quad
  \tfrac{1}{2}\bm\lambda = \left(RPR'\right)^{-1}\left(R\hat{\mathbf b}-r\right)
\]
En reportant~:
\[
  \tilde{\mathbf b} = \hat{\mathbf b} - PR'\left(RPR'\right)^{-1}\left(R\hat{\mathbf b}-r\right)
\]
L'estimateur contraint corrige l'estimateur libre proportionnellement à la
violation observée de la contrainte.

\question En prenant l'espérance et en utilisant
$\mathbb E[\hat{\mathbf b}]=\beta$~:
\[
  \mathbb E\left[\tilde{\mathbf b}\right] = \beta - PR'\left(RPR'\right)^{-1}\left(R\beta-r\right)
\]
L'estimateur contraint est sans biais \textbf{si et seulement si} la contrainte
est vraie, c'est-à-dire $R\beta=r$. Sinon le biais est proportionnel à l'erreur
de spécification $R\beta-r$.

\question Posons $A = PR'\left(RPR'\right)^{-1}R$. On a
$\tilde{\mathbf b}-\mathbb E[\tilde{\mathbf b}] = (I-A)\left(\hat{\mathbf b}-\beta\right)$,
donc $\mathbb V[\tilde{\mathbf b}] = \epsvar (I-A)P(I-A)'$. En développant, et
en remarquant que $AP = PR'(RPR')^{-1}RP$ est symétrique et que
$APA' = AP$~:
\[
  (I-A)P(I-A)' = P - AP - PA' + APA' = P - PR'\left(RPR'\right)^{-1}RP
\]
d'où~:
\[
  \mathbb V\left[\hat{\mathbf b}\right]-\mathbb V\left[\tilde{\mathbf b}\right]
  = \epsvar\, PR'\left(RPR'\right)^{-1}RP
\]
Cette matrice est semi-définie positive~: pour tout vecteur $u$, en posant
$w = RPu$, la forme quadratique vaut $\epsvar\,w'(RPR')^{-1}w\geq0$ puisque
$RPR'$ est définie positive. La contrainte réduit donc toujours la variance.

\question L'information apportée par la contrainte est utilisée par
l'estimateur, ce qui réduit sa variance --- \emph{que la contrainte soit vraie
ou fausse}, puisque le calcul de la question (4) ne fait pas intervenir
$R\beta-r$. Si la contrainte est fausse, on introduit un biais. L'erreur
quadratique moyenne de $\tilde{\mathbf b}$ peut néanmoins rester inférieure à
celle de $\hat{\mathbf b}$ lorsque le biais est petit devant le gain de
variance. C'est le même arbitrage qu'à l'exercice~\ref{ex:vraisemblance}~: une information
approximative vaut parfois mieux que pas d'information du tout.

\question Les rendements d'échelle sont constants si $\beta_2+\beta_3=1$, soit~:
\[
  R = \left(0\quad 1\quad 1\right), \qquad r = 1, \qquad m=1
\]
Sous la contrainte, $\beta_3 = 1-\beta_2$ et~:
\[
  \log Q_t = \beta_1 + \beta_2\log K_t + \left(1-\beta_2\right)\log L_t + \varepsilon_t
\]
soit, en retranchant $\log L_t$ des deux côtés~:
\[
  \log\frac{Q_t}{L_t} = \beta_1 + \beta_2\log\frac{K_t}{L_t} + \varepsilon_t
\]
La contrainte se met donc en œuvre par une simple reparamétrisation~: on régresse
la productivité du travail sur l'intensité capitalistique. On estime alors deux
paramètres au lieu de trois, et $\beta_3$ se déduit par $1-\hat\beta_2$.

\question Comme $m=1$, on peut tester par un $t$ de Student sur $\hat\beta_2+\hat\beta_3-1$,
ou de façon équivalente par un $F$ à $(1,T-3)$ degrés de liberté, les deux étant
liés par $F=t^2$ (question 5 de l'exercice~\ref{ex:tests}). En pratique, on compare les
sommes des carrés des résidus des modèles contraint et non contraint~:
\[
  F = \frac{\left(SSE_{\text{c}}-SSE\right)/m}{SSE/(T-K)}
\]
\textbf{Lien avec la vraisemblance.} Le cours établit que cette statistique de
Fisher \emph{est} un test du rapport des vraisemblances~: sous normalité, la
log-vraisemblance concentrée est une fonction décroissante de la somme des
carrés des résidus, de sorte que comparer $SSE_{\text{c}}$ et $SSE$ revient à
comparer les vraisemblances maximisées sous $\mathcal H_0$ et sous
$\mathcal H_1$. Le test de Fisher n'est donc pas un test \emph{ad hoc}~: il
s'inscrit dans la théorie générale du maximum de vraisemblance.

\nouvelexercice\label{ex:indicatrice}

\question La matrice $X$ a deux colonnes~: un vecteur de uns et le vecteur des
$d_i$. Comme $d_i^2 = d_i$ (une indicatrice est égale à son carré) et que
$\sum_i d_i = N_1$~:
\[
  X'X = \begin{pmatrix} N & N_1\\ N_1 & N_1\end{pmatrix},
  \qquad
  X'\mathbf y = \begin{pmatrix} \sum_i y_i\\ \sum_{i\in 1} y_i\end{pmatrix}
  = \begin{pmatrix} N_0\bar y_0 + N_1\bar y_1\\ N_1\bar y_1\end{pmatrix}
\]

\question Le déterminant vaut $NN_1-N_1^2 = N_1N_0$, d'où~:
\[
  (X'X)^{-1} = \frac{1}{N_0N_1}\begin{pmatrix} N_1 & -N_1\\ -N_1 & N\end{pmatrix}
\]
et en effectuant le produit~:
\[
  \hat{\mathbf b} = (X'X)^{-1}X'\mathbf y
  = \begin{pmatrix} \bar y_0\\ \bar y_1-\bar y_0\end{pmatrix}
\]
\textbf{Interprétation.} La constante estime la moyenne du groupe de référence,
et le coefficient de l'indicatrice l'\emph{écart} entre les deux moyennes. Une
régression sur une indicatrice n'est donc rien d'autre qu'une comparaison de
deux moyennes --- ce qui donne une lecture très concrète des MCO.

\question La prédiction vaut $\hat y_i = \hat b_1 = \bar y_0$ si $d_i=0$ et
$\hat y_i = \hat b_1+\hat b_2 = \bar y_1$ si $d_i=1$~: elle ne prend que deux
valeurs, les deux moyennes de groupe. Les résidus sont donc
$\hat\epsilon_i = y_i - \bar y_{g(i)}$ et~:
\[
  SSE = \sum_{g\in\{0,1\}}\ \sum_{i\in g}\left(y_i-\bar y_g\right)^2
\]
C'est la somme des carrés \emph{intra-groupes}.

\question Ici $K=2$, donc~:
\[
  s^2 = \frac{SSE}{N-2} = \frac{\sum_{i\in0}(y_i-\bar y_0)^2 + \sum_{i\in1}(y_i-\bar y_1)^2}{N_0+N_1-2}
\]
On reconnaît exactement l'estimateur de variance «~poolée~» du test de
comparaison de deux moyennes.

\question La variance de $\hat b_2$ est $s^2$ multiplié par l'élément $(2,2)$
de $(X'X)^{-1}$, soit $N/(N_0N_1)$~:
\[
  s_{\hat b_2}^2 = s^2\,\frac{N}{N_0N_1} = s^2\left(\frac{1}{N_0}+\frac{1}{N_1}\right)
\]
puisque $N/(N_0N_1) = (N_0+N_1)/(N_0N_1) = 1/N_0+1/N_1$. La statistique de
Student s'écrit donc~:
\[
  t = \frac{\hat b_2}{s_{\hat b_2}} = \frac{\bar y_1-\bar y_0}{s\sqrt{\dfrac{1}{N_0}+\dfrac{1}{N_1}}}\ \sim\ t(N-2)
\]
\textbf{C'est mot pour mot la statistique du test de Student de comparaison de
deux moyennes à variances égales}, degrés de liberté compris. Le test de
significativité d'un coefficient et le test de comparaison de moyennes sont donc
le même test~: les MCO généralisent un outil déjà connu.

\question Comme $SST = \sum_i (y_i-\bar y)^2$ et que $SSE$ est la somme
intra-groupes, $SSR = SST-SSE$ est la somme \emph{inter-groupes}~:
\[
  R^2 = \frac{N_0\left(\bar y_0-\bar y\right)^2 + N_1\left(\bar y_1-\bar y\right)^2}{\sum_i\left(y_i-\bar y\right)^2}
      = \frac{N_0N_1}{N}\,\frac{\left(\bar y_1-\bar y_0\right)^2}{\sum_i\left(y_i-\bar y\right)^2}
\]
(en utilisant $\bar y = (N_0\bar y_0+N_1\bar y_1)/N$). Le $R^2$ mesure ici la
part de la variance totale imputable à l'écart entre groupes~: c'est le rapport
de corrélation de l'analyse de la variance.

\question La somme des deux indicatrices vaut $d_i+e_i = 1$ pour tout $i$~:
elles reconstituent exactement la colonne de constante. La matrice $X$ n'est
donc pas de plein rang, $X'X$ n'est pas inversible, et l'estimateur des MCO
n'est pas défini~: c'est un cas de \textbf{colinéarité parfaite}, l'hypothèse
$\mathcal H_2$ étant violée (exercice~\ref{ex:multicolinearite}, question 5). On
l'appelle le «~piège des variables indicatrices~». Le remède est d'omettre soit
la constante, soit l'une des deux indicatrices~: on choisit alors un groupe de
référence, et les coefficients s'interprètent en écart à ce groupe. Avec $G$
modalités, on met $G-1$ indicatrices.

\nouvelexercice\label{ex:regression-inverse}

\question En notant $S_{xy} = \sum_i (x_i-\bar x)(y_i-\bar y)$,
$S_{xx} = \sum_i (x_i-\bar x)^2$ et $S_{yy} = \sum_i (y_i-\bar y)^2$~:
\[
  \hat b_{y|x} = \frac{S_{xy}}{S_{xx}},
  \qquad
  \hat b_{x|y} = \frac{S_{xy}}{S_{yy}}
\]
Les deux pentes ont le même numérateur mais des dénominateurs différents.

\question Le produit vaut~:
\[
  \hat b_{y|x}\times\hat b_{x|y} = \frac{S_{xy}^2}{S_{xx}S_{yy}} = r_{xy}^2 = R^2
\]
puisque dans une régression simple le coefficient de détermination est le carré
du coefficient de corrélation empirique (exercice~\ref{ex:r2}, question 3). Les
deux pentes ont donc le même signe, et leur produit est compris entre $0$ et
$1$.

\question Si l'on trace les deux droites dans le plan $(x,y)$, la première a
pour pente $\hat b_{y|x}$ et la seconde $1/\hat b_{x|y}$. Elles coïncident si et
seulement si $\hat b_{y|x} = 1/\hat b_{x|y}$, c'est-à-dire
$\hat b_{y|x}\hat b_{x|y} = 1$, soit $R^2=1$~: les points sont alors exactement
alignés. Par ailleurs les deux droites passent par $(\bar x,\bar y)$, puisque
chacune vérifie $\hat b_0 = \bar y - \hat b_1\bar x$ (respectivement l'analogue
en échangeant les rôles)~: elles se coupent donc toujours au point moyen.

\question Comme $\hat b_{y|x}\hat b_{x|y} = R^2\leq1$ et que les deux pentes
sont positives~:
\[
  \hat b_{y|x} = \frac{R^2}{\hat b_{x|y}} \leq \frac{1}{\hat b_{x|y}}
\]
\textbf{Interprétation.} La droite des MCO de $y$ sur $x$ est toujours moins
inclinée que celle qu'on obtiendrait en régressant $x$ sur $y$ puis en
retournant l'axe. La raison est que les MCO minimisent les écarts \emph{dans la
direction de la variable expliquée}~: le critère n'est pas symétrique, et
changer de variable expliquée change le problème.

\question C'est exactement le mécanisme de Galton
(exercice~\ref{ex:normale-bivariee}). Là-bas, avec $\sigma_1=\sigma_2$, la
pente valait $\rho<1$ dans un sens comme dans l'autre~: la taille des enfants
«~régresse~» vers la moyenne quand on part de celle des parents, mais aussi
l'inverse. L'apparente contradiction se dissipe dès qu'on comprend qu'il s'agit
de deux régressions différentes, et non d'une relation réversible. Ici, le
produit des deux pentes vaut $R^2 = \rho^2 < 1$, ce qui est la même chose.

\question \textbf{Réponse.} Ni l'un ni l'autre n'est «~le bon~» par pure
algèbre~: les deux calculs répondent à des questions différentes, et l'écart
entre eux est d'autant plus grand que le $R^2$ est faible. Le choix ne se
tranche pas statistiquement mais par un argument \emph{économique}~: il faut
savoir quelle variable est déterminée par l'autre, et laquelle porte la
perturbation. Ici, ni l'un ni l'autre ne convient d'ailleurs, puisque prix et
quantité sont déterminés \emph{simultanément} --- c'est très exactement le
problème de l'exercice~\ref{ex:offre-demande}, qui appelle des variables
instrumentales et non un choix de sens de régression.

\question Le sens de la régression cesse d'être arbitraire dès que l'on dispose
d'un \emph{modèle} précisant qui cause quoi et où se loge l'aléa. Deux
situations typiques~: \emph{(i)} une variable est déterminée en amont
(exogène) et l'autre en aval --- il faut alors régresser l'aval sur l'amont~;
\emph{(ii)} l'une des deux variables est mesurée avec erreur. Dans ce dernier
cas, l'exercice~\ref{ex:erreurs-mesure} montre que la régression de $y$ sur un
$x$ mal mesuré atténue la pente vers zéro, tandis que la régression inverse
l'exagère~: la vraie pente est \emph{encadrée} par les deux estimations. C'est
d'ailleurs un usage classique de ce résultat~: à défaut d'instrument, les deux
régressions fournissent un intervalle.

\nouvelexercice\label{ex:leave-one-out}

\question \textit{Démonstration.} Il suffit de vérifier que le membre de droite
est bien l'inverse. En posant $c = 1-\mathbf u'A^{-1}\mathbf u$ et en multipliant~:
\[
  \begin{split}
    \left(A-\mathbf u\mathbf u'\right)\left(A^{-1}+\frac{A^{-1}\mathbf u\mathbf u'A^{-1}}{c}\right)
    &= I + \frac{\mathbf u\mathbf u'A^{-1}}{c} - \mathbf u\mathbf u'A^{-1} - \frac{\mathbf u\left(\mathbf u'A^{-1}\mathbf u\right)\mathbf u'A^{-1}}{c}\\
    &= I + \mathbf u\mathbf u'A^{-1}\left(\frac{1}{c} - 1 - \frac{\mathbf u'A^{-1}\mathbf u}{c}\right)
  \end{split}
\]
et le terme entre parenthèses vaut
$\frac{1-c-\mathbf u'A^{-1}\mathbf u}{c} = \frac{1-c-(1-c)}{c} = 0$. D'où le
résultat. (C'est le cas particulier de rang un de la formule d'inversion par
blocs~: le complément de Schur se réduit ici au scalaire $c$.)

\question Retirer l'observation $t$ revient à retrancher sa contribution à la
matrice des moments~: $X_{(t)}'X_{(t)} = X'X - \mathbf x_t\mathbf x_t'$. En
appliquant la formule avec $A = X'X$ et $\mathbf u = \mathbf x_t$, et en
reconnaissant $h_{tt} = \mathbf x_t'(X'X)^{-1}\mathbf x_t$~:
\[
  \left(X_{(t)}'X_{(t)}\right)^{-1} = (X'X)^{-1} + \frac{(X'X)^{-1}\mathbf x_t\mathbf x_t'(X'X)^{-1}}{1-h_{tt}}
\]

\question De même $X_{(t)}'\mathbf y_{(t)} = X'\mathbf y - \mathbf x_t y_t$. En
posant $P = (X'X)^{-1}$ et en développant le produit~:
\[
  \begin{split}
    \hat{\mathbf b}_{(t)}
    &= \left(P + \frac{P\mathbf x_t\mathbf x_t'P}{1-h_{tt}}\right)\left(X'\mathbf y - \mathbf x_ty_t\right)\\
    &= \hat{\mathbf b} - P\mathbf x_ty_t + \frac{P\mathbf x_t\mathbf x_t'\hat{\mathbf b}}{1-h_{tt}} - \frac{P\mathbf x_t h_{tt}y_t}{1-h_{tt}}
  \end{split}
\]
En regroupant les trois derniers termes sur le dénominateur commun $1-h_{tt}$,
le numérateur devient
$P\mathbf x_t\left[-y_t(1-h_{tt}) + \mathbf x_t'\hat{\mathbf b} - h_{tt}y_t\right]
= P\mathbf x_t\left[\mathbf x_t'\hat{\mathbf b} - y_t\right] = -P\mathbf x_t\hat\epsilon_t$,
d'où~:
\[
  \hat{\mathbf b}_{(t)} = \hat{\mathbf b} - \frac{(X'X)^{-1}\mathbf x_t\,\hat\epsilon_t}{1-h_{tt}}
\]
\textbf{Remarque pratique.} Cette formule permet de calculer les $T$ estimations
«~sans l'observation $t$~» sans jamais réestimer le modèle~: une seule inversion
suffit.

\question En reportant~:
\[
  y_t - \mathbf x_t'\hat{\mathbf b}_{(t)}
  = y_t - \mathbf x_t'\hat{\mathbf b} + \frac{h_{tt}\hat\epsilon_t}{1-h_{tt}}
  = \hat\epsilon_t\left(1 + \frac{h_{tt}}{1-h_{tt}}\right)
  = \frac{\hat\epsilon_t}{1-h_{tt}}
\]
Comme $0\leq h_{tt}<1$, le facteur $1/(1-h_{tt})$ est supérieur à $1$~: l'erreur
de prévision hors échantillon est \emph{toujours} plus grande que le résidu
ordinaire. \textbf{Pourquoi}~: le résidu ordinaire est artificiellement petit
parce que l'observation $t$ a participé à l'ajustement de la droite --- elle
s'est en partie «~expliquée elle-même~». Le retirer révèle l'erreur réellement
commise. C'est la version analytique de la validation croisée.

\question Comme $H$ est symétrique idempotente, $H = HH$ donne pour l'élément
diagonal $h_{tt} = \sum_s H_{ts}^2 = h_{tt}^2 + \sum_{s\neq t}H_{ts}^2$, d'où
$h_{tt}\geq h_{tt}^2$, c'est-à-dire $0\leq h_{tt}\leq1$. Par ailleurs
$\sum_t h_{tt} = \trace(H) = K$ (exercice~\ref{ex:hc}). Le levier moyen vaut
donc $K/T$~: c'est la référence à laquelle on compare les leviers individuels,
d'où la règle usuelle signalant les observations telles que $h_{tt}>2K/T$.

\question Le levier $h_{tt}$ mesure à quel point l'observation $t$ est atypique
dans l'espace des \emph{explicatives} (et non des $y$)~: il est d'autant plus
grand que $\mathbf x_t$ est éloigné du centre de gravité du nuage. La formule de
la question (3) montre que l'influence d'une observation sur l'estimation
combine deux ingrédients~: son levier $h_{tt}$ et son résidu $\hat\epsilon_t$.
Lorsque $h_{tt}\rightarrow1$, la droite passe exactement par le point,
$\hat\epsilon_t\rightarrow0$, et retirer l'observation modifie radicalement
l'estimation~: un seul point détermine alors une partie du résultat. C'est
précisément la situation du quatrième jeu de données d'Anscombe
(exercice~\ref{ex:r2}, question 7).

\question Sous homoscédasticité,
$\mathbb E[\hat\epsilon_t^{\,2}] = \epsvar(1-h_{tt})$~: les résidus ordinaires
sous-estiment la dispersion, d'autant plus que le levier est élevé. Deux
corrections viennent naturellement à l'esprit. Diviser par $1-h_{tt}$ corrige
exactement ce biais en espérance~: c'est HC2. Mais on peut aussi raisonner en
termes de \emph{prévision}~: la question (4) montre que l'erreur qu'on
commettrait réellement sur l'observation $t$ est $\hat\epsilon_t/(1-h_{tt})$, et
c'est son carré, $\hat\epsilon_t^{\,2}/(1-h_{tt})^2$, qu'il faut alors utiliser.
C'est HC3, et l'on comprend ainsi son nom d'approximation du \emph{jackknife}~:
elle repose sur les résidus de validation croisée. Elle sur-corrige légèrement
le biais, ce qui la rend prudente en petit échantillon --- exactement la
propriété recherchée (exercice~\ref{ex:hc}).

\section*{Chapitre II~-- Violations des conditions idéales}

\nouvelexercice\label{ex:erreurs-non-centrees}

\question En substituant le DGP,
$\hat{\mathbf b} = \beta + (X'X)^{-1}X'\varepsilon$, d'où~:
\[
  \mathbb E\left[\hat{\mathbf b}\right] = \beta + (X'X)^{-1}X'\xi
\]
Le biais vaut $(X'X)^{-1}X'\xi$. Il ne dépend que de la partie systématique des
erreurs.

\question En divisant par $T$ au numérateur et au dénominateur~:
\[
  \hat{\mathbf b} \plimT \beta + Q^{-1}\lim_{T\rightarrow\infty}\frac{X'\xi}{T}
\]
L'estimateur n'est donc pas convergent, sauf si $X'\xi/T\rightarrow0$.

\question Ici $\hat\epsilon = M\varepsilon$ mais $\varepsilon$ n'est plus centré.
En décomposant $\varepsilon = \xi + u$ avec $\mathbb E[u]=0$ et
$\mathbb V[u]=\epsvar I_T$~:
\[
  \mathbb E\left[SSE\right] = \mathbb E\left[\varepsilon'M\varepsilon\right]
  = \tracarg{M\,\mathbb V[\varepsilon]} + \xi'M\xi
  = \epsvar(T-K) + \xi'M\xi
\]
d'où le résultat annoncé. Comme $M$ est semi-définie positive, $\xi'M\xi\geq0$~:
$s^2$ \textbf{sur-estime} $\epsvar$. \textbf{Conséquence pour les tests.} Les
écarts-types estimés sont trop grands, les statistiques de Student trop petites~:
le test rejette trop rarement. On est donc trop conservateur, et la puissance
diminue --- l'inverse du danger rencontré au chapitre~IV.

\question Si $X$ contient une constante, alors $\iota$ est la première colonne
de $X$, donc $\iota = Xe_1$ où $e_1$ est le premier vecteur de la base
canonique. Il vient~:
\[
  (X'X)^{-1}X'\xi = \mu\,(X'X)^{-1}X'Xe_1 = \mu\,e_1
\]
Seul le \emph{premier} coefficient --- la constante --- est biaisé, et le biais
vaut exactement $\mu$. Toutes les pentes restent estimées sans biais.

\question Si $\xi = X_1\gamma$ où $X_1$ regroupe certaines colonnes de $X$, on
peut écrire $X_1 = XS$ avec $S$ une matrice de sélection, et~:
\[
  (X'X)^{-1}X'\xi = (X'X)^{-1}X'XS\gamma = S\gamma
\]
Seuls les coefficients associés aux colonnes de $X_1$ sont décalés, de $\gamma$~;
les autres restent sans biais. La question (4) en est le cas particulier
$X_1=\iota$.

\question Dans le cas de la question (4), $\xi = \mu\iota = \mu Xe_1$ appartient
à l'espace engendré par les colonnes de $X$, donc $M\xi = \mu MXe_1 = 0$ et~:
\[
  \mathbb E\left[s^2\right] = \epsvar
\]
L'estimateur de la variance redevient sans biais. Le même argument vaut pour la
question (5).

\question \textbf{Conclusion.} On inclut systématiquement une constante parce
qu'elle \emph{absorbe} toute espérance non nulle des erreurs~: le problème
disparaît alors des pentes, qui portent l'essentiel de l'interprétation
économique, et de $s^2$. Le prix à payer est un problème d'identification~: la
constante estimée converge vers $\beta_1+\mu$, et il est impossible de séparer
la véritable constante du modèle de la moyenne des erreurs. Seules les pentes
conservent une interprétation structurelle.

\nouvelexercice\label{ex:variable-omise} On note $P_1 = X_1(X_1'X_1)^{-1}X_1'$ et $M_1 = I-P_1$.

\question En substituant le DGP dans l'estimateur du modèle tronqué~:
\[
  \begin{split}
    \hat{\mathbf b}_1 &= (X_1'X_1)^{-1}X_1'\mathbf y\\
                      &= (X_1'X_1)^{-1}X_1'\left(X_1\beta_1 + X_2\beta_2 + \varepsilon\right)\\
                      &= \beta_1 + (X_1'X_1)^{-1}X_1'X_2\,\beta_2 + (X_1'X_1)^{-1}X_1'\varepsilon
  \end{split}
\]
Les explicatives étant déterministes et $\mathbb E[\varepsilon]=0$, le dernier
terme est d'espérance nulle, d'où le résultat. \textbf{Interprétation de
$(X_1'X_1)^{-1}X_1'X_2$}~: c'est la matrice des coefficients des régressions
auxiliaires de chaque colonne de $X_2$ sur l'ensemble des colonnes de $X_1$.
Le biais est donc le produit de l'effet de la variable omise sur $\mathbf y$
($\beta_2$) par le lien entre variable omise et variables incluses.

\question L'estimateur est sans biais si et seulement si $\beta_2 = 0$ (la
variable omise n'a en réalité aucun effet) \textbf{ou} $X_1'X_2 = 0$ (la variable
omise est orthogonale aux variables incluses). On retrouve l'exercice~\ref{ex:fwl}.

\question Avec $K_1=K_2=1$, en divisant numérateur et dénominateur par $T$ et en
appliquant la loi des grands nombres~:
\[
  \mathrm{plim}\ \hat b_1 = \beta_1 + \beta_2\,\frac{\sigma_{xz}}{\sigma_x^2}
\]
Le terme $\sigma_{xz}/\sigma_x^2$ est précisément le coefficient de la
régression de $z$ sur $x$. \textbf{Signe du biais.} Il est du signe du produit
$\beta_2\,\sigma_{xz}$~: positif si les deux termes sont de même signe, négatif
sinon. \textbf{Application.} Dans l'équation de salaire, le talent $z$ augmente
le salaire ($\beta_2>0$) et est vraisemblablement corrélé positivement aux
études ($\sigma_{xz}>0$)~: les MCO \emph{sur-estiment} le rendement de
l'éducation, en lui attribuant une part de l'effet du talent.

\question La somme des carrés des résidus du modèle tronqué vaut
$SSE = \mathbf y'M_1\mathbf y$. Comme $M_1X_1 = 0$~:
\[
  M_1\mathbf y = M_1X_2\beta_2 + M_1\varepsilon
\]
d'où, en prenant l'espérance de la forme quadratique~:
\[
  \mathbb E\left[SSE\right] = \beta_2'X_2'M_1X_2\beta_2 + \epsvar\left(T-K_1\right)
\]
et donc $\mathbb E[s^2] = \epsvar + \frac{\beta_2'X_2'M_1X_2\beta_2}{T-K_1}$.
Comme $M_1$ est semi-définie positive, le terme additionnel est positif~: $s^2$
\textbf{sur-estime} la variance des erreurs. C'est logique~: la variance de la
variable omise, non expliquée par $X_1$, est rejetée dans le résidu.

\question Le DGP s'écrit maintenant
$\mathbf y = X\beta+\varepsilon$ avec $X = (X_1\mid X_2)$ et
$\beta = (\beta_1',\ 0')'$. L'estimateur des MCO du modèle augmenté vérifie~:
\[
  \mathbb E\left[\hat{\mathbf b}\right] = \beta = \begin{pmatrix}\beta_1\\ 0\end{pmatrix}
\]
donc $\hat{\mathbf b}_1$ est sans biais (et $\hat{\mathbf b}_2$ estime sans
biais zéro). \textbf{Ajouter une variable inutile ne biaise rien.}

\question Par le théorème de Frisch-Waugh (exercice~\ref{ex:fwl}), dans le modèle augmenté~:
\[
  \mathbb V\left[\hat{\mathbf b}_1\right] = \epsvar\left(X_1'M_2X_1\right)^{-1},
  \qquad M_2 = I - X_2(X_2'X_2)^{-1}X_2'
\]
alors que dans le modèle correct $\mathbb V[\hat{\mathbf b}_1] = \epsvar(X_1'X_1)^{-1}$.
Or~:
\[
  X_1'X_1 - X_1'M_2X_1 = X_1'X_2(X_2'X_2)^{-1}X_2'X_1 \succeq 0
\]
donc $X_1'M_2X_1 \preceq X_1'X_1$, et en passant aux inverses
$\left(X_1'M_2X_1\right)^{-1} \succeq \left(X_1'X_1\right)^{-1}$.
\textbf{Conclusion}~: la variance de $\hat{\mathbf b}_1$ est (faiblement) plus
grande dans le modèle augmenté. L'égalité n'a lieu que si $X_1'X_2=0$.

\question \textbf{L'arbitrage.} Omettre une variable pertinente crée un
\emph{biais} qui ne disparaît pas asymptotiquement. L'inclure inutilement ne
crée aucun biais mais \emph{gonfle la variance}, d'autant plus que la variable
ajoutée est corrélée aux autres. En cas de doute, la logique du cours penche
pour l'inclusion~: un estimateur imprécis mais centré reste interprétable, un
estimateur biaisé ne l'est pas. Mais l'exercice~\ref{ex:multicolinearite} montrera que cette règle a
elle aussi une limite.

\nouvelexercice\label{ex:multicolinearite}

\question Les variables étant centrées et réduites,
$\frac1T\sum_t x_{k,t} = 0$, $\frac1T\sum_t x_{k,t}^2 = 1$ et
$\frac1T\sum_t x_{2,t}x_{3,t} = \rho$. La constante étant orthogonale aux
variables centrées, le bloc de $\frac1TX'X$ relatif aux deux pentes s'écrit~:
\[
  \frac{1}{T}X'X = \begin{pmatrix}1 & \rho\\ \rho & 1\end{pmatrix}
\]

\question En inversant cette matrice $2\times2$, dont le déterminant vaut
$1-\rho^2$~:
\[
  \mathbb V\begin{pmatrix}\hat b_2\\ \hat b_3\end{pmatrix}
  = \epsvar\left(X'X\right)^{-1}
  = \frac{\epsvar}{T\left(1-\rho^2\right)}\begin{pmatrix}1 & -\rho\\ -\rho & 1\end{pmatrix}
\]

\question Lorsque $\rho\rightarrow1$, le facteur $1/(1-\rho^2)$ diverge~: les
deux variances explosent. La covariance,
$-\rho\epsvar/\left(T(1-\rho^2)\right)$, tend vers $-\infty$~: les deux
estimateurs deviennent parfaitement négativement corrélés.
\textbf{Interprétation.} Les données ne permettent plus de distinguer l'effet
propre de $x_2$ de celui de $x_3$~: si l'on attribue davantage à l'un, il faut
retrancher à l'autre. En revanche la \emph{somme} $\hat b_2+\hat b_3$ reste
estimée précisément, sa variance valant
$2\epsvar/\left(T(1+\rho)\right)$, qui ne diverge pas.

\question La régression de $x_2$ sur $x_3$ (et la constante) a pour coefficient
de détermination $R_2^2 = \rho^2$, donc
$\mathrm{VIF}_2 = 1/(1-\rho^2)$. Or la question (2) donne~:
\[
  \mathbb V\left[\hat b_2\right] = \frac{\epsvar}{T}\times\frac{1}{1-\rho^2}
  = \underbrace{\frac{\epsvar}{T}}_{\text{variance sans colinéarité}}\times\ \mathrm{VIF}_2
\]
Le facteur d'inflation de la variance porte donc bien son nom~: il mesure de
combien la colinéarité multiplie la variance par rapport au cas orthogonal.

\question Si $\rho=1$, le déterminant s'annule~: $X'X$ n'est plus inversible.
C'est l'hypothèse $\mathcal H_2$ (plein rang de $X$) qui est violée.
L'estimateur des MCO n'est alors pas défini~: les équations normales admettent
une infinité de solutions, car seule la combinaison $b_2+b_3$ est identifiée.

\question \textbf{Non}~: la multicolinéarité imparfaite ($|\rho|<1$) n'est
\emph{pas} une violation des conditions idéales. L'hypothèse $\mathcal H_2$ est
satisfaite, et toutes les propriétés du chapitre~I tiennent~: $\hat{\mathbf b}$
est sans biais, convergent, et reste BLUE par Gauss-Markov. Le théorème garantit
même que sa variance est la plus petite possible parmi les estimateurs linéaires
sans biais --- elle est simplement grande dans l'absolu. Ce n'est donc pas un
défaut du modèle ni de l'estimateur, mais un manque d'information dans les
données.

\question \textbf{Symptômes}~: $R^2$ élevé et statistique de Fisher jointe très
significative alors qu'aucune statistique de Student n'est significative
(exercice~\ref{ex:tests}, question 7)~; coefficients de signe ou d'ampleur invraisemblables~;
forte instabilité des estimations lorsqu'on ajoute ou retire quelques
observations ou une variable~; $\mathrm{VIF}$ élevés et mauvais conditionnement
de $X'X$ (rapport des valeurs propres extrêmes). \textbf{Remèdes}~: augmenter la
taille de l'échantillon, agréger ou remplacer des variables redondantes,
mobiliser une information \emph{a priori} sous forme de contrainte linéaire
(exercice~\ref{ex:contraintes}) --- ce qui réduit la variance au prix d'un biais éventuel. En
revanche, supprimer purement et simplement une variable pertinente réintroduit
le biais de variable omise de l'exercice~\ref{ex:variable-omise}~: on ne fait alors que déplacer le
problème.

\nouvelexercice\label{ex:mcg}

\question On a toujours
$\hat{\mathbf b} = \beta + (X'X)^{-1}X'\varepsilon$, et comme $X$ est
déterministe et $\mathbb E[\varepsilon]=0$~:
\[
  \mathbb E\left[\hat{\mathbf b}\right] = \beta
\]
L'absence de biais ne fait intervenir que la nullité de l'espérance des erreurs,
pas leur variance. La convergence est acquise dès que $X'X/T\rightarrow Q$
définie positive et que $X'\Omega X/T$ admet une limite finie, de sorte que
$\mathbb V[\hat{\mathbf b}]\rightarrow0$.

\question Comme $X$ est déterministe~:
\[
  \mathbb V\left[\hat{\mathbf b}\right] = (X'X)^{-1}X'\,\mathbb V[\varepsilon]\,X(X'X)^{-1}
  = \epsvar (X'X)^{-1}X'\Omega X(X'X)^{-1}
\]
L'appellation «~sandwich~» vient de la structure de l'expression~: le facteur
$X'\Omega X$ (la «~garniture~») est encadré par deux fois le même facteur
$(X'X)^{-1}$ (le «~pain~»). Si $\Omega=I_T$, la garniture vaut $X'X$ et
l'expression se réduit à $\epsvar(X'X)^{-1}$.

\question Comme au chapitre~I, $SSE = \varepsilon'M\varepsilon$, donc~:
\[
  \mathbb E\left[SSE\right] = \tracarg{M\,\mathbb V[\varepsilon]} = \epsvar\tracarg{M\Omega}
\]
Or $\tracarg{M\Omega}\neq\trace(M)=T-K$ en général~: $s^2$ est \textbf{biaisé},
et le biais peut aller dans les deux sens selon la structure de $\Omega$.
\textbf{Conséquence.} La quantité $s^2(X'X)^{-1}$ imprimée par les logiciels
n'estime ni $\mathbb V[\hat{\mathbf b}]$ (qui est le sandwich de la question 2),
ni rien d'autre de pertinent. Les écarts-types, les tests de Student et de
Fisher sont donc invalides, y compris asymptotiquement.

\question Le modèle transformé a pour erreurs $\Lambda\varepsilon$, d'espérance
nulle et de variance~:
\[
  \mathbb V\left[\Lambda\varepsilon\right] = \epsvar\Lambda\Omega\Lambda'
\]
Or $\Omega^{-1}=\Lambda'\Lambda$ entraîne $\Omega = \Lambda^{-1}\left(\Lambda'\right)^{-1}$,
donc $\Lambda\Omega\Lambda' = \Lambda\Lambda^{-1}\left(\Lambda'\right)^{-1}\Lambda' = I_T$.
Les erreurs transformées sont donc sphériques, et la matrice $\Lambda X$ reste
déterministe et de plein rang~: le modèle transformé satisfait les conditions
idéales.

\question Il suffit d'appliquer les MCO au modèle transformé~:
\[
  \hat{\mathbf b}_{\textsc{mcg}}
  = \left((\Lambda X)'\Lambda X\right)^{-1}(\Lambda X)'\Lambda\mathbf y
  = \left(X'\Omega^{-1}X\right)^{-1}X'\Omega^{-1}\mathbf y
\]
En substituant le DGP,
$\hat{\mathbf b}_{\textsc{mcg}} = \beta + \left(X'\Omega^{-1}X\right)^{-1}X'\Omega^{-1}\varepsilon$,
d'espérance $\beta$~: l'estimateur est sans biais.

\question La variance se lit sur le modèle transformé~:
\[
  \mathbb V\left[\hat{\mathbf b}_{\textsc{mcg}}\right] = \epsvar\left(X'\Omega^{-1}X\right)^{-1}
\]
Le théorème de Gauss-Markov, appliqué au modèle transformé (qui satisfait les
conditions idéales), assure que $\hat{\mathbf b}_{\textsc{mcg}}$ est BLUE parmi
les estimateurs linéaires sans biais de $\beta$. En particulier, il est au moins
aussi précis que les MCO~:
\[
  \mathbb V\left[\hat{\mathbf b}\right] - \mathbb V\left[\hat{\mathbf b}_{\textsc{mcg}}\right] \succeq 0
\]
\textbf{On retrouve donc le théorème de Gauss-Markov}, mais appliqué au bon
modèle~: les MCO ne sont plus BLUE, les MCG le sont. Les MCO restent sans biais,
mais gaspillent de l'information en traitant toutes les observations de la même
façon alors qu'elles n'ont pas la même précision.

\question Dans le modèle transformé, l'estimateur sans biais de $\epsvar$ est la
somme des carrés des résidus transformés divisée par $T-K$~:
\[
  \tilde s^2 = \frac{\left(\mathbf y - X\hat{\mathbf b}_{\textsc{mcg}}\right)'\Omega^{-1}\left(\mathbf y-X\hat{\mathbf b}_{\textsc{mcg}}\right)}{T-K}
\]

\nouvelexercice\label{ex:ar1-cochrane}

\question Par stationnarité, $\mathbb V[\varepsilon_t]$ est constante et vérifie
$\mathbb V[\varepsilon_t] = \varphi^2\mathbb V[\varepsilon_t]+\sigma_{\nu}^2$, d'où~:
\[
  \epsvar \equiv \mathbb V\left[\varepsilon_t\right] = \frac{\sigma_{\nu}^2}{1-\varphi^2}
\]
En itérant $\varepsilon_t = \varphi^s\varepsilon_{t-s} + \sum_{j=0}^{s-1}\varphi^j\nu_{t-j}$
et en remarquant que les $\nu$ postérieurs à $t-s$ sont indépendants de
$\varepsilon_{t-s}$~:
\[
  \mathrm{cov}\left(\varepsilon_t,\varepsilon_{t-s}\right) = \varphi^s\,\epsvar
\]
En normalisant par $\epsvar$, on obtient une matrice de Toeplitz~:
\[
  \Omega = \begin{pmatrix}
    1 & \varphi & \varphi^2 & \cdots & \varphi^{T-1}\\
    \varphi & 1 & \varphi & \cdots & \varphi^{T-2}\\
    \varphi^2 & \varphi & 1 & \cdots & \varphi^{T-3}\\
    \vdots & & & \ddots & \vdots\\
    \varphi^{T-1} & \cdots & & \varphi & 1
  \end{pmatrix},
  \qquad \Omega_{ij} = \varphi^{|i-j|}
\]

\question On vérifie que $\Lambda'\Lambda$ reproduit bien $\Omega^{-1}$. La
matrice $\Omega^{-1}$ est tridiagonale~: ses éléments diagonaux valent
$1+\varphi^2$ (sauf aux extrémités, où ils valent $1$) et ses éléments
sous-diagonaux $-\varphi$, le tout divisé par $1-\varphi^2$. Le produit
$\Lambda'\Lambda$ donne exactement cette structure~: la première ligne de
$\Lambda$ contribue $1-\varphi^2$ en position $(1,1)$, et les lignes suivantes
$\left(-\varphi,\ 1\right)$ engendrent les termes $\varphi^2+1$ sur la diagonale
et $-\varphi$ hors diagonale.

\question Les observations transformées sont~:
\[
  \tilde y_1 = \sqrt{1-\varphi^2}\,y_1,
  \qquad
  \tilde y_t = y_t - \varphi\,y_{t-1} \quad (t\geq 2)
\]
et de même pour chaque colonne de $X$. \textbf{Le traitement particulier de la
première observation} vient de ce qu'elle n'a pas de passé~: on ne peut pas lui
retrancher $\varphi y_0$. Le facteur $\sqrt{1-\varphi^2}$ ramène sa variance au
niveau de celle des autres observations transformées ($\sigma_{\nu}^2$).
Si on la néglige, on perd une observation~: c'est la différence entre la
procédure de Cochrane-Orcutt (qui l'écarte) et celle de Prais-Winsten (qui la
conserve). Asymptotiquement cela ne change rien, mais en petit échantillon --- et
surtout lorsque $\varphi$ est proche de $1$, car la première observation reçoit
alors un poids relatif important --- la perte peut être sensible.

\question \textbf{Procédure des MCQG.} \emph{(i)} Estimer le modèle par les MCO
et récupérer les résidus $\hat\epsilon_t$. \emph{(ii)} Estimer $\varphi$ par la
régression de $\hat\epsilon_t$ sur $\hat\epsilon_{t-1}$, soit
$\hat\varphi = \sum_{t=2}^T\hat\epsilon_t\hat\epsilon_{t-1}/\sum_{t=2}^T\hat\epsilon_{t-1}^{\,2}$.
\emph{(iii)} Construire $\hat\Lambda$ en remplaçant $\varphi$ par $\hat\varphi$,
transformer les données et ré-estimer par les MCO. \emph{(iv)} Éventuellement
itérer jusqu'à convergence de $\hat\varphi$.

\question L'estimateur des MCQG est \textbf{convergent} dès que~:
\[
  \frac{X'\hat\Omega^{-1}X}{T}\plimT Q
  \qquad\text{et}\qquad
  \frac{X'\hat\Omega^{-1}\varepsilon}{T}\plimT 0
\]
Il est \textbf{asymptotiquement équivalent} aux MCG si, de plus~:
\[
  \frac{X'\left(\hat\Omega^{-1}-\Omega^{-1}\right)X}{T}\plimT 0
  \qquad\text{et}\qquad
  \frac{X'\left(\hat\Omega^{-1}-\Omega^{-1}\right)\varepsilon}{\sqrt T}\plimT 0
\]
Ces conditions signifient que l'erreur d'estimation de $\Omega$ devient
négligeable assez vite pour ne pas polluer la loi limite.

\question \textbf{Non.} Ce sont des propriétés strictement asymptotiques. À
distance finie, $\hat\Omega$ est aléatoire et corrélé aux données, de sorte que
l'estimateur des MCQG est en général biaisé, que sa variance exacte est inconnue,
et qu'il n'est pas BLUE. Il n'y a aucune garantie qu'il batte les MCO en petit
échantillon, surtout si $\hat\varphi$ est mal estimé.

\question Avec $\varphi>0$ et un régresseur lui-même positivement autocorrélé,
la vraie variance $\epsvar(X'X)^{-1}X'\Omega X(X'X)^{-1}$ est \emph{supérieure} à
$\epsvar(X'X)^{-1}$, car $X'\Omega X > X'X$~: les observations successives
apportent moins d'information indépendante qu'il n'y paraît. L'économètre qui
utilise $s^2(X'X)^{-1}$ \textbf{sous-estime} donc les écarts-types. Les
statistiques de Student sont trop grandes, le test rejette trop souvent~: on
conclut à tort à la significativité. C'est le mécanisme des régressions
fallacieuses.

\nouvelexercice\label{ex:tcl} \textit{Exercice de simulation~; on donne les résultats attendus
et leur lecture.}

\question Les trois lois se construisent en centrant et en réduisant~: pour
l'uniforme, $\varepsilon_t = \sqrt{12}\left(u_t-\tfrac12\right)$ avec
$u_t\sim\mathcal U_{[0,1]}$~; pour l'exponentielle de paramètre $1$,
$\varepsilon_t = e_t-1$~; pour la Student à $5$ degrés de liberté, dont la
variance vaut $5/3$, $\varepsilon_t = t_t/\sqrt{5/3}$.

\question La moyenne empirique des $\hat b$ reste centrée sur $\beta = 2$ quelle
que soit la loi des erreurs. C'est attendu~: l'absence de biais ne repose que
sur $\mathbb E[\varepsilon]=0$ et le caractère déterministe de $X$, pas sur la
normalité. \textbf{Le théorème de Gauss-Markov reste également valide}~: il ne
suppose que des erreurs centrées, homoscédastiques et non corrélées. Les MCO
demeurent donc BLUE. C'est la \emph{loi exacte} de l'estimateur, et elle seule,
qui est perdue.

\question Pour $T=10$, la distribution de $\hat b-\beta$ hérite de la forme de
celle des erreurs~: nettement asymétrique avec des erreurs exponentielles,
à queues épaisses avec des erreurs de Student. Elle n'est pas gaussienne. Les
lois exactes de Student et de Fisher du chapitre~I, qui reposent sur
$\mathcal H_3$, ne s'appliquent donc pas.

\question Les histogrammes de $\sqrt T\left(\hat b-\beta\right)$ se rapprochent
d'une gaussienne à mesure que $T$ croît, et lui sont pratiquement confondus dès
$T=100$ pour l'uniforme, un peu plus tard pour l'exponentielle et la Student
(dont les queues sont plus lourdes). C'est le théorème central limite~:
l'estimateur est une moyenne pondérée des erreurs, et une somme convenablement
normalisée de variables i.i.d.\ de variance finie converge en loi vers une
gaussienne, \emph{quelle que soit} la loi des termes sommés.

\question La taille effective du test s'écarte nettement de $5\,\%$ pour $T=10$
(sur-rejet marqué avec des erreurs à queues épaisses), puis converge vers
$5\,\%$ quand $T$ croît. \textbf{Lecture.} L'hypothèse de normalité n'est pas
nécessaire pour que les tests soient \emph{asymptotiquement} valides~; elle l'est
pour qu'ils soient valides \emph{exactement}, à toute taille d'échantillon.

\question Avec des erreurs de Cauchy, rien ne se stabilise~: la distribution de
$\sqrt T(\hat b-\beta)$ ne converge pas, et la variance empirique des $\hat b$
explose avec le nombre de réplications. La raison est que la loi de Cauchy n'a
\textbf{ni espérance ni variance}~: l'hypothèse de moment d'ordre deux fini,
indispensable au théorème central limite de Lindeberg-Lévy, n'est pas
satisfaite. On retrouve le phénomène signalé au chapitre~III à propos des
instruments (exercice~\ref{ex:instruments-faibles}).

\question \textbf{Conclusion.} La normalité des erreurs est nécessaire pour les
résultats \emph{à distance finie}~: loi exacte de $\hat{\mathbf b}$, lois de
Student et de Fisher, équivalence avec le maximum de vraisemblance, efficacité
au sens de Cramer-Darmois-Fréchet-Rao. Elle n'est nécessaire ni pour l'absence
de biais, ni pour la convergence, ni pour Gauss-Markov, ni pour la validité
asymptotique des tests. Ce qui est réellement indispensable, en revanche, c'est
l'existence d'un moment d'ordre deux fini.

\nouvelexercice\label{ex:donnees-groupees}

\question En moyennant le modèle individuel à l'intérieur du groupe $g$, et
comme $\mathbf x_g$ ne varie pas dans le groupe~:
\[
  \bar y_g = \mathbf x_g'\beta + \bar\varepsilon_g,
  \qquad \bar\varepsilon_g = \frac{1}{n_g}\sum_{i=1}^{n_g}\varepsilon_{ig}
\]
Les $\varepsilon_{ig}$ étant non corrélés et de variance $\epsvar$~:
\[
  \mathbb V\left[\bar\varepsilon_g\right] = \frac{1}{n_g^2}\sum_{i=1}^{n_g}\mathbb V\left[\varepsilon_{ig}\right] = \frac{\epsvar}{n_g}
\]

\question Les erreurs agrégées restent non corrélées entre groupes, donc~:
\[
  \Omega = \mathrm{diag}\left(\frac{1}{n_1},\dots,\frac{1}{n_G}\right)
\]
\textbf{Ce qui rend la situation exceptionnelle}, c'est que $\Omega$ est
\emph{connue} --- entièrement, à partir des seuls effectifs observés, sans avoir
à estimer quoi que ce soit. Or c'est précisément l'hypothèse que faisait le
chapitre~II pour définir les MCG, et qu'on jugeait irréaliste~: ici elle est
littéralement satisfaite. On peut donc appliquer les MCG, et non les MCQG~: tous
les résultats à distance finie (absence de biais, BLUE, lois exactes sous
normalité) sont disponibles.

\question En reportant $\Omega^{-1} = \mathrm{diag}(n_1,\dots,n_G)$~:
\[
  \hat{\mathbf b}_{\textsc{mcg}} = \left(\sum_{g=1}^G n_g\,\mathbf x_g\mathbf x_g'\right)^{-1}\sum_{g=1}^G n_g\,\mathbf x_g\bar y_g
\]
Chaque groupe est pondéré par son effectif. De façon équivalente, on applique
les MCO aux données transformées $\sqrt{n_g}\,\bar y_g$ et
$\sqrt{n_g}\,\mathbf x_g$, dont les erreurs
$\sqrt{n_g}\,\bar\varepsilon_g$ ont toutes pour variance $\epsvar$.

\question Sur les données individuelles, les MCO donnent~:
\[
  \hat{\mathbf b} = \left(\sum_{g}\sum_{i}\mathbf x_g\mathbf x_g'\right)^{-1}\sum_g\sum_i \mathbf x_g y_{ig}
\]
Comme $\mathbf x_g$ ne dépend pas de $i$, la première somme vaut
$\sum_g n_g\mathbf x_g\mathbf x_g'$ et la seconde
$\sum_g \mathbf x_g\sum_i y_{ig} = \sum_g n_g\mathbf x_g\bar y_g$. Les deux
expressions sont \textbf{identiques}~:
\[
  \hat{\mathbf b}_{\textsc{mcg}}\ \text{(sur moyennes)} = \hat{\mathbf b}\ \text{(sur individus)}
\]
\textbf{Interprétation.} Correctement pondérée, l'agrégation ne détruit aucune
information sur $\beta$ --- ce qui est rassurant, et fournit une justification
très concrète de la pondération~: pondérer par $n_g$, c'est simplement
reconstituer le fait qu'un groupe de $1000$ individus contient mille fois plus
d'information qu'un groupe d'un seul.

\question L'estimateur reste \textbf{sans biais}~: l'absence de biais ne dépend
que de $\mathbb E[\varepsilon]=0$, pas de la pondération. Ce que l'on perd,
c'est l'\emph{efficacité}~: les MCO non pondérés traitent à égalité une moyenne
calculée sur trois individus et une moyenne calculée sur trois mille, alors que
la première est bien plus bruitée. Par le théorème de Gauss-Markov appliqué au
modèle transformé (exercice~\ref{ex:mcg}), les MCG font strictement mieux dès
que les $n_g$ ne sont pas tous égaux. Les écarts-types usuels sont de surcroît
faux, puisque calculés sous une hypothèse d'homoscédasticité fausse.

\question Avec $\mathbf x_g=1$, l'estimateur des MCG est la moyenne pondérée,
c'est-à-dire la moyenne de l'ensemble des individus~:
\[
  \mathbb V\left[\hat b_{\textsc{mcg}}\right] = \frac{\epsvar}{n_1+n_2} = \frac{\epsvar}{110} \simeq 0{,}0091\,\epsvar
\]
Les MCO non pondérés donnent la moyenne des deux moyennes de groupe~:
\[
  \mathbb V\left[\hat b\right] = \frac{1}{4}\left(\frac{\epsvar}{100}+\frac{\epsvar}{10}\right) = 0{,}0275\,\epsvar
\]
Le rapport vaut $\mathbf{3{,}03}$~: ignorer la pondération multiplie la variance
par trois, soit un écart-type $1{,}7$ fois plus grand. \textbf{La raison est
frappante}~: en donnant le même poids aux deux groupes, on laisse le petit
groupe de $10$ individus dicter la moitié du résultat.

\question Sans les $n_g$, on ne peut pas construire $\Omega$ et la pondération
est hors d'atteinte. Deux voies restent ouvertes~: estimer la structure de
variance à partir des résidus (MCQG, avec les réserves de
l'exercice~\ref{ex:ar1-cochrane}), ou --- plus prudent --- conserver les MCO non
pondérés et se contenter de corriger les écarts-types par la méthode robuste de
l'exercice~\ref{ex:hc}, qui ne demande à connaître ni les $n_g$ ni la forme de
$\Omega$. On accepte alors une perte d'efficacité en échange d'une inférence
valide.

\nouvelexercice\label{ex:mauvais-controle}

\question En substituant la première équation dans la seconde~:
\[
  y_i = \beta x_i + \delta\left(\gamma x_i + u_i\right) + \varepsilon_i
      = \underbrace{\left(\beta+\delta\gamma\right)}_{\text{effet total}}x_i + \left(\delta u_i+\varepsilon_i\right)
\]
L'\textbf{effet total} de $x$ sur $y$ vaut $\beta+\delta\gamma$~: il additionne
l'effet direct $\beta$ et l'effet indirect $\delta\gamma$ transitant par $m$.

\question La perturbation de la forme réduite, $\delta u_i+\varepsilon_i$, est
non corrélée à $x_i$ puisque $u$ et $\varepsilon$ le sont. La régression courte
est donc convergente~:
\[
  \mathrm{plim}\ \hat b_x = \beta+\delta\gamma
\]
Elle estime l'\textbf{effet total}.

\question Dans la régression longue, appliquons le théorème de Frisch-Waugh
(exercice~\ref{ex:fwl})~: le résidu de la régression de $m$ sur $x$ est
exactement $u$. Donc~:
\[
  \mathrm{plim}\ \hat b_m = \frac{\mathrm{cov}(u,y)}{\mathbb V[u]}
  = \frac{\mathrm{cov}\left(u,\ \beta x+\delta(\gamma x+u)+\varepsilon\right)}{\sigma_u^2}
  = \frac{\delta\sigma_u^2}{\sigma_u^2} = \delta
\]
et, par la formule de la variable omise reliant régression courte et longue
(exercice~\ref{ex:variable-omise}),
$\mathrm{plim}\,\hat b_x^{\text{court}} = \mathrm{plim}\,\hat b_x + \gamma\,\mathrm{plim}\,\hat b_m$,
d'où~:
\[
  \mathrm{plim}\ \hat b_x = \left(\beta+\delta\gamma\right) - \gamma\delta = \beta
\]
La régression longue estime l'\textbf{effet direct}.

\question Les deux estimateurs sont convergents, mais \emph{pour des paramètres
différents}~: il n'y a pas ici de «~bonne~» et de «~mauvaise~» régression au
sens statistique. Tout dépend de la question posée. Si l'on demande de combien
le salaire augmente lorsqu'on étudie un an de plus --- \emph{toutes conséquences
comprises} --- la réponse est l'\textbf{effet total} $\beta+\delta\gamma$, donc
la régression \textbf{courte}. La régression longue répond à une autre question~:
de combien le salaire augmente-t-il avec une année d'études supplémentaire,
\emph{à profession inchangée}.

\question Si $\mathrm{cov}(u,\varepsilon)=\sigma_{u\varepsilon}\neq0$, le calcul
de la question (3) devient~:
\[
  \mathrm{plim}\ \hat b_m = \frac{\mathrm{cov}(u,y)}{\sigma_u^2}
  = \frac{\delta\sigma_u^2+\sigma_{u\varepsilon}}{\sigma_u^2}
  = \delta + \frac{\sigma_{u\varepsilon}}{\sigma_u^2}
\]
tandis que la régression courte reste inchangée ($x$ est toujours indépendant de
$u$ et $\varepsilon$), donc égale à $\beta+\delta\gamma$. En reportant dans la
même relation~:
\[
  \mathrm{plim}\ \hat b_x = \left(\beta+\delta\gamma\right) - \gamma\left(\delta+\frac{\sigma_{u\varepsilon}}{\sigma_u^2}\right)
  = \beta - \gamma\,\frac{\sigma_{u\varepsilon}}{\sigma_u^2}
\]
\textbf{Ce résultat mérite d'être médité.} La variable $x$ est parfaitement
exogène~: la régression courte l'estimait sans biais. C'est \emph{l'ajout du
contrôle} $m$ --- variable elle-même endogène --- qui a introduit un biais sur le
coefficient de $x$. Ajouter une variable n'est donc pas une opération neutre.

\question Avec $x$ = années d'études, $m$ = profession et $y$ = salaire~: la
profession est un \emph{canal} par lequel l'éducation agit sur le salaire
(étudier permet d'accéder à des emplois mieux rémunérés). Contrôler par la
profession revient à comparer des individus \emph{exerçant le même métier}, donc
à neutraliser précisément ce canal. On mesure alors uniquement le supplément de
salaire obtenu à métier donné, ce qui sous-estime gravement le rendement de
l'éducation si l'on cherchait son effet total. Pire~: la profession dépend aussi
de déterminants inobservés (ambition, réseau) eux-mêmes liés au salaire, ce qui
place dans le cas de la question (5).

\question \textbf{La confrontation est instructive.} L'exercice~\ref{ex:variable-omise}
concluait qu'en cas de doute il vaut mieux inclure une variable, puisqu'omettre
biaise tandis qu'ajouter ne fait que gonfler la variance. Cette règle n'est
valable que pour une variable de contrôle \emph{antérieure} au traitement --- une
cause commune de $x$ et de $y$. Elle tombe pour une variable située \emph{en
aval}, sur le chemin causal~: on parle alors de «~mauvais contrôle~». La leçon
est qu'aucun critère statistique --- ni le $R^2$, ni la significativité du
coefficient de $m$, qui est ici parfaitement significatif dans les deux
scénarios --- ne permet de trancher. Le choix des variables de contrôle relève
d'un raisonnement \emph{causal} extérieur aux données~: il faut savoir ce qui
détermine quoi, et dans quel ordre. Une régression bien ajustée peut répondre
avec précision à une question qu'on ne se posait pas.

\nouvelexercice\label{ex:moulton}

\question Les deux composantes étant indépendantes~:
\[
  \mathbb V\left[\varepsilon_{ig}\right] = \sigma_v^2+\sigma_w^2 = \epsvar
\]
Pour deux individus distincts de la \emph{même} région, seul le choc régional
est commun~:
\[
  \mathrm{cov}\left(\varepsilon_{ig},\varepsilon_{jg}\right) = \mathbb V[v_g] = \sigma_v^2
  \qquad (i\neq j)
\]
et la covariance est nulle entre régions différentes. Le coefficient de
corrélation intra-groupe vaut donc~:
\[
  \rho = \frac{\sigma_v^2}{\sigma_v^2+\sigma_w^2}
\]

\question $\Omega$ est \textbf{bloc-diagonale}, chaque bloc correspondant à une
région~:
\[
  \Omega_g = \sigma_w^2 I_n + \sigma_v^2\,\iota_n\iota_n'
  = \epsvar\left[(1-\rho)I_n + \rho\,\iota_n\iota_n'\right]
\]
Elle n'est \textbf{pas diagonale}~: les termes hors diagonale valent $\sigma_v^2$
à l'intérieur de chaque bloc. C'est la différence décisive avec tout le
chapitre~IV, et elle va tout changer.

\question Comme $x$ ne varie qu'au niveau régional~:
\[
  \hat\beta - \beta = \frac{\sum_g\sum_i \left(x_g-\bar x\right)\varepsilon_{ig}}{\sum_g\sum_i\left(x_g-\bar x\right)^2}
  = \frac{\sum_g \left(x_g-\bar x\right)\sum_i\varepsilon_{ig}}{n\sum_g\left(x_g-\bar x\right)^2}
\]
Or $\sum_i\varepsilon_{ig} = nv_g + \sum_i w_{ig}$ a pour variance
$n^2\sigma_v^2 + n\sigma_w^2$, et les régions sont indépendantes. D'où~:
\[
  \mathbb V\left[\hat\beta\right]
  = \frac{\left(n^2\sigma_v^2+n\sigma_w^2\right)\sum_g\left(x_g-\bar x\right)^2}{n^2\left(\sum_g\left(x_g-\bar x\right)^2\right)^2}
  = \frac{\sigma_w^2+n\sigma_v^2}{n\sum_g\left(x_g-\bar x\right)^2}
\]

\question La formule usuelle, qui suppose $\Omega=\epsvar I_N$, donne~:
\[
  \frac{\epsvar}{\sum_{i,g}\left(x_g-\bar x\right)^2} = \frac{\epsvar}{n\sum_g\left(x_g-\bar x\right)^2}
\]
Le rapport de la vraie variance à celle-ci vaut donc~:
\[
  \frac{\sigma_w^2+n\sigma_v^2}{\epsvar}
  = \frac{\left(\epsvar-\sigma_v^2\right)+n\sigma_v^2}{\epsvar}
  = 1+(n-1)\frac{\sigma_v^2}{\epsvar}
  = \boxed{\,1+(n-1)\rho\,}
\]
C'est le \emph{facteur de Moulton}. Il ne dépend que de la taille des groupes et
de la corrélation intra-groupe.

\question Avec $n=30$ et $\rho=0{,}1$~:
\[
  1+(n-1)\rho = 1+29\times0{,}1 = 3{,}9
\]
La vraie variance est donc $3{,}9$ fois la variance annoncée~: les écarts-types
usuels sont sous-estimés d'un facteur $\sqrt{3{,}9}\simeq 2{,}0$, et les
statistiques de Student sont \textbf{deux fois trop grandes}. Un coefficient
affiché à $t=2{,}5$, apparemment significatif, a en réalité un $t$ voisin de
$1{,}3$~: il ne l'est pas. \textbf{Ce qui frappe, c'est l'ampleur du problème
pour une corrélation intra-groupe modeste}~: $\rho = 0{,}1$ paraît négligeable,
mais il est multiplié par $n-1$. Avec des groupes de $100$ individus, le facteur
passerait à $10{,}9$ et les écarts-types seraient sous-estimés d'un facteur
$3{,}3$. Augmenter le nombre d'individus \emph{par région} n'aide pas~: seul le
nombre de régions apporte de l'information sur $\beta$.

\question \textbf{Non, et c'est le point important.} L'estimateur sandwich de
l'exercice~\ref{ex:hc} estime la garniture par
$\sum_t\hat\varepsilon_t^{\,2}\mathbf x_t\mathbf x_t'$~: cette somme ne contient
que des termes «~diagonaux~», indexés par une seule observation. Elle est conçue
pour corriger l'inégalité des \emph{variances}, sous l'hypothèse que $\Omega$
est diagonale. Or ici le problème ne vient pas des variances --- qui sont toutes
égales à $\epsvar$, le modèle est parfaitement homoscédastique ! --- mais des
\emph{covariances} hors diagonale, que la formule ignore purement et simplement.
Appliquer HC0 ou HC3 ne corrigerait donc rien du tout.

\question Deux solutions. \emph{(i)} \textbf{Agréger}~: estimer le modèle sur
les $G$ moyennes régionales. On revient alors à
l'exercice~\ref{ex:donnees-groupees}, avec des erreurs
$v_g+\bar w_g$ indépendantes entre régions~: le problème disparaît, au prix
d'une perte apparente d'observations --- apparente seulement, puisque
l'information utile était bien de niveau régional. \emph{(ii)} Utiliser des
\textbf{écarts-types en grappes}, qui remplacent la garniture par
$\sum_g \left(X_g'\hat\varepsilon_g\right)\left(X_g'\hat\varepsilon_g\right)'$~:
on agrège les résidus \emph{à l'intérieur} de chaque grappe avant de les élever
au carré, ce qui capture précisément les covariances intra-groupe. Cette
correction est hors du programme du cours, mais l'exercice montre pourquoi elle
est devenue systématique en économie appliquée. Elle exige un nombre de grappes
suffisant~: c'est $G$, et non $N$, qui gouverne la qualité de l'approximation
asymptotique.

\section*{Chapitre III~-- Régresseurs aléatoires et variables instrumentales}

\nouvelexercice\label{ex:lie}

\question L'estimateur des MCO de la pente dans un modèle simple s'écrit~:
\[
  \hat b = \frac{\sum_{t=1}^T\left(y_t-\bar y\right)\left(x_t-\bar x\right)}{\sum_{t=1}^T\left(x_t-\bar x\right)^2}
\]

\question En substituant $y_t-\bar y = \beta\left(x_t-\bar x\right)+\left(\varepsilon_t-\bar\varepsilon\right)$~:
\[
  \hat b = \beta + \frac{\sum_t\left(\varepsilon_t-\bar\varepsilon\right)\left(x_t-\bar x\right)}{\sum_t\left(x_t-\bar x\right)^2}
\]
Contrairement au chapitre~I, le second terme est le rapport de deux variables
aléatoires~: le dénominateur n'est plus déterministe.

\question La densité conditionnelle de $V$ sachant $U$ est
$f_{V|U}(v|u) = f_{U,V}(u,v)/f_U(u)$, et l'espérance conditionnelle est la
fonction~:
\[
  \Psi(u) = \mathbb E\left[V|U=u\right] = \int v\,f_{V|U}(v|u)\,\mathrm dv
\]
C'est bien une fonction de $u$~; $\Psi(U)$ est donc une variable aléatoire.

\question \textit{Démonstration (loi des espérances itérées).}
\[
  \begin{split}
    \mathbb E\left[\Psi(U)\right]
      &= \int \Psi(u) f_U(u)\,\mathrm du
       = \int\left(\int v\,f_{V|U}(v|u)\,\mathrm dv\right)f_U(u)\,\mathrm du\\
      &= \int\int v\,\underbrace{f_{V|U}(v|u)f_U(u)}_{=f_{U,V}(u,v)}\,\mathrm dv\,\mathrm du
       = \int v\left(\int f_{U,V}(u,v)\,\mathrm du\right)\mathrm dv\\
      &= \int v\,f_V(v)\,\mathrm dv = \mathbb E[V]
  \end{split}
\]
On note $\mathbb E\left[\mathbb E[V|U]\right] = \mathbb E[V]$. C'est l'outil
central du chapitre~: il permet de passer d'un énoncé conditionnel (où $x$ est
traité comme fixe, donc où l'on retrouve le chapitre~I) à un énoncé marginal.

\question En appliquant la loi des espérances itérées avec $U=x_s$ et
$V=\varepsilon_t$~:
\[
  \mathbb E\left[\varepsilon_t\right] = \mathbb E\left[\mathbb E\left[\varepsilon_t|x_s\right]\right] = \mathbb E[0] = 0
\]

\question De même~: une fois que l'on s'est donné $x_s$, la quantité $x_s$
est connue et se comporte donc comme une constante, que l'on peut sortir de
l'espérance conditionnelle~:
\[
  \mathbb E\left[\varepsilon_t x_s\right]
  = \mathbb E\left[\mathbb E\left[\varepsilon_t x_s|x_s\right]\right]
  = \mathbb E\left[x_s\underbrace{\mathbb E\left[\varepsilon_t|x_s\right]}_{=0}\right] = 0
\]
pour tout $t,s$. Comme $\mathbb E[\varepsilon_t]=0$, cela signifie aussi que
$\mathrm{cov}(\varepsilon_t,x_s)=0$~: erreurs et régresseurs sont non corrélés,
à toutes les dates.

\question Puisque $\mathbb E[\varepsilon_t]=0$,
$\mathbb V[\varepsilon_t] = \mathbb E[\varepsilon_t^2]$, et~:
\[
  \mathbb V\left[\varepsilon_t\right]
  = \mathbb E\left[\mathbb E\left[\varepsilon_t^2|x_s\right]\right]
  = \mathbb E\left[\epsvar\right] = \epsvar
\]
la variance conditionnelle étant constante par hypothèse. La variance marginale
ne dépend donc pas de $t$~: les erreurs sont homoscédastiques, non seulement
conditionnellement mais aussi inconditionnellement.

\question En conditionnant par rapport à l'ensemble de l'échantillon
$x = (x_1,\dots,x_T)$, le dénominateur devient déterministe et~:
\[
  \mathbb E\left[\hat b\,\middle|\,x\right]
  = \beta + \frac{\sum_t\left(x_t-\bar x\right)\mathbb E\left[\varepsilon_t-\bar\varepsilon\,\middle|\,x\right]}{\sum_t\left(x_t-\bar x\right)^2}
  = \beta
\]
puis, par la loi des espérances itérées,
$\mathbb E[\hat b] = \mathbb E\left[\mathbb E[\hat b|x]\right] = \beta$.
\textbf{L'estimateur des MCO est donc sans biais}, bien que le régresseur soit
aléatoire.

\smallskip

\textbf{Remarque importante.} Le passage à l'espérance conditionnelle par
rapport au \emph{vecteur} $x$ tout entier suppose
$\mathbb E[\varepsilon_t|x_1,\dots,x_T]=0$, hypothèse dite d'\emph{exogénéité
stricte}, légèrement plus forte que la condition $\mathbb E[\varepsilon_t|x_s]=0$
prise deux à deux dans l'énoncé. C'est bien ainsi qu'il faut lire l'exercice.
Cette distinction n'est pas anodine~: l'exercice~\ref{ex:ar1-biais} exhibe précisément un cas où
l'orthogonalité contemporaine est satisfaite mais l'exogénéité stricte ne l'est
pas, et où l'estimateur est alors biaisé tout en restant convergent.

\nouvelexercice\label{ex:vi-simple}

\question \textbf{Non.} Le raisonnement de l'exercice précédent tombe~: comme
$\mathrm{cov}(x,\varepsilon)\neq0$, on n'a plus
$\mathbb E[\varepsilon_t|x]=0$, et l'espérance du second terme de
$\hat b = \beta + \sum_t(\varepsilon_t-\bar\varepsilon)(x_t-\bar x)/\sum_t(x_t-\bar x)^2$
n'est pas nulle. L'estimateur est biaisé.

\question En divisant numérateur et dénominateur par $T$ et en appliquant la loi
des grands nombres puis le théorème de Slutsky~:
\[
  \hat b \plimT \beta + \frac{\mathrm{cov}(x,\varepsilon)}{\mathbb V[x]} \neq \beta
\]
L'estimateur des MCO n'est \textbf{pas convergent}~: le problème ne disparaît pas
avec la taille de l'échantillon. C'est bien plus grave qu'un simple biais.

\question Les deux équations sont les contreparties empiriques des conditions de
moments que satisfait le vrai modèle~: $\mathbb E[\varepsilon_t]=0$ (les erreurs
sont centrées) et $\mathbb E[z_t\varepsilon_t]=0$ ($z$ n'est pas corrélé aux
erreurs). On remplace les espérances par des moyennes empiriques et les
paramètres inconnus par leurs estimateurs~: c'est la méthode des moments. On
n'impose \emph{pas} $\frac1T\sum_t x_t\hat\epsilon_t = 0$, précisément parce que
cette condition est fausse dans la population.

\question La première équation donne
$\hat a_{\textsc{vi}} = \bar y - \hat b_{\textsc{vi}}\bar x$. En reportant dans
la seconde~:
\[
  \frac{1}{T}\sum_t z_t\left[\left(y_t-\bar y\right) - \hat b_{\textsc{vi}}\left(x_t-\bar x\right)\right] = 0
\]
Comme $\sum_t\bar z\left[(y_t-\bar y)-\hat b_{\textsc{vi}}(x_t-\bar x)\right]=0$,
on peut centrer $z_t$~:
\[
  \hat b_{\textsc{vi}} = \frac{\sum_t\left(z_t-\bar z\right)\left(y_t-\bar y\right)}{\sum_t\left(z_t-\bar z\right)\left(x_t-\bar x\right)}
\]
L'estimateur existe dès que $\mathrm{cov}(z,x)\neq0$~: c'est la condition de
\emph{pertinence}.

\question En substituant le DGP au numérateur~:
\[
  \hat b_{\textsc{vi}} = \beta + \frac{\frac1T\sum_t\left(z_t-\bar z\right)\left(\varepsilon_t-\bar\varepsilon\right)}{\frac1T\sum_t\left(z_t-\bar z\right)\left(x_t-\bar x\right)}
  \plimT \beta + \frac{\mathrm{cov}(z,\varepsilon)}{\mathrm{cov}(z,x)} = \beta
\]
puisque $\mathrm{cov}(z,\varepsilon)=0$ (\emph{exogénéité}) et
$\mathrm{cov}(z,x)\neq0$ (\emph{pertinence}). L'estimateur à variable
instrumentale est donc convergent, là où les MCO ne l'étaient pas.

\question Le théorème central limite appliqué au numérateur donne~:
\[
  \mathbb V_{\infty}\left[\sqrt T\left(\hat b_{\textsc{vi}}-\beta\right)\right]
  = \frac{\epsvar\,\mathbb V[z]}{\mathrm{cov}(z,x)^2}
\]
\textbf{Comparaison avec les MCO.} Si $x$ était exogène, la variance asymptotique
des MCO serait $\epsvar/\mathbb V[x]$. Le rapport des deux vaut~:
\[
  \frac{\mathbb V_{\infty}\left[\sqrt T\left(\hat b_{\textsc{vi}}-\beta\right)\right]}{\mathbb V_{\infty}\left[\sqrt T\left(\hat b-\beta\right)\right]}
  = \frac{\mathbb V[z]\,\mathbb V[x]}{\mathrm{cov}(z,x)^2}
  = \frac{1}{\rho_{xz}^2} > 1
\]
puisque $\rho_{xz}\in\left]-1,1\right[$. \textbf{Interprétation.} L'estimateur à
variables instrumentales est \emph{toujours} moins précis que les MCO, et
d'autant plus que l'instrument est faiblement corrélé au régresseur. C'est le
prix de la convergence~: on n'exploite que la part de la variation de $x$
qu'explique $z$. On y reviendra à l'exercice~\ref{ex:instruments-faibles}.

\nouvelexercice\label{ex:vi-matriciel}

\question On impose les $K$ conditions de moments empiriques
$Z'\left(\mathbf y - X\hat{\mathbf b}_{\textsc{vi}}\right) = 0$, soit
$Z'X\hat{\mathbf b}_{\textsc{vi}} = Z'\mathbf y$. Comme $Z'X$ est carrée
($K\times K$) et supposée inversible~:
\[
  \hat{\mathbf b}_{\textsc{vi}} = \left(Z'X\right)^{-1}Z'\mathbf y
\]
On retrouve les MCO en prenant $Z=X$.

\question Les deux conditions du théorème principal sont~: \emph{(i)}
\textbf{pertinence}, $\frac{Z'X}{T}\plimT Q_{ZX}$ finie et non singulière~;
\emph{(ii)} \textbf{exogénéité}, $\frac{Z'\varepsilon}{T}\plimT 0$. En
substituant le DGP~:
\[
  \hat{\mathbf b}_{\textsc{vi}} = \beta + \left(\frac{Z'X}{T}\right)^{-1}\frac{Z'\varepsilon}{T}
  \plimT \beta + Q_{ZX}^{-1}\cdot 0 = \beta
\]
en utilisant le théorème de Slutsky (l'inversion matricielle étant continue au
voisinage d'une matrice non singulière).

\question Si de plus $\frac{Z'\varepsilon}{\sqrt T}\dlimT\normal{0}{\Psi}$,
alors~:
\[
  \sqrt T\left(\hat{\mathbf b}_{\textsc{vi}}-\beta\right)\dlimT\normal{0}{Q_{ZX}^{-1}\Psi\left(Q_{ZX}'\right)^{-1}}
\]
Dans le cas usuel où les erreurs sont i.i.d.\ et homoscédastiques,
$\Psi = \epsvar Q_{ZZ}$, et la variance asymptotique se simplifie. En pratique on
l'estime par
$\hat\sigma^2(Z'X)^{-1}Z'Z(X'Z)^{-1}$, et les tests de Student et de Fisher ne
valent plus qu'\emph{asymptotiquement}.

\question Avec $L>K$ instruments, le système $Z'(\mathbf y - X\mathbf b) = 0$
compte $L$ équations pour $K$ inconnues~: il est sur-déterminé et n'admet en
général aucune solution exacte. On ne peut donc plus annuler toutes les
conditions de moments~; on cherche à les rendre \emph{aussi petites que
possible} en minimisant une forme quadratique~:
\[
  \hat{\mathbf b} = \arg\min_{\mathbf b}\ \left(Z'\left(\mathbf y-X\mathbf b\right)\right)'W\left(Z'\left(\mathbf y-X\mathbf b\right)\right)
\]
C'est la méthode des moments généralisée. Le choix $W = (Z'Z)^{-1}$ conduit à
l'estimateur des \textbf{doubles moindres carrés}~:
\[
  \hat{\mathbf b}_{\textsc{dmc}} = \left(X'P_ZX\right)^{-1}X'P_Z\mathbf y,
  \qquad P_Z = Z(Z'Z)^{-1}Z'
\]

\nouvelexercice\label{ex:ar1-biais}

\question Le régresseur $y_{t-1}$ n'est manifestement \textbf{pas
déterministe}~: c'est une variable aléatoire, fonction des chocs passés. Il est
\textbf{indépendant de $\varepsilon_t$}, puisque
$y_{t-1} = \sum_{j\geq0}\rho^j\varepsilon_{t-1-j}$ ne fait intervenir que les
innovations antérieures à $t$. En revanche il n'est \textbf{pas indépendant de
$\varepsilon_{t-1}$}, qui figure explicitement dans son écriture. La matrice $X$
des régresseurs et le vecteur $\varepsilon$ ne sont donc pas indépendants~: on
n'est pas dans le «~cas 1~» du cours.

\question L'orthogonalité contemporaine découle de l'indépendance~:
$\mathbb E[y_{t-1}\varepsilon_t] = \mathbb E[y_{t-1}]\mathbb E[\varepsilon_t] = 0$.
Pourtant~:
\[
  \hat\rho = \rho + \frac{\sum_t y_{t-1}\varepsilon_t}{\sum_t y_{t-1}^2}
\]
et $\mathbb E[\hat\rho]\neq\rho$. \textbf{Il n'y a pas de contradiction}~: le
numérateur est bien d'espérance nulle, mais l'estimateur est un \emph{rapport},
et l'espérance d'un rapport n'est pas le rapport des espérances. Le dénominateur
$\sum_t y_{t-1}^2$ est aléatoire et corrélé au numérateur --- un choc
$\varepsilon_{t}$ élevé accroît simultanément le numérateur à la date $t$ et le
dénominateur à la date $t+1$. C'est précisément l'exogénéité stricte qui fait
défaut, et elle seule permettrait de sortir le dénominateur de l'espérance.

\question En divisant par $T$~:
\[
  \hat\rho = \rho + \frac{\frac1T\sum_t y_{t-1}\varepsilon_t}{\frac1T\sum_t y_{t-1}^2}
\]
Le processus étant stationnaire, ses moments ne dépendent pas de la date et la
loi des grands nombres s'applique aux moyennes temporelles~:
$\frac1T\sum_t y_{t-1}\varepsilon_t\plimT\mathbb E[y_{t-1}\varepsilon_t]=0$ et
$\frac1T\sum_t y_{t-1}^2\plimT\mathbb E[y_{t-1}^2]=\epsvar/(1-\rho^2)>0$. Par
Slutsky, $\hat\rho\plimT\rho$~: l'estimateur est \textbf{convergent}. La seule
condition mobilisée est l'\emph{orthogonalité contemporaine}
$\mathbb E[y_{t-1}\varepsilon_t]=0$, bien plus faible que l'exogénéité stricte.

\question Le biais valant approximativement $-2\rho/T$, il est \emph{négatif}
pour $\rho>0$~: l'estimateur des MCO \textbf{sous-estime} la persistance du
processus. Le biais est d'autant plus marqué que $\rho$ est grand et que $T$ est
petit --- ce qui est fâcheux, car c'est justement en présence de forte
persistance qu'on souhaite mesurer $\rho$ avec précision.

\question La simulation confirme l'ordre de grandeur~: pour $\rho=0{,}9$, le
biais théorique $-2\rho/T$ vaut environ $-0{,}072$ pour $T=25$, $-0{,}036$ pour
$T=50$, $-0{,}018$ pour $T=100$ et $-0{,}0036$ pour $T=500$. On observe bien une
moyenne des $\hat\rho$ inférieure à $0{,}9$, l'écart étant divisé par deux
lorsque $T$ double, et devenant négligeable pour $T=500$.

\question C'est le cas \textbf{intermédiaire} entre les deux situations
extrêmes du chapitre~: le cas~1, où $X$ et $\varepsilon$ sont indépendants et où
toutes les propriétés du chapitre~I valent conditionnellement à $X$ (absence de
biais comprise)~; et le cas~2, où $X$ est corrélé contemporainement à
$\varepsilon$ et où l'estimateur n'est même plus convergent. Ici l'estimateur
est \emph{biaisé mais convergent}~: on perd les propriétés à distance finie, on
conserve les propriétés asymptotiques.

\question Si les $\varepsilon_t$ sont eux-mêmes autocorrélés, alors
$\varepsilon_t$ est corrélé à $\varepsilon_{t-1}$, donc à $y_{t-1}$ qui en
dépend~: l'orthogonalité contemporaine $\mathbb E[y_{t-1}\varepsilon_t]=0$
\textbf{tombe}. On bascule dans le cas~2 et l'estimateur des MCO devient
\emph{inconvergent}. C'est un point important~: dans un modèle dynamique,
l'autocorrélation des erreurs n'est plus une simple perte d'efficacité (comme au
chapitre~II) mais une source d'inconvergence. Il faut alors instrumenter (par
exemple par $y_{t-2}$) ou modéliser explicitement la dynamique des erreurs.

\nouvelexercice\label{ex:erreurs-mesure}

\question En remplaçant $x_t = x_t^{*}-u_t$ dans le DGP puis en substituant~:
\[
  y_t^{*} = y_t + v_t = \alpha + \beta\left(x_t^{*}-u_t\right) + \varepsilon_t + v_t
          = \alpha + \beta x_t^{*} + w_t
\]
avec la perturbation composée $w_t = \varepsilon_t + v_t - \beta u_t$.

\question Le régresseur observé et la perturbation composée partagent le terme
$u_t$~:
\[
  \mathrm{cov}\left(x_t^{*},w_t\right)
  = \mathrm{cov}\left(x_t+u_t,\ \varepsilon_t+v_t-\beta u_t\right)
  = -\beta\,\sigma_u^2
\]
tous les autres termes étant nuls par indépendance. Le régresseur est donc
\textbf{endogène} dès que $\beta\neq0$ et $\sigma_u^2>0$~: l'erreur de mesure sur
la variable explicative crée mécaniquement une corrélation entre régresseur et
perturbation.

\question Comme $\mathbb V[x_t^{*}] = \sigma_x^2+\sigma_u^2$~:
\[
  \mathrm{plim}\ \hat b = \beta + \frac{\mathrm{cov}(x^{*},w)}{\mathbb V[x^{*}]}
  = \beta - \frac{\beta\sigma_u^2}{\sigma_x^2+\sigma_u^2}
  = \beta\,\frac{\sigma_x^2}{\sigma_x^2+\sigma_u^2}
\]
Le facteur $\sigma_x^2/(\sigma_x^2+\sigma_u^2)$ appartient à $\left]0,1\right[$~:
la limite est donc de même signe que $\beta$ mais de \textbf{module plus
petit}. D'où le nom d'\emph{atténuation}~: le biais pousse systématiquement
l'estimateur \emph{vers zéro}, et l'on sous-estime l'effet en valeur absolue. Le
cas extrême $\sigma_u^2\gg\sigma_x^2$ donne $\mathrm{plim}\,\hat b\approx0$~: un
régresseur presque entièrement composé de bruit de mesure paraît sans effet.

\question Si $\sigma_u^2 = 0$, le facteur d'atténuation vaut $1$ et
$\mathrm{plim}\,\hat b = \beta$. \textbf{L'erreur de mesure sur la seule variable
expliquée est inoffensive} pour la convergence~: elle se contente de grossir la
variance de la perturbation ($\epsvar+\sigma_v^2$ au lieu de $\epsvar$), donc
les écarts-types, sans introduire de biais. C'est une asymétrie remarquable~:
mal mesurer $y$ coûte de la précision, mal mesurer $x$ coûte la convergence.

\question Vérifions les deux conditions pour $x^{(2)}_t = x_t+u_t^{(2)}$~:
\begin{itemize}
\item \textbf{Pertinence}~:
  $\mathrm{cov}\left(x^{(2)},x^{*}\right) = \mathrm{cov}\left(x+u^{(2)},x+u\right) = \sigma_x^2\neq0$.
  Les deux mesures sont corrélées parce qu'elles mesurent la même grandeur.
\item \textbf{Exogénéité}~:
  $\mathrm{cov}\left(x^{(2)},w\right) = \mathrm{cov}\left(x+u^{(2)},\ \varepsilon+v-\beta u\right) = 0$,
  car $u^{(2)}$ est indépendant de $u$, $v$ et $\varepsilon$, et $x$ est
  indépendant des trois.
\end{itemize}
La seconde mesure est donc un \textbf{instrument valide}. C'est un des rares cas
où la validité de l'instrument se justifie par construction plutôt que par un
argument économique discutable.

\question \textbf{Non.} Le résultat d'atténuation est propre à la régression
simple. En régression multiple, l'erreur de mesure sur une seule variable
contamine \emph{tous} les coefficients, y compris ceux des variables
correctement mesurées, par le jeu des corrélations entre régresseurs. De plus le
sens du biais n'est en général plus prévisible~: on ne peut même plus affirmer
que le coefficient de la variable mal mesurée est atténué.

\question La simulation montre une décroissance monotone de
$\mathrm{plim}\,\hat b/\beta$ en fonction du rapport bruit sur signal
$\sigma_u^2/\sigma_x^2$~: le facteur d'atténuation vaut
$1/\left(1+\sigma_u^2/\sigma_x^2\right)$, soit $0{,}9$ pour un rapport de
$0{,}11$, $0{,}5$ pour un rapport de $1$, et $0{,}1$ pour un rapport de $9$.

\nouvelexercice\label{ex:offre-demande}

\question En égalisant offre et demande~:
\[
  \alpha_0+\alpha_1P_t+\alpha_2R_t+u_t = \gamma_0+\gamma_1P_t+\gamma_2W_t+v_t
\]
soit $\left(\alpha_1-\gamma_1\right)P_t = \left(\gamma_0-\alpha_0\right)+\gamma_2W_t-\alpha_2R_t+\left(v_t-u_t\right)$
et donc~:
\[
  P_t = \frac{\left(\gamma_0-\alpha_0\right)+\gamma_2W_t-\alpha_2R_t+\left(v_t-u_t\right)}{\alpha_1-\gamma_1}
\]
Le prix d'équilibre dépend des \emph{deux} chocs~: c'est la forme réduite.

\question Le prix contient $u_t$, donc~:
\[
  \mathbb E\left[P_tu_t\right] = \frac{-\sigma_u^2}{\alpha_1-\gamma_1}\neq0
\]
Comme $\alpha_1<0<\gamma_1$, on a $\alpha_1-\gamma_1<0$ et cette covariance est
\emph{positive}. Estimer l'équation de demande par les MCO donne donc un
estimateur \textbf{inconvergent} de $\alpha_1$, biaisé vers le haut~: la pente
de la demande paraît moins négative qu'elle ne l'est. \textbf{Intuition}~: un
choc positif de demande fait monter le prix \emph{et} la quantité, ce qui
brouille la relation décroissante que l'on cherche à mesurer. C'est le problème
de \emph{simultanéité}, ou de double causalité.

\question Pour l'\textbf{équation de demande}, il faut un instrument qui déplace
l'offre sans figurer dans la demande~: c'est $W_t$ (la météo). Pertinence~: $W_t$
apparaît dans la forme réduite de $P_t$ avec le coefficient
$\gamma_2/(\alpha_1-\gamma_1)\neq0$. Exogénéité~: $W_t$ n'est pas corrélé à
$u_t$, le choc de demande. Symétriquement, pour l'\textbf{équation d'offre},
l'instrument est $R_t$ (le revenu), qui déplace la demande sans figurer dans
l'offre. Dans les deux cas, la variable exogène propre à l'équation estimée est
son propre instrument~: pour la demande, $Z_t' = \left(1,\ R_t,\ W_t\right)$.
\textbf{C'est le déplacement d'une courbe qui permet d'identifier l'autre.}

\question Avec $L$ instruments et $K$ régresseurs~:
\[
  \hat{\mathbf b}_{\textsc{dmc}} = \left(X'P_ZX\right)^{-1}X'P_Z\mathbf y
\]
Comme $P_Z$ est un projecteur, $X'P_ZX = \hat X'\hat X$ et
$X'P_Z\mathbf y = \hat X'\mathbf y$ avec $\hat X = P_ZX = Z\hat\Pi$ et
$\hat\Pi = (Z'Z)^{-1}Z'X$. D'où la \textbf{lecture en deux étapes}~:
\emph{(i)} régresser chaque colonne de $X$ sur $Z$ et retenir les valeurs
ajustées $\hat X$ --- c'est la part de $X$ expliquée par les instruments, donc
purgée de la corrélation avec $\varepsilon$~; \emph{(ii)} régresser $\mathbf y$
sur $\hat X$.

\question Si $L=K$, la matrice $Z'X$ est carrée et inversible, et~:
\[
  \begin{split}
    \left(X'P_ZX\right)^{-1}X'P_Z\mathbf y
      &= \left(X'Z(Z'Z)^{-1}Z'X\right)^{-1}X'Z(Z'Z)^{-1}Z'\mathbf y\\
      &= \left(Z'X\right)^{-1}\left(Z'Z\right)\left(X'Z\right)^{-1}X'Z(Z'Z)^{-1}Z'\mathbf y\\
      &= \left(Z'X\right)^{-1}Z'\mathbf y = \hat{\mathbf b}_{\textsc{vi}}
  \end{split}
\]
Les doubles moindres carrés généralisent donc bien l'estimateur à variables
instrumentales~: ils coïncident en juste identification.

\question \textbf{Test de Sargan.} On calcule les résidus des doubles moindres
carrés $\hat u = \mathbf y - X\hat{\mathbf b}_{\textsc{dmc}}$, on les régresse
sur l'ensemble des instruments $Z$, et on forme~:
\[
  S = T R^2 \dlimT \chi^2(L-K)
\]
Un $R^2$ élevé signale que les instruments expliquent encore les résidus, donc
qu'ils ne sont pas orthogonaux à la perturbation. \textbf{Ce que l'on teste
exactement}~: seulement les $L-K$ restrictions \emph{sur-identifiantes},
c'est-à-dire la \emph{cohérence mutuelle} des instruments. Les $K$ conditions
nécessaires à l'identification ne sont pas testables~: si tous les instruments
étaient invalides «~de la même façon~», le test ne détecterait rien. D'où les
$L-K$ degrés de liberté, et l'impossibilité de tester quoi que ce soit en juste
identification.

\question \textbf{Non.} Le test est un test global~: un rejet indique qu'au
moins une restriction sur-identifiante est violée, sans désigner le coupable. Il
est d'ailleurs possible qu'un instrument valide soit «~accusé~» parce qu'il est
comparé à des instruments invalides. Réciproquement, un non-rejet n'est pas une
preuve de validité~: il peut refléter un manque de puissance, ou le cas
mentionné ci-dessus où tous les instruments partagent le même défaut. Le test de
Sargan est un garde-fou, pas une certification.

\nouvelexercice\label{ex:instruments-faibles}

\question En reprenant la question (5) de l'exercice~\ref{ex:vi-simple} et en notant que
$\mathrm{cov}(z,x) = \pi\sigma_z^2$~:
\[
  \mathrm{plim}\ \hat b_{\textsc{vi}} = \beta + \frac{\mathrm{cov}(z,\varepsilon)}{\pi\,\sigma_z^2}
\]
Si l'instrument est \emph{exactement} exogène, $\mathrm{cov}(z,\varepsilon)=0$ et
la limite vaut $\beta$ quel que soit $\pi\neq0$. Mais si l'exogénéité n'est
qu'approximativement satisfaite, le terme de biais est divisé par $\pi$~: il
\textbf{explose lorsque $\pi\rightarrow0$}. Un instrument faible amplifie
démesurément la moindre violation de l'exogénéité --- au point que
l'estimateur à variables instrumentales peut être plus biaisé que les MCO.

\question La variance asymptotique vaut $\epsvar\mathbb V[z]/\mathrm{cov}(z,x)^2 = \epsvar/(\pi^2\sigma_z^2)$,
qui diverge en $1/\pi^2$ lorsque $\pi\rightarrow0$. De façon équivalente, le
rapport $1/\rho_{xz}^2$ de l'exercice~\ref{ex:vi-simple} tend vers l'infini. L'estimateur
devient donc à la fois très imprécis et très sensible.

\question Le résultat rappelé dans l'énoncé montre que
$\mathbb E[Z_t\varepsilon_t]=0$ \textbf{n'implique pas} que la moyenne empirique
$\frac1T\sum_t Z_t\varepsilon_t$ converge vers quoi que ce soit~: dans cet
exemple, le produit $Z_t\varepsilon_t$ suit une loi de Cauchy, qui n'a pas
d'espérance, et la loi des grands nombres ne s'applique pas.
\textbf{Leçon.} La non-corrélation est un énoncé sur un moment de population~;
la convergence exige en outre que ce moment \emph{existe} et que la loi des
grands nombres soit applicable. Les conditions du théorème des variables
instrumentales sont donc formulées en termes de convergence
($Z'\varepsilon/T\plimT0$), et non de simple non-corrélation --- ce n'est pas une
coquetterie de rédaction.

\question \textbf{Non.} En juste identification, l'estimateur à variables
instrumentales est un rapport de deux formes quadratiques corrélées, et ne
possède \emph{aucun moment} --- pas même une espérance. Parler de son «~biais~»
n'a donc littéralement pas de sens à distance finie. C'est une différence
essentielle avec les MCO, et cela explique que les simulations sur instruments
faibles produisent des distributions à queues très épaisses, parfois bimodales.

\question On utilise la statistique de \textbf{Fisher de la première étape},
testant la nullité conjointe des coefficients des instruments dans la régression
de $x$ sur $Z$. La règle empirique usuelle est de considérer les instruments
comme faibles si $F<10$. \textbf{Limites}~: ce seuil provient d'un calcul
particulier (borner le biais relatif des doubles moindres carrés à $10\,\%$ dans
un cadre homoscédastique) et n'a aucune validité générale~; il se dégrade
lorsque les instruments sont nombreux ou les erreurs hétéroscédastiques~; enfin,
sélectionner le modèle sur la base de $F$ puis conduire l'inférence comme si de
rien n'était introduit un biais de pré-test.

\question \textbf{Test de Hausman.} L'idée est de comparer deux estimateurs~:
$\hat{\mathbf b}$ (MCO), efficace mais convergent seulement sous exogénéité, et
$\hat{\mathbf b}_{\textsc{vi}}$, convergent dans les deux cas mais moins précis.
Sous $\mathcal H_0$ (exogénéité) les deux convergent vers $\beta$ et leur écart
tend vers zéro~; sous $\mathcal H_1$ ils divergent. D'où~:
\[
  H = \left(\hat{\mathbf b}_{\textsc{vi}}-\hat{\mathbf b}\right)'
      \left[\hat{\mathbb V}\left[\hat{\mathbf b}_{\textsc{vi}}\right]-\hat{\mathbb V}\left[\hat{\mathbf b}\right]\right]^{-}
      \left(\hat{\mathbf b}_{\textsc{vi}}-\hat{\mathbf b}\right)
  \dlimT \chi^2(K_1)
\]
où $K_1$ est le nombre de régresseurs dont on teste l'exogénéité.
\textbf{Pourquoi un inverse généralisé}~: sous $\mathcal H_0$, $\hat{\mathbf b}$
est efficace, ce qui garantit que la différence des variances est semi-définie
positive, mais elle est \emph{singulière} --- son rang vaut $K_1$ et non $K$,
car les composantes correspondant aux régresseurs exogènes sont estimées
identiquement par les deux méthodes. L'inverse ordinaire n'existe donc pas.

\question \textbf{Variante de Durbin-Wu.} On régresse $x$ sur l'ensemble des
instruments et l'on récupère les résidus $\hat{\mathbf v}$ de cette première
étape~; on ajoute $\hat{\mathbf v}$ comme régresseur supplémentaire dans
l'équation d'intérêt~:
\[
  y_t = a + b\,x_t + c\,\hat v_t + \text{erreur}
\]
et l'on teste $\mathcal H_0: c=0$ par un simple $t$ de Student. En effet, $c$ ne
peut être non nul que si la part de $x$ non expliquée par les instruments est
corrélée à la perturbation, c'est-à-dire si $x$ est endogène.
\textbf{Avantages pratiques}~: le test se ramène à une statistique de Student
standard, il évite le calcul d'un inverse généralisé, il ne souffre pas du
problème d'une différence de variances estimée non semi-définie positive en
petit échantillon (qui rend parfois $H$ négatif), et il se robustifie
immédiatement à l'hétéroscédasticité en utilisant un écart-type robuste.

\nouvelexercice\label{ex:wald}

\question D'après l'exercice~\ref{ex:vi-simple}~:
\[
  \hat b_{\textsc{vi}} = \frac{\sum_i\left(z_i-\bar z\right)\left(y_i-\bar y\right)}{\sum_i\left(z_i-\bar z\right)\left(x_i-\bar x\right)}
\]

\question Notons $N_1 = Np$ le nombre d'observations telles que $z_i=1$ et
$N_0 = N(1-p)$ les autres, de sorte que $\bar z = p$. En séparant la somme selon
les deux groupes, et comme $z_i-\bar z$ vaut $1-p$ sur le premier et $-p$ sur le
second~:
\[
  \begin{split}
    \sum_i\left(z_i-\bar z\right)\left(y_i-\bar y\right)
      &= (1-p)\sum_{i:z_i=1}\left(y_i-\bar y\right) - p\sum_{i:z_i=0}\left(y_i-\bar y\right)\\
      &= (1-p)N_1\left(\bar y_1-\bar y\right) - pN_0\left(\bar y_0-\bar y\right)
  \end{split}
\]
Or $\bar y = p\bar y_1+(1-p)\bar y_0$, donc
$\bar y_1-\bar y = (1-p)(\bar y_1-\bar y_0)$ et
$\bar y_0-\bar y = -p(\bar y_1-\bar y_0)$. En reportant~:
\[
  \sum_i\left(z_i-\bar z\right)\left(y_i-\bar y\right)
  = \left[(1-p)^2Np + p^2N(1-p)\right]\left(\bar y_1-\bar y_0\right)
  = N\,p(1-p)\left(\bar y_1-\bar y_0\right)
\]
puisque $(1-p)^2p+p^2(1-p) = p(1-p)$. Le même calcul avec $x$ donne
$N\,p(1-p)\left(\bar x_1-\bar x_0\right)$.

\question Le facteur commun $Np(1-p)$ se simplifie dans le rapport~:
\[
  \hat b_{\textsc{vi}} = \frac{\bar y_1-\bar y_0}{\bar x_1-\bar x_0}
\]
C'est l'\textbf{estimateur de Wald}.

\question L'estimateur est le rapport de deux différences de moyennes~: au
numérateur, l'écart de $y$ entre les deux groupes définis par l'instrument~; au
dénominateur, l'écart de $x$ entre ces mêmes groupes. \textbf{Interprétation.}
L'instrument déplace $x$ d'une quantité $\bar x_1-\bar x_0$, ce qui déplace $y$
d'une quantité $\bar y_1-\bar y_0$~: le rapport mesure «~combien de $y$ par unité
de $x$~», en n'utilisant que la variation de $x$ imputable à $z$. La
\emph{condition de pertinence} $\mathrm{cov}(z,x)\neq0$ devient simplement
$\bar x_1\neq\bar x_0$~: l'instrument doit effectivement déplacer $x$. On voit
aussi immédiatement pourquoi la division est dangereuse si ce dénominateur est
petit.

\question Ici $x_i\in\{0,1\}$, donc $\bar x_1$ est le taux de participation
parmi les personnes tirées au sort et $\bar x_0$ celui parmi les autres.
\begin{itemize}
\item \textbf{Numérateur}~: $\bar y_1-\bar y_0$ compare le résultat moyen des
  personnes \emph{à qui l'on a proposé} la formation à celui des autres.
  Comme $z$ est tiré au sort, les deux groupes sont comparables~: cette
  différence est l'effet de l'\emph{intention de traiter}. C'est ce que l'on
  mesurerait en régressant $y$ sur $z$.
\item \textbf{Dénominateur}~: $\bar x_1-\bar x_0$ est l'écart des taux de
  participation effective entre les deux groupes, c'est-à-dire le
  \emph{taux de conformité} au tirage. C'est la régression de première étape.
\end{itemize}
L'estimateur de Wald corrige donc l'effet d'intention de traiter par le taux de
conformité~: si seuls $50\,\%$ des personnes sollicitées participent
effectivement, l'effet mesuré sur l'ensemble du groupe doit être doublé pour
retrouver l'effet sur ceux qui participent réellement.

\question Si le taux de conformité est faible, le dénominateur $\bar x_1-\bar x_0$
est proche de zéro~: l'estimateur devient très instable, sa variance
asymptotique explose, et le moindre défaut d'exogénéité de $z$ est amplifié dans
le rapport. C'est exactement la situation d'instruments faibles de
l'exercice~\ref{ex:instruments-faibles}~: on y voit ici de façon très concrète
d'où vient le problème, puisqu'il s'agit littéralement d'une division par un
petit nombre. À la limite $\bar x_1=\bar x_0$ (personne ne modifie son
comportement suite au tirage), l'estimateur n'est pas défini~: l'expérience
n'apporte aucune information sur $\beta$.

\question \textbf{Le tirage au sort porte sur $z$, pas sur $x$.} Par
construction, $z$ est indépendant de tout ce qui caractérise les individus,
donc de $\varepsilon$~: l'exogénéité de l'instrument est garantie \emph{par le
protocole}, ce qui est rare et précieux --- ailleurs, elle repose toujours sur un
argument discutable. En revanche $x$, la participation \emph{effective},
résulte d'une décision individuelle~: les personnes qui acceptent la formation
sont vraisemblablement les plus motivées, et cette motivation influence aussi
$y$. Donc $\mathrm{cov}(x,\varepsilon)\neq0$. Régresser $y$ sur $x$ produirait
un estimateur inconvergent (exercice~\ref{ex:vi-simple}), confondant l'effet de
la formation avec celui de la motivation~: on attribuerait au programme le mérite
de sélectionner les meilleurs. Le tirage au sort ne suffit donc pas~: encore
faut-il l'utiliser comme \emph{instrument} et non comme traitement.

\section*{Chapitre IV~-- Hétéroscédasticité}

\nouvelexercice\label{ex:cout-hetero}

\question En substituant le DGP,
$\hat{\mathbf b} = \beta + (X'X)^{-1}X'\varepsilon$. Comme $X$ est déterministe
et $\mathbb E[\varepsilon]=0$, $\mathbb E[\hat{\mathbf b}]=\beta$~:
\textbf{l'absence de biais est préservée}, car elle ne fait intervenir que
l'espérance des erreurs, jamais leur variance. L'estimateur reste convergent dès
que $X'X/T\plimT Q$ définie positive et que $X'\Omega X/T$ admet une limite
finie.

\question Comme $\mathbb V[\varepsilon]=\Omega$~:
\[
  \mathbb V\left[\hat{\mathbf b}\right] = (X'X)^{-1}X'\Omega X(X'X)^{-1}
\]
Elle diffère de $\epsvar(X'X)^{-1}$ parce que le calcul usuel remplace
$\mathbb V[\varepsilon]$ par $\epsvar I_T$, ce qui permet de sortir $\epsvar$ du
produit et de simplifier $X'X$. Ici $\Omega$ ne peut pas être sortie de
$X'\Omega X$, et la simplification n'a plus lieu. C'est le cas diagonal de la
variance sandwich du chapitre~II.

\question \textbf{Non.} Le théorème de Gauss-Markov suppose
$\mathbb V[\varepsilon]=\epsvar I_T$, c'est-à-dire à la fois l'homoscédasticité
(éléments diagonaux tous égaux) et l'absence de corrélation (éléments hors
diagonale nuls). C'est ici la première partie qui est violée. Il existe donc un
estimateur linéaire sans biais de variance plus petite~: l'estimateur des MCG,
qui prend ici la forme des moindres carrés pondérés (exercice~\ref{ex:mcp}).

\question \textbf{Non}, $s^2(X'X)^{-1}$ n'estime rien de pertinent~: $s^2$
converge vers une moyenne pondérée des $\sigma_t^2$, et le produit par
$(X'X)^{-1}$ ne reconstitue pas $X'\Omega X$ encadré. Il n'y a aucune raison
pour que le résultat soit proche de la vraie variance --- il peut la sur-estimer
comme la sous-estimer. Par conséquent les tests de Student et de Fisher du
chapitre~I sont invalides, et ils le restent \textbf{asymptotiquement}~: ce n'est
pas un problème de petit échantillon que davantage d'observations résoudrait.

\question \textbf{Le danger n'est pas dans l'estimateur, qui reste juste~: il
est dans la formule d'écart-type imprimée à côté.} C'est la phrase à retenir du
chapitre. Elle explique aussi pourquoi la première réponse au problème
(exercice~\ref{ex:hc}) consiste à ne rien changer à $\hat{\mathbf b}$ et à corriger
uniquement son écart-type.

\question Si $y_g$ est la moyenne de $n_g$ observations individuelles
homoscédastiques de variance $\epsvar$ et non corrélées entre elles~:
\[
  \mathbb V\left[\varepsilon_g\right] = \mathbb V\left[\frac{1}{n_g}\sum_{j=1}^{n_g}\varepsilon_{gj}\right] = \frac{\epsvar}{n_g}
\]
Les erreurs restent non corrélées entre groupes, donc $\Omega$ est bien
diagonale, mais ses éléments diffèrent dès que les groupes n'ont pas la même
taille. C'est un cas où $\Omega$ est \emph{connue} à un facteur près, et où la
pondération de l'exercice~\ref{ex:mcp} s'impose naturellement (chaque groupe est pondéré
par $\sqrt{n_g}$).

\question Trois exemples classiques~: \emph{(i)} les équations de consommation
ou d'épargne sur données individuelles, où la dispersion des comportements croît
avec le revenu --- les ménages aisés ont beaucoup plus de latitude~; \emph{(ii)}
les modèles portant sur des entreprises de tailles très différentes, où la
variabilité des résultats croît avec l'échelle d'activité~; \emph{(iii)} les
données agrégées ou moyennées par zone géographique, comme dans la question
précédente. On peut y ajouter les modèles à variable dépendante limitée, où la
variance dépend mécaniquement de l'espérance.

\nouvelexercice\label{ex:breusch-pagan}

\question Le test est un test du multiplicateur de Lagrange, évalué
\emph{sous} $\mathcal H_0:\alpha=0$. Au voisinage de ce point, seul compte le
développement au premier ordre de $h$~: la dérivée $h'(\alpha_0)$ intervient
comme un facteur d'échelle, qui se simplifie dans la statistique une fois
celle-ci normalisée. La seule exigence est donc $h'(\alpha_0)\neq0$, c'est-à-dire
que la variance réagisse effectivement à l'indice. Le test a ainsi la même forme
que $\sigma_t^2$ soit linéaire, exponentielle ou quadratique en $z_t'\alpha$~:
c'est ce qui fait sa robustesse et son élégance.

\question On régresse la variable
$\hat\varepsilon_t^{\,2}/\hat\sigma^2$ sur une constante et sur $z_t$, où
$\hat\sigma^2 = \frac1T\sum_t\hat\varepsilon_t^{\,2}$ (noter la division par $T$
et non par $T-K$). On retient la somme des carrés expliquée $SCE$ de cette
régression auxiliaire.

\question Sous normalité, $\varepsilon_t/\sigma_t\sim\normal{0}{1}$, donc
$\left(\varepsilon_t/\sigma_t\right)^2\sim\chi^2(1)$. Or la variance d'une loi
du khi-deux à un degré de liberté vaut $2\times1 = 2$~:
\[
  \mathbb V\left[\frac{\varepsilon_t^2}{\sigma_t^2}\right] = 2
\]
\textbf{D'où vient le facteur $2$}~: la variable expliquée de la régression
auxiliaire a, sous $\mathcal H_0$, une variance \emph{connue} et égale à $2$. La
statistique $LM$ est la somme des carrés expliquée divisée par cette variance,
comme toute statistique de score. C'est parce que cette variance est connue --- et
qu'elle vaut $2$ --- que la normalité est indispensable.

\question Sous $\mathcal H_0$~:
\[
  LM = \frac{SCE}{2}\ \dlimT\ \chi^2(p)
\]
où $p$ est le nombre de variables dans $z_t$ (hors constante). On rejette
l'homoscédasticité si $LM$ dépasse le quantile d'ordre $1-\alpha$ de cette loi.

\question Si les erreurs ne sont pas gaussiennes,
$\mathbb V[\varepsilon_t^2/\sigma_t^2]\neq2$ en général (elle dépend du moment
d'ordre quatre), et diviser par $2$ fausse la statistique~: le test n'a plus la
taille annoncée, et il sur-rejette typiquement en présence de queues épaisses.
\textbf{Version de Koenker.} On régresse directement $\hat\varepsilon_t^{\,2}$
sur $(1,z_t)$ et l'on retient~:
\[
  LM_K = T R^2 \ \dlimT\ \chi^2(p)
\]
Le $R^2$ normalise implicitement par la variance \emph{empirique} de
$\hat\varepsilon_t^{\,2}$ au lieu de lui imposer la valeur théorique $2$~: c'est
en ce sens que le test est «~studentisé~». Il ne requiert plus la normalité.

\question Sous normalité, la version originale est \emph{plus puissante}~: elle
exploite une information vraie (la variance vaut exactement $2$) au lieu de
l'estimer. La version de Koenker est \emph{plus robuste}~: elle conserve la
bonne taille quelle que soit la loi des erreurs. En pratique, la normalité étant
rarement crédible sur données individuelles, c'est la version de Koenker qui est
utilisée par défaut dans les logiciels.

\question \textbf{Non.} Le test n'a de puissance que contre les alternatives où
la variance dépend de $z_t$ à travers l'indice linéaire $z_t'\alpha$. Une
hétéroscédasticité qui dépendrait d'une variable absente de $z_t$, ou qui
dépendrait de $z_t$ de façon non monotone (par exemple en $z_t^2$ avec un
minimum intérieur, de sorte que la corrélation entre $\varepsilon^2$ et $z$ soit
nulle), passerait inaperçue. C'est précisément ce qui motive le test de White.

\nouvelexercice\label{ex:white}

\question On régresse le carré des résidus des MCO, $\hat\varepsilon_t^{\,2}$,
sur une constante, sur les $J=K-1$ régresseurs du modèle initial, sur leurs
carrés, et sur tous leurs produits croisés deux à deux.

\question Supposons que la variance soit une fonction inconnue des explicatives,
$\sigma_t^2 = f(\mathbf x_t)$. Un développement de Taylor d'ordre deux au
voisinage de $\bar{\mathbf x}$ donne~:
\[
  f(\mathbf x_t) \simeq f(\bar{\mathbf x})
  + \nabla f(\bar{\mathbf x})'\left(\mathbf x_t-\bar{\mathbf x}\right)
  + \tfrac12\left(\mathbf x_t-\bar{\mathbf x}\right)'\nabla^2 f(\bar{\mathbf x})\left(\mathbf x_t-\bar{\mathbf x}\right)
\]
En développant, on fait apparaître une constante, les niveaux $x_{k,t}$, les
carrés $x_{k,t}^2$ et les produits croisés $x_{k,t}x_{l,t}$. La régression
auxiliaire de White est donc l'approximation quadratique générique d'une
fonction de variance \emph{quelconque}~: c'est ce qui lui donne sa généralité.

\question On compte $J$ niveaux, $J$ carrés et $\binom{J}{2} = J(J-1)/2$ produits
croisés distincts, soit~:
\[
  q = J + J + \frac{J(J-1)}{2} = J + \frac{J(J+1)}{2}
\]

\question La statistique et sa loi asymptotique sous $\mathcal H_0$~:
\[
  W = TR^2\ \dlimT\ \chi^2(q)
\]
où $R^2$ est celui de la régression auxiliaire.

\question Avec $J=3$ régresseurs ($educ$, $exper$, $exper^2$), le décompte naïf
donnerait $q = 3 + \frac{3\times4}{2} = 9$. Mais le carré de $exper$ est
$exper^2$, qui figure \emph{déjà} parmi les niveaux~: la colonne
«~$exper\times exper$~» de la régression auxiliaire est exactement colinéaire à
la colonne «~$exper^2$~». Une colonne est donc redondante et doit être écartée~:
\[
  q = 9 - 1 = 8
\]
\textbf{Leçon générale.} Dès que le modèle contient une transformation
polynomiale d'un régresseur, ou une indicatrice (dont le carré est elle-même),
la régression auxiliaire de White comporte des colinéarités et le degré de
liberté effectif est inférieur à la formule. Il faut lire le rang effectif de la
matrice auxiliaire, ce que font les logiciels --- et ne pas appliquer la formule
mécaniquement.

\question Le test de White est \textbf{plus général}~: il ne suppose ni forme
fonctionnelle particulière, ni normalité, ni le choix préalable de variables
$z_t$. Le test de Breusch-Pagan est \textbf{plus puissant} lorsque la forme de
l'hétéroscédasticité est effectivement celle qu'il postule, car il concentre sa
puissance sur $p$ degrés de liberté au lieu de la disperser sur $q$, en général
bien plus grand. C'est l'arbitrage habituel entre généralité et puissance.

\question \textbf{Non, il faut être prudent.} La régression auxiliaire de White
contient les niveaux, les carrés et les produits croisés des explicatives~: ce
sont précisément les termes qu'on aurait ajoutés au modèle initial pour tester
une \emph{erreur de forme fonctionnelle}. Si le vrai modèle est non linéaire et
qu'on l'estime linéairement, la partie non modélisée se retrouve dans les
résidus et sera captée par la régression auxiliaire, provoquant un rejet sans
qu'il y ait la moindre hétéroscédasticité. Le test de White est donc, en
pratique, un test de mauvaise spécification au sens large. Avant de conclure à
l'hétéroscédasticité, il est prudent de vérifier la forme fonctionnelle du
modèle.

\nouvelexercice\label{ex:goldfeld-quandt}

\question On ordonne les observations selon $z$, croissant. On écarte les $c$
observations centrales, puis on estime le modèle \emph{séparément} sur le
premier sous-échantillon (faibles $z$) et sur le second (forts $z$), chacun de
taille $m = (T-c)/2$. On compare les deux sommes des carrés des résidus.

\question \textbf{Pourquoi écarter les observations centrales}~: elles sont
celles pour lesquelles la variance est intermédiaire~; les inclure atténuerait le
contraste entre les deux groupes et rapprocherait le rapport de $1$ sous
$\mathcal H_1$. Les retirer accentue la différence et augmente la puissance.
\textbf{Effet du choix de $c$}~: augmenter $c$ accroît le contraste mais réduit
$m$, donc les degrés de liberté et la précision de chaque $SCR$. Il y a un
optimum intermédiaire, d'où la recommandation $c\simeq T/4$.

\question Sous $\mathcal H_0$ et sous normalité, chaque sous-échantillon relève
du chapitre~I~: $SCR_i/\epsvar\sim\chi^2(m-K)$. Les deux sous-échantillons étant
disjoints et les erreurs indépendantes, les deux formes quadratiques sont
indépendantes. Le rapport de deux khi-deux indépendants divisés par leurs degrés
de liberté suit une loi de Fisher, et $\epsvar$ se simplifie~:
\[
  F = \frac{SCR_2/(m-K)}{SCR_1/(m-K)} \sim \mathcal F(m-K,\,m-K)
\]

\question C'est une loi \textbf{exacte}, valable à toute taille d'échantillon --- 
c'est le seul des trois tests dans ce cas, les deux autres n'ayant qu'une
justification asymptotique. En contrepartie, la \textbf{normalité des erreurs
est indispensable}~: sans elle, $SCR_i/\epsvar$ ne suit plus une loi du khi-deux.

\question Le test est \textbf{unilatéral} parce que l'hypothèse alternative est
orientée~: on postule que la variance \emph{croît} avec $z$, donc que
$SCR_2>SCR_1$. On ordonne les sous-échantillons de sorte que le second
corresponde aux fortes valeurs de $z$, et l'on rejette pour les grandes valeurs
de $F$. Si l'on soupçonnait une variance décroissante, il suffirait d'inverser
l'ordre (ou de placer au numérateur la plus grande des deux sommes, en doublant
le seuil).

\question La \textbf{limite principale} est la double exigence de monotonie et
d'unidimensionnalité~: il faut savoir à l'avance selon quelle variable ordonner,
et postuler que la variance en est une fonction monotone. Le test est aveugle à
toute hétéroscédasticité dépendant de plusieurs variables, ou non monotone. Il
exige de surcroît la normalité et sacrifie $c$ observations. Les tests de
Breusch-Pagan et de White n'ont aucune de ces contraintes.

\question
\begin{center}
\begin{tabular}{p{2.4cm}p{3.3cm}p{1.9cm}p{2.0cm}p{2.1cm}}
\hline
 & \textbf{Hétéroscéd. testée} & \textbf{Normalité} & \textbf{Loi sous $\mathcal H_0$} & \textbf{Degrés de liberté}\\
\hline
Breusch-Pagan (original) & $\sigma_t^2=h(\alpha_0+z_t'\alpha)$, forme d'indice & requise & $\chi^2$ asympt. & $p$\\[1mm]
Koenker & idem & non requise & $\chi^2$ asympt. & $p$\\[1mm]
White & quelconque (approx. quadratique) & non requise & $\chi^2$ asympt. & $q=J+\frac{J(J+1)}{2}$, moins les colinéarités\\[1mm]
Goldfeld-Quandt & monotone en une variable $z$ & requise & $\mathcal F$ exacte & $(m-K,\,m-K)$\\
\hline
\end{tabular}
\end{center}

\nouvelexercice\label{ex:hc}

\question On pourrait croire qu'il faut estimer les $T$ variances
$\sigma_1^2,\dots,\sigma_T^2$ à partir de $T$ observations~: le nombre de
paramètres croîtrait aussi vite que l'échantillon et aucune convergence ne
serait possible. \textbf{Mais ce n'est pas ce dont on a besoin.} La variance de
$\hat{\mathbf b}$ ne fait intervenir les $\sigma_t^2$ qu'à travers la matrice
$K\times K$~:
\[
  S = \frac{1}{T}\sum_{t=1}^T \sigma_t^2\,\mathbf x_t\mathbf x_t'
\]
qui ne contient que $K(K+1)/2$ paramètres distincts, indépendamment de $T$.
C'est une moyenne, et une moyenne s'estime bien. Toute l'astuce de White est là.

\question On remplace chaque $\sigma_t^2$ inconnue par le carré du résidu
observé~:
\[
  \widehat{\mathbb V}_{HC0}
  = \left(\sum_t \mathbf x_t\mathbf x_t'\right)^{-1}
    \left(\sum_t \hat\varepsilon_t^{\,2}\,\mathbf x_t\mathbf x_t'\right)
    \left(\sum_t \mathbf x_t\mathbf x_t'\right)^{-1}
\]
Pris isolément, $\hat\varepsilon_t^{\,2}$ est un très mauvais estimateur de
$\sigma_t^2$ (un seul point pour une variance !), mais c'est sans importance~:
seule la \emph{moyenne} pondérée doit converger.

\question Le résultat du cours est~: si $X'X/T\plimT Q$ définie positive et
$\frac1T\sum_t\sigma_t^2\mathbf x_t\mathbf x_t'\plimT S$, alors~:
\[
  T\,\widehat{\mathbb V}_{HC0} \plimT Q^{-1}SQ^{-1}
\]
\textbf{Attention}~: l'énoncé porte sur $T\widehat{\mathbb V}_{HC0}$, et non sur
$\widehat{\mathbb V}_{HC0}$. Il ne peut pas en être autrement~:
$\widehat{\mathbb V}_{HC0}$ et $\mathbb V[\hat{\mathbf b}]$ tendent toutes deux
vers la matrice nulle quand $T\rightarrow\infty$, et dire de l'une qu'elle
converge vers l'autre serait vide de sens. La normalisation par $T$ est ce qui
donne un contenu à l'énoncé --- c'est la même précaution que pour
$\sqrt T(\hat{\mathbf b}-\beta)$.

\question Pour la trace~:
\[
  \sum_t h_{tt} = \trace(H) = \tracarg{X(X'X)^{-1}X'} = \tracarg{(X'X)^{-1}X'X} = \trace(I_K) = K
\]
Le levier moyen vaut donc $K/T$. Par ailleurs $\hat\varepsilon = M\varepsilon$
avec $M = I-H$, donc sous homoscédasticité~:
\[
  \mathbb E\left[\hat\varepsilon_t^{\,2}\right] = \epsvar\,M_{tt} = \epsvar\left(1-h_{tt}\right)
\]

\question Comme $0\leq h_{tt}\leq1$, on a
$\mathbb E[\hat\varepsilon_t^{\,2}]\leq\epsvar$~: les résidus estimés sont
\emph{systématiquement trop petits}, parce que les MCO ajustent la droite en
partie sur le bruit. HC0, qui les utilise tels quels, \textbf{sous-estime} donc
la variance. Le biais est d'autant plus marqué que le levier $h_{tt}$ est élevé,
c'est-à-dire pour les observations atypiques --- précisément celles qui pèsent le
plus dans la somme. En petit échantillon, les tests fondés sur HC0 rejettent
donc trop souvent.

\question En corrigeant chaque terme~:
\[
  HC1:\ \frac{T}{T-K}\,\hat\varepsilon_t^{\,2}
  \qquad
  HC2:\ \frac{\hat\varepsilon_t^{\,2}}{1-h_{tt}}
  \qquad
  HC3:\ \frac{\hat\varepsilon_t^{\,2}}{\left(1-h_{tt}\right)^2}
\]
HC1 applique la même correction de degrés de liberté à toutes les observations.
HC2 corrige exactement le biais calculé à la question (4), observation par
observation. HC3 (approximation du \emph{jackknife}) sur-corrige délibérément~:
comme $(1-h_{tt})^2 < 1-h_{tt}$, c'est la correction la plus forte des trois.

\question La règle du cours est~: retenir HC3 si $T<250$, ou si
$\max_t h_{tt}$ est nettement supérieur à $2K/T$ (présence de points à fort
levier). Dans le doute, HC3 est le choix par défaut prudent~: sur grand
échantillon les quatre variantes coïncident, de sorte qu'on ne perd rien à la
retenir systématiquement. \textbf{Simulation.} Pour $T=30$ et une
hétéroscédasticité marquée, un test au seuil nominal de $5\,\%$ rejette de
l'ordre de $10$ à $12\,\%$ du temps avec HC0, $8$ à $9\,\%$ avec HC1, et revient
au voisinage de $5$ à $6\,\%$ avec HC3. Les écarts-types usuels, eux, peuvent
donner des taux de rejet bien plus élevés encore.

\nouvelexercice\label{ex:mcp}

\question Si $\Omega = \mathrm{diag}(\sigma_1^2,\dots,\sigma_T^2)$, alors
$\Omega^{-1} = \mathrm{diag}(1/\sigma_1^2,\dots,1/\sigma_T^2)$ et les produits
matriciels se lisent comme des sommes pondérées~:
\[
  X'\Omega^{-1}X = \sum_{t=1}^T\frac{\mathbf x_t\mathbf x_t'}{\sigma_t^2},
  \qquad
  X'\Omega^{-1}\mathbf y = \sum_{t=1}^T\frac{\mathbf x_ty_t}{\sigma_t^2}
\]
d'où l'expression annoncée. Chaque observation entre dans les sommes affectée du
poids $1/\sigma_t^2$~: d'où le nom de \emph{moindres carrés pondérés}. De façon
équivalente, l'estimateur minimise $\sum_t (y_t-\mathbf x_t'\mathbf b)^2/\sigma_t^2$,
somme des carrés des résidus pondérée.

\question Il suffit de diviser chaque observation par son écart-type~:
\[
  \tilde y_t = \frac{y_t}{\sigma_t},
  \qquad
  \tilde{\mathbf x}_t = \frac{\mathbf x_t}{\sigma_t}
\]
ce qui correspond à $\Lambda = \Omega^{-1/2} = \mathrm{diag}(1/\sigma_t)$. Les
erreurs transformées $\varepsilon_t/\sigma_t$ ont alors une variance unitaire
pour tout $t$~: le modèle transformé est homoscédastique, et les MCO qu'on y
applique satisfont les conditions idéales.

\question Les observations de \emph{faible} variance reçoivent le poids le plus
élevé. C'est parfaitement intuitif~: une observation dont on sait qu'elle est
mesurée avec précision est plus informative sur la position de la droite qu'une
observation très dispersée, et doit peser davantage. Les MCO, en donnant le même
poids à toutes, gaspillent cette information --- d'où leur perte d'efficacité.

\question Parce que la fonction exponentielle est strictement positive~:
$\sigma_t^2 = \sigma^2\exp(z_t'\gamma)>0$ quels que soient $\gamma$ et $z_t$. Une
spécification linéaire $\sigma_t^2 = \sigma^2 + z_t'\gamma$ ne le garantirait pas
et pourrait produire des variances estimées négatives, donc des poids négatifs
et un estimateur dépourvu de sens. La forme multiplicative rend en outre
l'estimation linéaire après passage au logarithme.

\question \textbf{Procédure en quatre étapes.} \emph{(i)} Estimer le modèle par
les MCO et récupérer les résidus $\hat\varepsilon_t$. \emph{(ii)} Régresser
$\log\hat\varepsilon_t^{\,2}$ sur une constante et $z_t$~: la pente estime
$\gamma$. \emph{(iii)} Former les variances ajustées
$\hat\sigma_t^2 = \exp(\hat c + z_t'\hat\gamma)$. \emph{(iv)} Estimer le modèle
par moindres carrés pondérés avec les poids $1/\hat\sigma_t^2$.

\question La constante $-1{,}2704$ (qui n'est autre que
$\mathbb E[\log\chi^2(1)]$) est \emph{la même pour toutes les observations}~:
elle ne modifie donc que la \textbf{constante} de la régression auxiliaire, et
laisse les pentes $\gamma$ inchangées. Quant à $\breve{\mathbf b}$, il n'est pas
affecté non plus~: multiplier tous les poids par une même constante positive
$\lambda$ multiplie $X'\Omega^{-1}X$ et $X'\Omega^{-1}\mathbf y$ par $\lambda$,
facteur qui se simplifie dans le produit
$\left(X'\Omega^{-1}X\right)^{-1}X'\Omega^{-1}\mathbf y$. \textbf{Seuls les
poids relatifs comptent}, jamais leur échelle --- ce qui est heureux, puisque
$\sigma^2$ lui-même est inconnu.

\question \textbf{Le risque} est que le modèle de variance soit faux. Les poids
sont alors mal choisis. L'estimateur reste \emph{convergent} --- l'espérance
conditionnelle est toujours correctement spécifiée, et pondérer ne fait que
repondérer des conditions de moments valides --- mais il n'est plus efficace, et
surtout ses écarts-types conventionnels, calculés comme si les poids étaient
justes, sont \textbf{invalides}. On aurait alors remplacé un problème par un
autre. \textbf{Comparaison.} Les écarts-types robustes ne demandent aucun modèle
de la variance~: ils sont agnostiques, donc quasiment sans risque, au prix d'un
estimateur non efficace. La pondération peut gagner beaucoup en précision, mais
seulement si le modèle de variance est crédible. En pratique, on privilégie les
écarts-types robustes par défaut, et l'on ne pondère que lorsque la forme de
$\Omega$ est connue par construction --- typiquement le cas des données groupées
de l'exercice~\ref{ex:cout-hetero}, où $\sigma_g^2=\epsvar/n_g$ avec $n_g$ observé. Une prudence
supplémentaire consiste à pondérer \emph{et} à calculer des écarts-types
robustes sur le modèle pondéré.

\nouvelexercice\label{ex:mv-recap}

\question Les erreurs étant indépendantes et gaussiennes de variances
$\sigma_t^2$~:
\[
  \ell\left(\beta,\Omega;\mathbf y\right)
  = -\frac{T}{2}\log(2\pi) - \frac{1}{2}\sum_{t=1}^T\log\sigma_t^2
    - \frac{1}{2}\sum_{t=1}^T\frac{\left(y_t-\mathbf x_t'\beta\right)^2}{\sigma_t^2}
\]
soit, sous forme matricielle,
$\ell = -\frac{T}{2}\log(2\pi) - \frac12\log|\Omega| - \frac12\left(\mathbf y-X\beta\right)'\Omega^{-1}\left(\mathbf y-X\beta\right)$.

\question Si $\Omega$ est connue, les deux premiers termes ne dépendent pas de
$\beta$~: maximiser $\ell$ en $\beta$ revient exactement à \emph{minimiser} la
forme quadratique $\left(\mathbf y-X\beta\right)'\Omega^{-1}\left(\mathbf y-X\beta\right)$,
qui est le programme des MCG. D'où~:
\[
  \check{\mathbf b} = \hat{\mathbf b}_{\textsc{mcg}} = \left(X'\Omega^{-1}X\right)^{-1}X'\Omega^{-1}\mathbf y
\]
et ceci \textbf{exactement}, à toute taille d'échantillon --- il ne s'agit pas
d'une équivalence asymptotique. C'est la transposition au cas non sphérique de
l'équivalence MCO/MV du chapitre~I.

\question La dérivée croisée de la log-vraisemblance fait intervenir~:
\[
  \frac{\partial^2\ell}{\partial\beta\,\partial\gamma'}
  = -\sum_t \frac{\left(y_t-\mathbf x_t'\beta\right)}{\sigma_t^2}\,\mathbf x_t z_t'
  = -\sum_t \frac{\varepsilon_t}{\sigma_t^2}\,\mathbf x_t z_t'
\]
dont l'espérance est nulle puisque $\mathbb E[\varepsilon_t]=0$ et que
$\mathbf x_t$, $z_t$ et $\sigma_t^2$ sont déterministes. Le bloc croisé de la
matrice d'information est donc nul~: elle est \textbf{bloc-diagonale}. C'est un
moment d'ordre impair qui s'annule --- le même mécanisme qu'au chapitre~I entre
$\beta$ et $\epsvar$.

\question La bloc-diagonalité signifie que $\beta$ et $\gamma$ sont orthogonaux
au sens de l'information~: l'inverse de la matrice d'information est
bloc-diagonale, de sorte que la borne pour $\beta$ est la même selon que
$\gamma$ est connu ou estimé. \textbf{Estimer $\Omega$ ne coûte donc rien
asymptotiquement}~: l'incertitude sur les poids ne se propage pas, au premier
ordre, à l'estimation des paramètres d'intérêt.

\question Il en résulte que l'estimateur des MCQG (qui utilise $\hat\Omega$) et
l'estimateur du maximum de vraisemblance (qui estime tout conjointement) ont la
\textbf{même distribution asymptotique}~:
\[
  \sqrt T\left(\breve{\mathbf b}-\beta\right)\ \dlimT\
  \normal{0}{\left(\mathrm{plim}\ \frac{X'\Omega^{-1}X}{T}\right)^{-1}}
\]
et tous deux sont asymptotiquement équivalents aux MCG à $\Omega$ connue. Il n'y
a donc pas lieu, asymptotiquement, de préférer l'un à l'autre.

\question
\begin{center}
\begin{tabular}{p{1.6cm}p{2.3cm}p{2.0cm}p{3.2cm}p{3.0cm}}
\hline
 & \textbf{$\Omega$} & \textbf{Normalité} & \textbf{Optimalité} & \textbf{Mise en œuvre}\\
\hline
MCG & connue & non requise & BLUE à distance finie (Gauss-Markov) & directe, mais $\Omega$ l'est rarement\\[1mm]
MCQG & estimée & non requise & asymptotique seulement & procédure en plusieurs étapes\\[1mm]
MV & paramétrée et estimée & requise & efficace asympt. (borne CDFR) & optimisation numérique\\
\hline
\end{tabular}
\end{center}

\question L'estimateur sandwich de l'exercice~\ref{ex:hc} repose sur l'hypothèse que
$\Omega$ est \textbf{diagonale}~: la garniture $\sum_t\hat\varepsilon_t^{\,2}\mathbf x_t\mathbf x_t'$
ne contient que des termes en $t$, aucun terme croisé en $(t,s)$. Il corrige donc
l'inégalité des variances, mais \emph{pas} la présence de covariances hors
diagonale. Or~:
\begin{itemize}
\item En \textbf{série temporelle}, les erreurs sont fréquemment autocorrélées
  (chapitre~II)~: $\Omega$ a des termes hors diagonale, et $X'\Omega X$ comporte
  des sommes doubles que le sandwich ignore. La correction de Newey-West ajoute
  précisément ces termes croisés, pondérés de façon décroissante avec l'écart
  temporel pour garantir une matrice semi-définie positive.
\item En \textbf{données groupées} (élèves d'une même classe, salariés d'une même
  entreprise), les erreurs sont corrélées à l'intérieur des grappes~: $\Omega$ est
  bloc-diagonale et non diagonale. Les écarts-types en grappes agrègent les
  résidus au niveau du groupe plutôt que de l'observation.
\end{itemize}
Dans les deux cas, ignorer la corrélation conduit typiquement à sous-estimer
lourdement les écarts-types.

\nouvelexercice\label{ex:salaire} \textit{Les valeurs ci-dessous ont été obtenues sur
\texttt{data/chapitre-1/wage1.csv} ($T=526$).}

\question L'estimation par les MCO donne ($R^2 = 0{,}40$, $s = 0{,}413$)~:
\[
  \widehat{\log wage} = \underset{(0{,}102)}{0{,}390}
  + \underset{(0{,}0070)}{0{,}0841}\,educ
  + \underset{(0{,}0048)}{0{,}0389}\,exper
  - \underset{(0{,}00011)}{0{,}000686}\,exper^2
  - \underset{(0{,}036)}{0{,}337}\,female
\]
\textbf{Pourquoi le logarithme}~: le salaire est une grandeur positive à
distribution très asymétrique, et la théorie du capital humain (Mincer) prédit
un effet \emph{multiplicatif} de l'éducation. En log, le modèle devient additif
et les coefficients s'interprètent en variation relative~: $\hat\beta_2 = 0{,}084$
signifie qu'une année d'études supplémentaire est associée à un salaire horaire
supérieur d'environ $8{,}4\,\%$. Le terme quadratique en $exper$, de coefficient
négatif, traduit des rendements décroissants de l'expérience (le maximum est
atteint vers $28$ ans). Le coefficient de $female$ indique un salaire inférieur
d'environ $29\,\%$ ($e^{-0{,}337}-1$) pour les femmes, à caractéristiques
observées données.

\question Les deux graphiques montrent un net évasement des résidus~: leur
dispersion croît avec le niveau d'éducation et avec le salaire prédit. Les
diplômés du supérieur ont des salaires beaucoup plus hétérogènes que les
non-diplômés. C'est la signature graphique de l'hétéroscédasticité, et elle
suggère une variance croissante en $educ$ --- ce qui oriente vers le choix de
$z_t$ pour Breusch-Pagan et vers la variable d'ordonnancement de Goldfeld-Quandt.

\question Les trois tests donnent~:
\begin{itemize}
\item \textbf{Breusch-Pagan} (version de Koenker, $z_t$ = les quatre
  régresseurs)~: $LM = 9{,}98$, $\chi^2(4)$, $p = 0{,}041$. On rejette
  l'homoscédasticité au seuil de $5\,\%$.
\item \textbf{White} (régression auxiliaire complète)~: $W = 19{,}93$. Le
  décompte naïf donnerait $q = 4+\frac{4\times5}{2} = 14$, mais deux colonnes
  sont colinéaires ($exper\times exper = exper^2$, et $female\times female = female$
  puisque $female$ est une indicatrice), d'où $q = 12$ et
  $p = 0{,}068$. On ne rejette \emph{pas} tout à fait au seuil de $5\,\%$.
\item \textbf{Goldfeld-Quandt} (tri selon $educ$, un quart des observations
  centrales écartées)~: $F = 1{,}33$, $p = 0{,}026$. On rejette à $5\,\%$.
\end{itemize}
\textbf{Les conclusions ne concordent pas tout à fait}, et c'est instructif.
Breusch-Pagan et Goldfeld-Quandt rejettent, White non. L'explication tient à la
puissance~: White disperse sa puissance sur $12$ degrés de liberté pour détecter
n'importe quelle forme d'hétéroscédasticité, alors que les deux autres
concentrent la leur sur une alternative monotone en $educ$ --- qui se trouve être
la bonne. C'est exactement l'arbitrage généralité/puissance de l'exercice~\ref{ex:white}.

\question Écarts-types comparés~:
\begin{center}
\begin{tabular}{lccccc}
\hline
 & usuel & HC0 & HC1 & HC2 & HC3\\
\hline
$educ$   & $0{,}00696$ & $0{,}00767$ & $0{,}00770$ & $0{,}00773$ & $0{,}00781$\\
$exper$  & $0{,}00482$ & $0{,}00470$ & $0{,}00472$ & $0{,}00472$ & $0{,}00475$\\
$female$ & $0{,}03634$ & $0{,}03604$ & $0{,}03621$ & $0{,}03623$ & $0{,}03643$\\
\hline
\end{tabular}
\end{center}
La correction est modeste et \emph{ne va pas dans le même sens pour toutes les
variables}~: l'écart-type de $educ$ augmente d'environ $12\,\%$ (sa statistique
de Student passe de $12{,}1$ à $10{,}8$), tandis que ceux de $exper$ et de
$female$ diminuent légèrement. \textbf{Aucun coefficient ne change de statut}~:
tous restent très significatifs. C'est un rappel utile --- l'hétéroscédasticité
ne gonfle pas mécaniquement tous les écarts-types.

\question Ici $T = 526$, très au-delà du seuil de $250$, et les leviers restent
faibles ($\max_t h_{tt} = 0{,}054$ pour un levier moyen $K/T = 0{,}0095$). Les
quatre variantes sont d'ailleurs presque confondues, ce qui confirme le
diagnostic. \textbf{N'importe laquelle convient}, et l'on retiendra HC3 par
défaut puisqu'elle ne coûte rien sur un échantillon de cette taille.

\question Avec $\sigma_i^2 = \sigma^2\exp(\gamma_1 + \gamma_2\,educ_i + \gamma_3\,exper_i)$,
la régression auxiliaire de $\log\hat\varepsilon_i^{\,2}$ donne
$\hat\gamma_2 = 0{,}097$ ($t = 2{,}5$) et $\hat\gamma_3 = 0{,}021$ ($t = 2{,}7$)~:
la variance croît bien avec l'éducation et l'expérience, ce qui confirme les
tests. Les MCQG donnent alors~:
\begin{center}
\begin{tabular}{lcccc}
\hline
 & \multicolumn{2}{c}{MCO} & \multicolumn{2}{c}{MCQG}\\
 & coef. & éc.-type & coef. & éc.-type\\
\hline
$educ$   & $0{,}0841$ & $0{,}0070$ & $0{,}0743$ & $0{,}0065$\\
$exper$  & $0{,}0389$ & $0{,}0048$ & $0{,}0415$ & $0{,}0047$\\
$female$ & $-0{,}3372$ & $0{,}0363$ & $-0{,}3005$ & $0{,}0355$\\
\hline
\end{tabular}
\end{center}
Le gain de précision est réel mais modeste (environ $6\,\%$ sur $educ$).
\textbf{Un point mérite attention}~: le coefficient de $educ$ passe de $0{,}084$
à $0{,}074$, soit un écart de près d'un écart-type et demi. Sous spécification
correcte, MCO et MCQG sont tous deux convergents et ne devraient différer que
par l'aléa d'échantillonnage. Un écart systématique de cette ampleur est un
signal de type Hausman~: il invite à soupçonner une mauvaise spécification de
l'espérance conditionnelle --- forme fonctionnelle, ou variable omise
(le talent, cf. exercice~\ref{ex:variable-omise}) --- plutôt qu'à se réjouir du gain d'efficacité.

\question \textbf{Conclusion.} L'hétéroscédasticité est ici statistiquement
détectable mais \emph{économiquement inoffensive}~: les corrections déplacent
les écarts-types de quelques pourcents, aucun coefficient ne change de signe ni
de statut de significativité, et les conclusions de l'étude --- rendement de
l'éducation d'environ $8\,\%$ par année, rendements décroissants de
l'expérience, écart salarial femmes-hommes d'environ $30\,\%$ --- sont
inchangées. \textbf{Que privilégier}~: les écarts-types robustes. Ils ne coûtent
rien, ne demandent aucune hypothèse sur la forme de $\Omega$, et suffisent
lorsque l'enjeu est l'inférence. La pondération n'apporte ici qu'un gain
d'efficacité marginal, au prix d'un modèle de variance arbitraire --- et elle
soulève, comme on vient de le voir, une question de spécification plus sérieuse
que celle qu'elle prétendait régler. C'est bien le message du chapitre~: le
danger n'était pas dans l'estimateur, mais dans l'écart-type imprimé à côté.

\nouvelexercice\label{ex:mpl}

\question Puisque $y_i$ vaut $0$ ou $1$, l'erreur ne peut prendre que deux
valeurs~:
\[
  \varepsilon_i =
  \begin{cases}
    1-\mathbf x_i'\beta & \text{avec probabilité } \mathbf x_i'\beta\\
    -\mathbf x_i'\beta  & \text{avec probabilité } 1-\mathbf x_i'\beta
  \end{cases}
\]
L'erreur est donc, elle aussi, une variable de Bernoulli translatée --- très loin
d'une gaussienne.

\question L'espérance conditionnelle vaut~:
\[
  \mathbb E\left[\varepsilon_i|\mathbf x_i\right]
  = \left(1-\mathbf x_i'\beta\right)\mathbf x_i'\beta + \left(-\mathbf x_i'\beta\right)\left(1-\mathbf x_i'\beta\right) = 0
\]
Pour la variance, en posant $\pi_i = \mathbf x_i'\beta$~:
\[
  \mathbb V\left[\varepsilon_i|\mathbf x_i\right]
  = \mathbb E\left[\varepsilon_i^2|\mathbf x_i\right]
  = \left(1-\pi_i\right)^2\pi_i + \pi_i^2\left(1-\pi_i\right)
  = \pi_i\left(1-\pi_i\right)\left[\left(1-\pi_i\right)+\pi_i\right]
  = \pi_i\left(1-\pi_i\right)
\]
(c'est simplement la variance d'une Bernoulli de paramètre $\pi_i$).
\textbf{Commentaire.} La variance dépend de $\mathbf x_i$ à travers $\pi_i$~:
elle est maximale ($1/4$) pour $\pi_i=1/2$ et tend vers zéro lorsque $\pi_i$
s'approche de $0$ ou de $1$. L'hétéroscédasticité n'est donc pas ici une
imperfection des données que l'on pourrait espérer absente~: elle est une
\emph{conséquence logique} du caractère binaire de $y$. Aucun test n'est
nécessaire pour la diagnostiquer --- sauf à vouloir vérifier la spécification.
C'est le seul cas du cours où la forme de $\Omega$ est connue \emph{comme
fonction des paramètres à estimer}.

\question L'estimateur des MCO est \textbf{sans biais et convergent}~: ces deux
propriétés ne reposent que sur $\mathbb E[\varepsilon|\mathbf x]=0$, qui est ici
vraie par construction dès lors que l'espérance conditionnelle est effectivement
linéaire (exercice~\ref{ex:cout-hetero}, question 1). En revanche il n'est
\textbf{pas efficace}~: le théorème de Gauss-Markov ne s'applique pas, puisque
les erreurs sont hétéroscédastiques. Et surtout, les écarts-types usuels
$s^2(X'X)^{-1}$ sont \textbf{faux}.

\question Les poids optimaux sont l'inverse des variances~:
\[
  \omega_i = \frac{1}{\pi_i\left(1-\pi_i\right)} = \frac{1}{\mathbf x_i'\beta\left(1-\mathbf x_i'\beta\right)}
\]
\textbf{On ne peut pas les utiliser directement} parce qu'ils dépendent de
$\beta$, le paramètre que l'on cherche précisément à estimer~: c'est un
raisonnement circulaire. D'où la procédure en deux étapes (MCQG)~: \emph{(i)}
estimer $\beta$ par les MCO ordinaires, ce qui est légitime puisque l'estimateur
est convergent, et former les probabilités ajustées
$\hat\pi_i = \mathbf x_i'\hat{\mathbf b}$~; \emph{(ii)} ré-estimer par moindres
carrés pondérés avec les poids $\hat\omega_i = 1/[\hat\pi_i(1-\hat\pi_i)]$.

\question Rien dans les MCO ne contraint $\mathbf x_i'\hat{\mathbf b}$ à rester
dans $[0,1]$~: c'est une combinaison linéaire non bornée des explicatives, et
pour des valeurs extrêmes de $\mathbf x_i$ la prédiction sort de l'intervalle.
On obtient alors des «~probabilités~» négatives ou supérieures à $1$, ce qui est
déjà gênant en soi, mais devient \emph{fatal} pour la pondération~: si
$\hat\pi_i\notin[0,1]$, alors $\hat\pi_i(1-\hat\pi_i)<0$ et le poids
$\hat\omega_i$ est \textbf{négatif}. L'estimateur pondéré n'a plus aucun sens --- 
une observation contribuerait négativement à la somme des carrés. \textbf{Deux
remèdes pratiques}~: écarter les observations concernées (au prix d'une
sélection arbitraire de l'échantillon, ce qui peut biaiser), ou tronquer les
probabilités ajustées dans un intervalle du type $[0{,}01\ ;\ 0{,}99]$ (au prix
d'un choix tout aussi arbitraire). Aucune de ces solutions n'est satisfaisante.

\question Précisément à cause de la question précédente. La pondération promet un
gain d'efficacité, mais elle se heurte à un obstacle structurel --- des poids
possiblement négatifs --- et ne peut être sauvée que par des bricolages
arbitraires. Les écarts-types robustes, eux, ne demandent rien~: ils
n'exigent ni de connaître $\Omega$, ni de la modéliser, ni que les probabilités
ajustées soient bien rangées. On conserve donc les MCO, dont on sait qu'ils sont
convergents, et l'on se contente de corriger leur variance par l'estimateur
sandwich. C'est exactement la stratégie recommandée à
l'exercice~\ref{ex:mcp}~: pondérer seulement lorsque le modèle de variance est
crédible, et se rabattre sinon sur les écarts-types robustes. Le modèle de
probabilité linéaire s'estime ainsi en pratique~: MCO plus écarts-types robustes,
systématiquement.

\question \textbf{Comparaison.}
\begin{itemize}
\item \textbf{Avantages du modèle de probabilité linéaire}~: les coefficients
  s'interprètent directement comme des effets marginaux sur la probabilité
  (une unité de $x_k$ en plus fait varier $\mathbb P(y=1)$ de $\beta_k$), sans
  calcul supplémentaire~; l'estimation est linéaire, donc immédiate, sans
  optimisation numérique ni problème de convergence~; il se combine sans peine
  aux variables instrumentales, aux effets fixes et aux données de panel~; et
  il fournit une bonne approximation de l'effet marginal moyen lorsque les
  probabilités estimées restent au voisinage de $1/2$.
\item \textbf{Inconvénients}~: les probabilités prédites peuvent sortir de
  $[0,1]$~; l'effet marginal est supposé constant, alors qu'il devrait
  logiquement s'atténuer près de $0$ et de $1$~; et l'hypothèse de linéarité de
  l'espérance conditionnelle est forte.
\item \textbf{Probit et logit} corrigent ces défauts en spécifiant
  $\mathbb P(y_i=1|\mathbf x_i) = F(\mathbf x_i'\beta)$ avec $F$ une fonction de
  répartition, ce qui contraint la prédiction à $]0,1[$ et rend l'effet marginal
  variable. Le prix à payer est une estimation par maximum de vraisemblance
  (donc non linéaire), des coefficients qui ne s'interprètent plus directement,
  et une plus grande sensibilité à la spécification retenue pour $F$.
\end{itemize}
En pratique, les deux approches donnent des effets marginaux moyens très proches
--- ce que confirment les simulations de \texttt{exercices/E03.py} --- et le choix
relève souvent de la commodité.

\end{document}

%%% Local Variables:
%%% mode: latex
%%% TeX-master: t
%%% End:
